\chapter{群论图片转写待审核稿}

\begin{dxtips}

本文件为图片与 Markdown 笔记整理出的待审核转写稿，已尽量按 elegantbook 环境归类；正式并入正文前仍建议逐条复核。

\end{dxtips}

\section*{1.10阿贝尔群}

\begin{note}[来源上级主题]

1.2群 (1\%202\%E7\%BE\%A4\%201ecd2efa3f4280cf8df9c881d186f2a4.md)

\end{note}

\begin{note}[1.10阿贝尔群]

子级 项目: 引理1.1幺半群的逆元构成群 (\%E5\%BC\%95\%E7\%90\%861\%201\%E5\%B9\%BA\%E5\%8D\%8A\%E7\%BE\%A4\%E7\%9A\%84\%E9\%80\%86\%E5\%85\%83\%E6\%9E\%84\%E6\%88\%90\%E7\%BE\%A4\%201edd2efa3f4280479e2be2b835b83ca6.md)

常见的群，加群最为普遍

\end{note}

\begin{note}[关联图片]
以下图片已纳入本条整理，当前版本优先依据 Markdown 文字整理，图片细节可在审核时对照：

1. dd068fda-ce9f-4c25-982c-dccc2d5878e8.png\\

\end{note}

\begin{note}[来源路径]

近世代数/近世代数/无标题/1 10阿贝尔群 1cbd2efa3f4280bfa73dc117dd0d632c.md

\end{note}

\section*{1.1幺半群}

\begin{note}[1.1幺半群]

子级 项目: 幺半群的定义 定义1.2 定义1.3 (\%E5\%B9\%BA\%E5\%8D\%8A\%E7\%BE\%A4\%E7\%9A\%84\%E5\%AE\%9A\%E4\%B9\%89\%20\%E5\%AE\%9A\%E4\%B9\%891\%202\%20\%E5\%AE\%9A\%E4\%B9\%891\%203\%201cad2efa3f42800ea08fdf383be45a9a.md), 定义1.3群中的幂 (\%E5\%AE\%9A\%E4\%B9\%891\%203\%E7\%BE\%A4\%E4\%B8\%AD\%E7\%9A\%84\%E5\%B9\%82\%201cad2efa3f4280fa8cbac5abca50f603.md), 广义结合律与数学归纳法 (\%E5\%B9\%BF\%E4\%B9\%89\%E7\%BB\%93\%E5\%90\%88\%E5\%BE\%8B\%E4\%B8\%8E\%E6\%95\%B0\%E5\%AD\%A6\%E5\%BD\%92\%E7\%BA\%B3\%E6\%B3\%95\%201cad2efa3f42803e8230cae3ea9f46af.md), 定义1.4幺半群的子群还是幺半群 (\%E5\%AE\%9A\%E4\%B9\%891\%204\%E5\%B9\%BA\%E5\%8D\%8A\%E7\%BE\%A4\%E7\%9A\%84\%E5\%AD\%90\%E7\%BE\%A4\%E8\%BF\%98\%E6\%98\%AF\%E5\%B9\%BA\%E5\%8D\%8A\%E7\%BE\%A4\%201cad2efa3f4280a186e5ff06d952436b.md), 定义1.5幺半群的同态 (\%E5\%AE\%9A\%E4\%B9\%891\%205\%E5\%B9\%BA\%E5\%8D\%8A\%E7\%BE\%A4\%E7\%9A\%84\%E5\%90\%8C\%E6\%80\%81\%201cad2efa3f4280878f7ace8395a0a565.md), 定义1.6由子集生成的自幺半群 (\%E5\%AE\%9A\%E4\%B9\%891\%206\%E7\%94\%B1\%E5\%AD\%90\%E9\%9B\%86\%E7\%94\%9F\%E6\%88\%90\%E7\%9A\%84\%E8\%87\%AA\%E5\%B9\%BA\%E5\%8D\%8A\%E7\%BE\%A4\%201cad2efa3f4280c8b8e4ca2c6bd26b3d.md), 命题1.5证明A生成的子幺半群是包含A的最小子幺半群 (\%E5\%91\%BD\%E9\%A2\%981\%205\%E8\%AF\%81\%E6\%98\%8EA\%E7\%94\%9F\%E6\%88\%90\%E7\%9A\%84\%E5\%AD\%90\%E5\%B9\%BA\%E5\%8D\%8A\%E7\%BE\%A4\%E6\%98\%AF\%E5\%8C\%85\%E5\%90\%ABA\%E7\%9A\%84\%E6\%9C\%80\%E5\%B0\%8F\%E5\%AD\%90\%E5\%B9\%BA\%E5\%8D\%8A\%E7\%BE\%A4\%201cad2efa3f4280d7aeccdd9429149cb4.md), 定义1.7幺半群的同构 (\%E5\%AE\%9A\%E4\%B9\%891\%207\%E5\%B9\%BA\%E5\%8D\%8A\%E7\%BE\%A4\%E7\%9A\%84\%E5\%90\%8C\%E6\%9E\%84\%201cbd2efa3f4280e49eafe83547a86380.md), 命题1.7证明方法把单位元拆解然后用结合律引出新的东西唯一性证明 (\%E5\%91\%BD\%E9\%A2\%981\%207\%E8\%AF\%81\%E6\%98\%8E\%E6\%96\%B9\%E6\%B3\%95\%E6\%8A\%8A\%E5\%8D\%95\%E4\%BD\%8D\%E5\%85\%83\%E6\%8B\%86\%E8\%A7\%A3\%E7\%84\%B6\%E5\%90\%8E\%E7\%94\%A8\%E7\%BB\%93\%E5\%90\%88\%E5\%BE\%8B\%E5\%BC\%95\%E5\%87\%BA\%E6\%96\%B0\%E7\%9A\%84\%E4\%B8\%9C\%E8\%A5\%BF\%E5\%94\%AF\%E4\%B8\%80\%E6\%80\%A7\%E8\%AF\%81\%E6\%98\%8E\%201cbd2efa3f428061b991eda66a78e31d.md)

\end{note}

\begin{note}[来源路径]

近世代数/近世代数/无标题/1 1幺半群 1cad2efa3f4280628789c382df95fdc5.md

\end{note}

\section*{1.2群}

\begin{note}[1.2群]

子级 项目: 定义1.9幺半群加逆元就是群 (\%E5\%AE\%9A\%E4\%B9\%891\%209\%E5\%B9\%BA\%E5\%8D\%8A\%E7\%BE\%A4\%E5\%8A\%A0\%E9\%80\%86\%E5\%85\%83\%E5\%B0\%B1\%E6\%98\%AF\%E7\%BE\%A4\%201cbd2efa3f4280789737e0626b159cd6.md), 定义1.8幺半群的逆元 (\%E5\%AE\%9A\%E4\%B9\%891\%208\%E5\%B9\%BA\%E5\%8D\%8A\%E7\%BE\%A4\%E7\%9A\%84\%E9\%80\%86\%E5\%85\%83\%201cbd2efa3f4280ab8ab0c76329b68410.md), 1.10阿贝尔群 (1\%2010\%E9\%98\%BF\%E8\%B4\%9D\%E5\%B0\%94\%E7\%BE\%A4\%201cbd2efa3f4280bfa73dc117dd0d632c.md), 定义1.11n阶一般线性群 (\%E5\%AE\%9A\%E4\%B9\%891\%2011n\%E9\%98\%B6\%E4\%B8\%80\%E8\%88\%AC\%E7\%BA\%BF\%E6\%80\%A7\%E7\%BE\%A4\%201cbd2efa3f4280dda90ce3802a0d1200.md), 定义1.12n阶特殊线性群行列式是1 (\%E5\%AE\%9A\%E4\%B9\%891\%2012n\%E9\%98\%B6\%E7\%89\%B9\%E6\%AE\%8A\%E7\%BA\%BF\%E6\%80\%A7\%E7\%BE\%A4\%E8\%A1\%8C\%E5\%88\%97\%E5\%BC\%8F\%E6\%98\%AF1\%201cbd2efa3f42805e8f6be357d2dce70f.md), 定义1.13定义子群根本目的是证明它是一个群 (\%E5\%AE\%9A\%E4\%B9\%891\%2013\%E5\%AE\%9A\%E4\%B9\%89\%E5\%AD\%90\%E7\%BE\%A4\%E6\%A0\%B9\%E6\%9C\%AC\%E7\%9B\%AE\%E7\%9A\%84\%E6\%98\%AF\%E8\%AF\%81\%E6\%98\%8E\%E5\%AE\%83\%E6\%98\%AF\%E4\%B8\%80\%E4\%B8\%AA\%E7\%BE\%A4\%201cbd2efa3f4280ee903ac1e652f44dec.md), 定义1.15群同态下定义 核是映射到e‘的原像集 imf是像集 (\%E5\%AE\%9A\%E4\%B9\%891\%2015\%E7\%BE\%A4\%E5\%90\%8C\%E6\%80\%81\%E4\%B8\%8B\%E5\%AE\%9A\%E4\%B9\%89\%20\%E6\%A0\%B8\%E6\%98\%AF\%E6\%98\%A0\%E5\%B0\%84\%E5\%88\%B0e\%E2\%80\%98\%E7\%9A\%84\%E5\%8E\%9F\%E5\%83\%8F\%E9\%9B\%86\%20imf\%E6\%98\%AF\%E5\%83\%8F\%E9\%9B\%86\%201d1d2efa3f42808daa56cdc3a30925f6.md), 定义1.16群同态 (\%E5\%AE\%9A\%E4\%B9\%891\%2016\%E7\%BE\%A4\%E5\%90\%8C\%E6\%80\%81\%201d1d2efa3f4280dfbb71fb6ae82fca39.md), 命题1.17同态就是拆括号的游戏 (\%E5\%91\%BD\%E9\%A2\%981\%2017\%E5\%90\%8C\%E6\%80\%81\%E5\%B0\%B1\%E6\%98\%AF\%E6\%8B\%86\%E6\%8B\%AC\%E5\%8F\%B7\%E7\%9A\%84\%E6\%B8\%B8\%E6\%88\%8F\%201d1d2efa3f42800d9bf2efcc79be14a6.md), 定义1.18两个群直积还是群 (\%E5\%AE\%9A\%E4\%B9\%891\%2018\%E4\%B8\%A4\%E4\%B8\%AA\%E7\%BE\%A4\%E7\%9B\%B4\%E7\%A7\%AF\%E8\%BF\%98\%E6\%98\%AF\%E7\%BE\%A4\%201d1d2efa3f4280189fb9c52e96214d0c.md)

\end{note}

\begin{note}[来源路径]

近世代数/近世代数/无标题/1 2群 1ecd2efa3f4280cf8df9c881d186f2a4.md

\end{note}

\section*{1.3有限群}

\begin{note}[1.3有限群]

子级 项目: 定义1.19一族群的值积 定义1.20还需再看 (\%E5\%AE\%9A\%E4\%B9\%891\%2019\%E4\%B8\%80\%E6\%97\%8F\%E7\%BE\%A4\%E7\%9A\%84\%E5\%80\%BC\%E7\%A7\%AF\%20\%E5\%AE\%9A\%E4\%B9\%891\%2020\%E8\%BF\%98\%E9\%9C\%80\%E5\%86\%8D\%E7\%9C\%8B\%201d1d2efa3f4280bbb5e6c6260092ce70.md), 定义1.22 x的幂 (\%E5\%AE\%9A\%E4\%B9\%891\%2022\%20x\%E7\%9A\%84\%E5\%B9\%82\%201d1d2efa3f42801fbf26e402d99d9fbb.md), 定义1.23由x生成子群 (\%E5\%AE\%9A\%E4\%B9\%891\%2023\%E7\%94\%B1x\%E7\%94\%9F\%E6\%88\%90\%E5\%AD\%90\%E7\%BE\%A4\%201d1d2efa3f428063a543c4b613e77eac.md), 定义1.2.4子集生成群证明 (\%E5\%AE\%9A\%E4\%B9\%891\%202\%204\%E5\%AD\%90\%E9\%9B\%86\%E7\%94\%9F\%E6\%88\%90\%E7\%BE\%A4\%E8\%AF\%81\%E6\%98\%8E\%201d1d2efa3f42805fa62fd5a0548e8041.md), 定义1.25循环群 (\%E5\%AE\%9A\%E4\%B9\%891\%2025\%E5\%BE\%AA\%E7\%8E\%AF\%E7\%BE\%A4\%201d1d2efa3f42803f81baf1a7f0cbd190.md), 命题1.2.6 带余除法和阶n的结合 (\%E5\%91\%BD\%E9\%A2\%981\%202\%206\%20\%E5\%B8\%A6\%E4\%BD\%99\%E9\%99\%A4\%E6\%B3\%95\%E5\%92\%8C\%E9\%98\%B6n\%E7\%9A\%84\%E7\%BB\%93\%E5\%90\%88\%201d2d2efa3f4280f4aba7f41bf31cebda.md), n阶循环群互相同构 无限循环群彼此同构 命题1.2.9 (n\%E9\%98\%B6\%E5\%BE\%AA\%E7\%8E\%AF\%E7\%BE\%A4\%E4\%BA\%92\%E7\%9B\%B8\%E5\%90\%8C\%E6\%9E\%84\%20\%E6\%97\%A0\%E9\%99\%90\%E5\%BE\%AA\%E7\%8E\%AF\%E7\%BE\%A4\%E5\%BD\%BC\%E6\%AD\%A4\%E5\%90\%8C\%E6\%9E\%84\%20\%E5\%91\%BD\%E9\%A2\%981\%202\%209\%201d2d2efa3f4280ebb6eaf07515b209e7.md), 命题1.30循环群的生成元的幂指数和循环群阶数互素 (\%E5\%91\%BD\%E9\%A2\%981\%2030\%E5\%BE\%AA\%E7\%8E\%AF\%E7\%BE\%A4\%E7\%9A\%84\%E7\%94\%9F\%E6\%88\%90\%E5\%85\%83\%E7\%9A\%84\%E5\%B9\%82\%E6\%8C\%87\%E6\%95\%B0\%E5\%92\%8C\%E5\%BE\%AA\%E7\%8E\%AF\%E7\%BE\%A4\%E9\%98\%B6\%E6\%95\%B0\%E4\%BA\%92\%E7\%B4\%A0\%201d2d2efa3f4280cfaefbd295458b0a99.md), 子群的阶定义拉格朗日定理 (\%E5\%AD\%90\%E7\%BE\%A4\%E7\%9A\%84\%E9\%98\%B6\%E5\%AE\%9A\%E4\%B9\%89\%E6\%8B\%89\%E6\%A0\%BC\%E6\%9C\%97\%E6\%97\%A5\%E5\%AE\%9A\%E7\%90\%86\%201d2d2efa3f4280abb951d3f9b92ef54b.md), 左陪集定义 (\%E5\%B7\%A6\%E9\%99\%AA\%E9\%9B\%86\%E5\%AE\%9A\%E4\%B9\%89\%201d2d2efa3f4280fdbd2fd7e28903354a.md), 命题1.3.3子群陪集要不然相等要不然无交 (\%E5\%91\%BD\%E9\%A2\%981\%203\%203\%E5\%AD\%90\%E7\%BE\%A4\%E9\%99\%AA\%E9\%9B\%86\%E8\%A6\%81\%E4\%B8\%8D\%E7\%84\%B6\%E7\%9B\%B8\%E7\%AD\%89\%E8\%A6\%81\%E4\%B8\%8D\%E7\%84\%B6\%E6\%97\%A0\%E4\%BA\%A4\%201d2d2efa3f42806da9c1eaa41dc2aa66.md), 定义1.28商集与指数 (\%E5\%AE\%9A\%E4\%B9\%891\%2028\%E5\%95\%86\%E9\%9B\%86\%E4\%B8\%8E\%E6\%8C\%87\%E6\%95\%B0\%201d2d2efa3f4280779550ef45dc732f24.md), 拉格朗日定理 (\%E6\%8B\%89\%E6\%A0\%BC\%E6\%9C\%97\%E6\%97\%A5\%E5\%AE\%9A\%E7\%90\%86\%201d2d2efa3f42801fa5fdf15cafbdcb4e.md), 子群里面的元素当陪集的x，则这个配集还是原来子群xH=H (\%E5\%AD\%90\%E7\%BE\%A4\%E9\%87\%8C\%E9\%9D\%A2\%E7\%9A\%84\%E5\%85\%83\%E7\%B4\%A0\%E5\%BD\%93\%E9\%99\%AA\%E9\%9B\%86\%E7\%9A\%84x\%EF\%BC\%8C\%E5\%88\%99\%E8\%BF\%99\%E4\%B8\%AA\%E9\%85\%8D\%E9\%9B\%86\%E8\%BF\%98\%E6\%98\%AF\%E5\%8E\%9F\%E6\%9D\%A5\%E5\%AD\%90\%E7\%BE\%A4xH=H\%201d2d2efa3f4280d29dd3d8346705ba89.md), 命题1.34 1.35嵌套子群的配集 (\%E5\%91\%BD\%E9\%A2\%981\%2034\%201\%2035\%E5\%B5\%8C\%E5\%A5\%97\%E5\%AD\%90\%E7\%BE\%A4\%E7\%9A\%84\%E9\%85\%8D\%E9\%9B\%86\%201d2d2efa3f428089be3ecd9b4c4dbf6c.md)

\end{note}

\begin{note}[来源路径]

近世代数/近世代数/无标题/1 3有限群 1ead2efa3f4280c39a49d8b0af9aec03.md

\end{note}

\section*{1.4正规子群}

\begin{note}[1.4正规子群]

子级 项目: 定义1.29正规子群 (\%E5\%AE\%9A\%E4\%B9\%891\%2029\%E6\%AD\%A3\%E8\%A7\%84\%E5\%AD\%90\%E7\%BE\%A4\%201d2d2efa3f428009a478d7345a44f7a0.md), 命题1.36良定义 (\%E5\%91\%BD\%E9\%A2\%981\%2036\%E8\%89\%AF\%E5\%AE\%9A\%E4\%B9\%89\%201d8d2efa3f4280f894b7ef35415e3855.md), 命题1.37 商群 (\%E5\%91\%BD\%E9\%A2\%981\%2037\%20\%E5\%95\%86\%E7\%BE\%A4\%201d8d2efa3f4280a2ab8bc71b5decdd72.md), 引理1.4正规子群 (\%E5\%BC\%95\%E7\%90\%861\%204\%E6\%AD\%A3\%E8\%A7\%84\%E5\%AD\%90\%E7\%BE\%A4\%201d8d2efa3f4280599a3fd672105fee0c.md), 命题1.38正规子群的交集还是正规子群 (\%E5\%91\%BD\%E9\%A2\%981\%2038\%E6\%AD\%A3\%E8\%A7\%84\%E5\%AD\%90\%E7\%BE\%A4\%E7\%9A\%84\%E4\%BA\%A4\%E9\%9B\%86\%E8\%BF\%98\%E6\%98\%AF\%E6\%AD\%A3\%E8\%A7\%84\%E5\%AD\%90\%E7\%BE\%A4\%201d8d2efa3f428001ba66de6c2e2fde67.md), 命题1.39 1.40e和自己都是自己的正规子群交换群的子群都是正规子群 (\%E5\%91\%BD\%E9\%A2\%981\%2039\%201\%2040e\%E5\%92\%8C\%E8\%87\%AA\%E5\%B7\%B1\%E9\%83\%BD\%E6\%98\%AF\%E8\%87\%AA\%E5\%B7\%B1\%E7\%9A\%84\%E6\%AD\%A3\%E8\%A7\%84\%E5\%AD\%90\%E7\%BE\%A4\%E4\%BA\%A4\%E6\%8D\%A2\%E7\%BE\%A4\%E7\%9A\%84\%E5\%AD\%90\%E7\%BE\%A4\%E9\%83\%BD\%E6\%98\%AF\%E6\%AD\%A3\%E8\%A7\%84\%E5\%AD\%90\%E7\%BE\%A4\%201d8d2efa3f4280beb4f3c1126c211d2b.md)

\end{note}

\begin{note}[来源路径]

近世代数/近世代数/无标题/1 4正规子群 1ead2efa3f428095a412e65ae0862a5d.md

\end{note}

\section*{4元群}

\begin{note}[来源上级主题]

3.4 (3\%204\%20209d2efa3f4280eca043fa61064594c9.md)

\end{note}

\begin{note}[4元群]

特征是加法所有非零元素的阶   阶数

特征的定义  对加法所有非零元素共同的阶

拉格朗日定理

子群的阶定义拉格朗日定理

有限域一定有阶 阶只能是4或2 只有2是素数只能是2

域一定有特征 特征只能2

\end{note}

\begin{note}[关联图片]
以下图片已纳入本条整理，当前版本优先依据 Markdown 文字整理，图片细节可在审核时对照：

1. 11021749127514\_.pic.jpg\\
2. image.png\\
3. image 1.png\\

\end{note}

\begin{note}[来源路径]

近世代数/近世代数/无标题/4元群 209d2efa3f42801c9f8ef20ec54114ef.md

\end{note}

\section*{a和a的逆阶数一样}

\begin{note}[来源上级主题]

定义1.20元素的阶 (\%E5\%AE\%9A\%E4\%B9\%891\%2020\%E5\%85\%83\%E7\%B4\%A0\%E7\%9A\%84\%E9\%98\%B6\%201afd2efa3f42804abe31e7c2129a941d.md)

\end{note}

\begin{note}[a和a的逆阶数一样]

设a的阶数是m a逆的m次方也是e m大于等于a逆的阶数

同理设a逆的阶数是n a的n次方也是e n大于等于a的阶数

m大于等于n n大于等于m

\end{note}

\begin{note}[关联图片]
以下图片已纳入本条整理，当前版本优先依据 Markdown 文字整理，图片细节可在审核时对照：

1. image.png\\

\end{note}

\begin{note}[来源路径]

近世代数/近世代数/无标题/a和a的逆阶数一样 1b1d2efa3f4280c9848ef6cf1a5a43d6.md

\end{note}

\section*{n次对称群S\_n的阶是n！}

\begin{note}[来源上级主题]

置换群 (\%E7\%BD\%AE\%E6\%8D\%A2\%E7\%BE\%A4\%201bad2efa3f4280228c17c29195a2c525.md)

\end{note}

\begin{note}[n次对称群S\_n的阶是n！]

EH 1 nwo S, Hen.

RERNRZB-BRA-TERWES. ROA SHRAMM, KN

ASS — At. RINB—-T+ BM

T: A; —>aQ,, 5  1=1, 2, 5 n

RP BRT RAM ERASE ATU C1, Rid, (2, hed, or, Cn, kh, KN

WEAR. RNR BRYN B—THEMAEY EXT BRS MK

1  2  eee  n

(,, kp om .):

EXMRATER, BHA n SRE KPERRATARA, WARWE

BY « Fe 11th By A

2 1 3  n

(ok bm kL)

KRM. ATRNAMRS HIRE 1, 2, +, nXPKF, HK RAEA

Bey AA. FVAMTAS—A Gi.

\end{note}

\begin{note}[图片 OCR 摘录，待人工复核]
以下内容来自源图片识别，便于审核时对照：

\textbf{图片 1}（image.png）
\begin{Verbatim}[fontsize=\small]
EH 1 nwo S, Hen.
RERNRZB-BRA-TERWES. ROA SHRAMM, KN
ASS — At. RINB—-T+ BM
T: A; —>aQ,, 5  1=1, 2, 5 n
RP BRT RAM ERASE ATU C1, Rid, (2, hed, or, Cn, kh, KN
WEAR. RNR BRYN B—THEMAEY EXT BRS MK
1  2  eee  n
(,, kp om .):
EXMRATER, BHA n SRE KPERRATARA, WARWE
BY « Fe 11th By A
2 1 3  n
(ok bm kL)
KRM. ATRNAMRS HIRE 1, 2, +, nXPKF, HK RAEA
Bey AA. FVAMTAS—A Gi.
\end{Verbatim}

\end{note}

\begin{note}[来源路径]

近世代数/近世代数/无标题/n次对称群S\_n的阶是n！ 1bad2efa3f4280749db7c5ede391c5b3.md

\end{note}

\section*{不变子群的中心}

\begin{note}[不变子群的中心]

Sl2 NAIA SHCG HAA PHM n,

na=an, AE a eG HBP

MBA NE GCH-THEFE. AA N2 ec, ALN RESH. KM

NA=AN,y NA=AN>N N.a=N, any =aN, Ny

na=an>n ‘a=n ann! =n"'nan'=an!

mEN, mE N>nn EN; n©N>n!€N

Hill, 8, \#1, NE—-TTHR. AGHS tio TUNA

B70 n Hk, UREA Na=aN, El NBREFR,

i “PY AB EE FRE AY file G AY AR.

\end{note}

\begin{note}[图片 OCR 摘录，待人工复核]
以下内容来自源图片识别，便于审核时对照：

\textbf{图片 1}（image.png）
\begin{Verbatim}[fontsize=\small]
Sl2 NAIA SHCG HAA PHM n,
na=an, AE a eG HBP

MBA NE GCH-THEFE. AA N2 ec, ALN RESH. KM
NA=AN,y NA=AN>N N.a=N, any =aN, Ny
na=an>n ‘a=n ann! =n"'nan'=an!
mEN, mE N>nn EN; n©N>n!€N

Hill, 8, #1, NE—-TTHR. AGHS tio TUNA

B70 n Hk, UREA Na=aN, El NBREFR,
i “PY AB EE FRE AY file G AY AR.
\end{Verbatim}

\end{note}

\begin{note}[来源路径]

近世代数/近世代数/无标题/不变子群的中心 1cad2efa3f428043b357e2e9c7c9ce6b.md

\end{note}

\section*{不变子群的陪集叫商群}

\begin{note}[不变子群的陪集叫商群]

商群的阶是

\end{note}

\begin{note}[图片 OCR 摘录，待人工复核]
以下内容来自源图片识别，便于审核时对照：

\textbf{图片 1}（image.png）
\begin{Verbatim}[fontsize=\small]
EH3 —THFEPTRORRMF EDM EH RARER
— HF.

TEAR ATER EMRE. 1. ON. VaR.

1. oR,

I.     (zNyN) zN= [(ry) N] zN= (zryz) N

xN(yNzN) =xN((yz) N] = (xyz) N
lV.          eNaN= (er) N=2xN
V.         x 'NxN= (2x7!r) N=eN           iE st
EX EG INDABA BES BT TE BY Yi

TSAR. MERITS G/N Rem.
\end{Verbatim}

\end{note}

\begin{note}[来源路径]

近世代数/近世代数/无标题/不变子群的陪集叫商群 1cad2efa3f428038a9add2a529077225.md

\end{note}

\section*{交换群 阿贝尔群}

\begin{note}[来源上级主题]

群的基础知识 (\%E7\%BE\%A4\%E7\%9A\%84\%E5\%9F\%BA\%E7\%A1\%80\%E7\%9F\%A5\%E8\%AF\%86\%201aed2efa3f428009ac4eca74062880ff.md)

\end{note}

\begin{note}[交换群 阿贝尔群]

群的定义 +交换律=abel群

交换群就是证明ab=ba

\end{note}

\begin{note}[关联图片]
以下图片已纳入本条整理，当前版本优先依据 Markdown 文字整理，图片细节可在审核时对照：

1. image.png\\
2. image 1.png\\
3. image 2.png\\

\end{note}

\begin{note}[来源路径]

近世代数/近世代数/无标题/交换群 阿贝尔群 1aed2efa3f4280cdbde1e3cd83d58bc8.md

\end{note}

\section*{代数运算}

\begin{note}[来源上级主题]

基本概念 (\%E5\%9F\%BA\%E6\%9C\%AC\%E6\%A6\%82\%E5\%BF\%B5\%201c3d2efa3f4280648bdcd081c8508c78.md)

\end{note}

\begin{note}[代数运算]

子级 项目: 代数运算定义 (\%E4\%BB\%A3\%E6\%95\%B0\%E8\%BF\%90\%E7\%AE\%97\%E5\%AE\%9A\%E4\%B9\%89\%201a1d2efa3f42801bb126f612f365d13b.md), 运算表 (\%E8\%BF\%90\%E7\%AE\%97\%E8\%A1\%A8\%201a1d2efa3f4280498713e856fc016324.md), 封闭 二元运算 (\%E5\%B0\%81\%E9\%97\%AD\%20\%E4\%BA\%8C\%E5\%85\%83\%E8\%BF\%90\%E7\%AE\%97\%201a1d2efa3f4280c1abc9d4dc388cc80c.md), 结合律 交换律 分配律 (\%E7\%BB\%93\%E5\%90\%88\%E5\%BE\%8B\%20\%E4\%BA\%A4\%E6\%8D\%A2\%E5\%BE\%8B\%20\%E5\%88\%86\%E9\%85\%8D\%E5\%BE\%8B\%201a1d2efa3f428005bcfac88aa2af2b6e.md)

\end{note}

\begin{note}[来源路径]

近世代数/近世代数/无标题/代数运算 1aed2efa3f4280d597defe353a625299.md

\end{note}

\section*{代数运算定义}

\begin{note}[来源上级主题]

代数运算 (\%E4\%BB\%A3\%E6\%95\%B0\%E8\%BF\%90\%E7\%AE\%97\%201aed2efa3f4280d597defe353a625299.md)

\end{note}

\begin{definition}[代数运算定义]

两个集合的积到D的映射叫代数运算

例子

（a，b）=A*B A*B得到的不是a/b 对应的是a/b

\end{definition}

\begin{note}[关联图片]
以下图片已纳入本条整理，当前版本优先依据 Markdown 文字整理，图片细节可在审核时对照：

1. image.png\\
2. image 1.png\\

\end{note}

\begin{note}[来源路径]

近世代数/近世代数/无标题/代数运算定义 1a1d2efa3f42801bb126f612f365d13b.md

\end{note}

\section*{例子}

\begin{note}[来源上级主题]

子群的陪集 (\%E5\%AD\%90\%E7\%BE\%A4\%E7\%9A\%84\%E9\%99\%AA\%E9\%9B\%86\%201c3d2efa3f42806a85b0f0113028b207.md)

\end{note}

\begin{example}[例子]

循环置换

（13）代表一种映射123——321加入（12）变成312 3-2 2-1 1-3

\end{example}

\begin{note}[关联图片]
以下图片已纳入本条整理，当前版本优先依据 Markdown 文字整理，图片细节可在审核时对照：

1. image.png\\

\end{note}

\begin{note}[来源路径]

近世代数/近世代数/无标题/例子 1c4d2efa3f428049872dd61c51dda43e.md

\end{note}

\section*{关系的定义 二元关系}

\begin{note}[来源上级主题]

分类与关系 (\%E5\%88\%86\%E7\%B1\%BB\%E4\%B8\%8E\%E5\%85\%B3\%E7\%B3\%BB\%201aed2efa3f42802491cdeb2b56b5f1f7.md)

\end{note}

\begin{definition}[关系的定义 二元关系]

二元关系 每一个映射都叫做A的一个关系 D中对错是固定的集合 D不一定要

S 的每个划分都以自然的方式产生一个关系 R，即对于 x，y ∈ S，当且仅当 x 和 y 处于划分的同一个单元时，xRy。在集合表示法中，我们将写成 xRy，表示 （x, y） ∈ R      relation

R可以代表所有比较符号甚至定义的关系；

R也可以是A*A的一个子集；

我们引入了集合 A 的元素与集合 B 的元素相关的概念，我们可能会用 aRb 表示。符号 aRb 展示了元素 a 和元素 b 从左到右的顺序，就像 A×B 中元素 （a, b） 的顺序一样。这使我们对关系 R 作为集的定义如下；

从定义上来将aRb的意思是（a，b）属于集合R

\end{definition}

\begin{note}[关联图片]
以下图片已纳入本条整理，当前版本优先依据 Markdown 文字整理，图片细节可在审核时对照：

1. image.png\\
2. IMG\_2643.jpg\\

\end{note}

\begin{note}[来源路径]

近世代数/近世代数/无标题/关系的定义 二元关系 1a7d2efa3f42802f89cbe2c78b33686d.md

\end{note}

\section*{凯莱定理任何群都是同一个群的同构}

\begin{note}[来源上级主题]

变换群 (\%E5\%8F\%98\%E6\%8D\%A2\%E7\%BE\%A4\%201c3d2efa3f4280a3a416f8493ad4dae2.md)

\end{note}

\begin{theorem}[凯莱定理任何群都是同一个群的同构]

transform group置换群

第一步构造映射g到gx这个操作叫做taox

每一个元都有一种变换 再构建变换和元素之间的映射

x对应的是

前面的假定也是taox的一部分

gx之间是由定义的运算连接的而不是乘法

同构

这步证明的是g（xoy）=g（x）og（y）

变换群

\end{theorem}

\begin{note}[关联图片]
以下图片已纳入本条整理，当前版本优先依据 Markdown 文字整理，图片细节可在审核时对照：

1. image.png\\
2. image 1.png\\
3. image 2.png\\
4. image 3.png\\

\end{note}

\begin{note}[来源路径]

近世代数/近世代数/无标题/凯莱定理任何群都是同一个群的同构 1b5d2efa3f42805cb494c5d6269c2742.md

\end{note}

\section*{分类}

\begin{note}[来源上级主题]

分类与关系 (\%E5\%88\%86\%E7\%B1\%BB\%E4\%B8\%8E\%E5\%85\%B3\%E7\%B3\%BB\%201aed2efa3f42802491cdeb2b56b5f1f7.md)

\end{note}

\begin{note}[分类]

1.A=UAi  2.Ai非空    3.任意Ai和Aj交集是空

分类是所有类的全体

类里面一个元素就是代表

一个子集叫一个类     所有属于一个类的元素组成的集合叫做A的一个分类

我们可以把 0 看作是能被 2 整除的整数，而 1 则是当被 2 整除时余数为 1 的整数。这个概念同样适用于 2 以外的正整数。例如，我们可以把 Z 分割成三个单元：

\end{note}

\begin{note}[关联图片]
以下图片已纳入本条整理，当前版本优先依据 Markdown 文字整理，图片细节可在审核时对照：

1. image.png\\
2. image 1.png\\

\end{note}

\begin{note}[来源路径]

近世代数/近世代数/无标题/分类 1a7d2efa3f428051977beee16bb06dae.md

\end{note}

\section*{分类与关系}

\begin{note}[来源上级主题]

基本概念 (\%E5\%9F\%BA\%E6\%9C\%AC\%E6\%A6\%82\%E5\%BF\%B5\%201c3d2efa3f4280648bdcd081c8508c78.md)

\end{note}

\begin{note}[分类与关系]

子级 项目: 关系的定义 二元关系 (\%E5\%85\%B3\%E7\%B3\%BB\%E7\%9A\%84\%E5\%AE\%9A\%E4\%B9\%89\%20\%E4\%BA\%8C\%E5\%85\%83\%E5\%85\%B3\%E7\%B3\%BB\%201a7d2efa3f42802f89cbe2c78b33686d.md), 分类 (\%E5\%88\%86\%E7\%B1\%BB\%201a7d2efa3f428051977beee16bb06dae.md), 等价关系 (\%E7\%AD\%89\%E4\%BB\%B7\%E5\%85\%B3\%E7\%B3\%BB\%201a7d2efa3f4280a0ae90db267cbf3328.md), 每一个分类都决定了A的一个等价关系 (\%E6\%AF\%8F\%E4\%B8\%80\%E4\%B8\%AA\%E5\%88\%86\%E7\%B1\%BB\%E9\%83\%BD\%E5\%86\%B3\%E5\%AE\%9A\%E4\%BA\%86A\%E7\%9A\%84\%E4\%B8\%80\%E4\%B8\%AA\%E7\%AD\%89\%E4\%BB\%B7\%E5\%85\%B3\%E7\%B3\%BB\%201a7d2efa3f4280eeabf8d5f851e06cc2.md), 集合A的任一个等价关系都可确定A的一个分类 (\%E9\%9B\%86\%E5\%90\%88A\%E7\%9A\%84\%E4\%BB\%BB\%E4\%B8\%80\%E4\%B8\%AA\%E7\%AD\%89\%E4\%BB\%B7\%E5\%85\%B3\%E7\%B3\%BB\%E9\%83\%BD\%E5\%8F\%AF\%E7\%A1\%AE\%E5\%AE\%9AA\%E7\%9A\%84\%E4\%B8\%80\%E4\%B8\%AA\%E5\%88\%86\%E7\%B1\%BB\%201a8d2efa3f4280d68789ddcfda4d73ac.md), 同余关系 模n的剩余类 (\%E5\%90\%8C\%E4\%BD\%99\%E5\%85\%B3\%E7\%B3\%BB\%20\%E6\%A8\%A1n\%E7\%9A\%84\%E5\%89\%A9\%E4\%BD\%99\%E7\%B1\%BB\%201a8d2efa3f4280ffbd5ffb8a8131fb1b.md)

\end{note}

\begin{note}[来源路径]

近世代数/近世代数/无标题/分类与关系 1aed2efa3f42802491cdeb2b56b5f1f7.md

\end{note}

\section*{剩余类加群}

\begin{note}[来源上级主题]

群的基础知识 (\%E7\%BE\%A4\%E7\%9A\%84\%E5\%9F\%BA\%E7\%A1\%80\%E7\%9F\%A5\%E8\%AF\%86\%201aed2efa3f428009ac4eca74062880ff.md)

\end{note}

\begin{note}[剩余类加群]

剩余类加群有A必有n-A

\end{note}

\begin{note}[图片 OCR 摘录，待人工复核]
以下内容来自源图片识别，便于审核时对照：

\textbf{图片 1}（image.png）
\begin{Verbatim}[fontsize=\small]
Z, = {(0],[11,[2],--,[ — 11}
bie [a]+[o]=[a+ 5]
PARI AAI :
*[a]=[4'],[4]=[4']

=> [a]+[o]=[a+5],[a']+[o']=[a'+5"]
fg UE AY [¢+ 5] =[a'+ 5"
“nja—a',n|b—b' .ni|a-—a'+b—b'
n|(a+b)—(a'+b') => [a+b] =[a'+5']
\end{Verbatim}

\textbf{图片 2}（image 1.png）
\begin{Verbatim}[fontsize=\small]
(lL) 28s
(2) GRACE
(a + [4]) + [c] oe [a +b+ c] a [a] + ([5] + [e])
(3) BAAS ft7c [0]
(]+[4]=[4]
(4) @#zB I [-4]
[-2]+[4]=[9]
TEER, PERN RRA.
ive (Z.,+) BUICE Z,
\end{Verbatim}

\end{note}

\begin{note}[来源路径]

近世代数/近世代数/无标题/剩余类加群 1aed2efa3f4280428135f1eab5e513d0.md

\end{note}

\section*{加群定义}

\begin{note}[来源上级主题]

第三章 (\%E7\%AC\%AC\%E4\%B8\%89\%E7\%AB\%A0\%201ecd2efa3f4280c7a6b0ea0de416c09f.md)

\end{note}

\begin{definition}[加群定义]

分割

\end{definition}

\begin{note}[图片 OCR 摘录，待人工复核]
以下内容来自源图片识别，便于审核时对照：

\textbf{图片 1}（image.png）
\begin{Verbatim}[fontsize=\small]
EM — PACA — PIE, TR RT a RBs FY
%, FAAS +RRM.
\end{Verbatim}

\textbf{图片 2}（image 1.png）
\begin{Verbatim}[fontsize=\small]
Hi BWAZH SREB AHN RBS Fy ik ST FAA A. A
KEMREBRNABERAA SHOE. FSW, FSITRMMNERYR
th BR a BS. BLE FR TT fai A HP A PS PT A.

AFMBHMEBAAAE, nS3Ga1> Aer **%s An NMA R SM, Pd
MRAVER AES Dla, HAR,

Sa, S04, ag “TH 2% +P Bae

— 7S Dn FF BS ME — A BA OR OO HRN, FARE UMS. RNAW
PISA :

Ota=at0=a (a HERI).                 (1)

Toa WHE MMR INA —a KRM, HAE MMe HAT MRM a).
Ttat(—bd) RNS mab Sma WM). Aix SE MURR HM,
\end{Verbatim}

\end{note}

\begin{note}[来源路径]

近世代数/近世代数/无标题/加群定义 1ecd2efa3f4280c5a839d3c295d23a6b.md

\end{note}

\section*{单位元}

\begin{note}[来源上级主题]

群的基础知识 (\%E7\%BE\%A4\%E7\%9A\%84\%E5\%9F\%BA\%E7\%A1\%80\%E7\%9F\%A5\%E8\%AF\%86\%201aed2efa3f428009ac4eca74062880ff.md)

\end{note}

\begin{note}[单位元]

在代数中，一个集合上的操作可能具有另一个重要的性质。0 和 1 作为实数具有特殊的作用，因为对于任何实数 a，a + 0 = a 和 a × 1 = a。由于这些性质，0 被称为 R 中加法单位元，1 被称为 R 中乘法单位元。通常，我们有以下单位元的定义。

定义 LetS是一个具有二元运算∗的集合。如果对于集合S中的所有a，a∗e=e∗a=a，那么e被称为是二元运算∗的单位元

定理（恒等式的唯一性）具有二元运算*的集合最多有一个单位元

示例 继续使用示例 1.7，设 F 为所有将实数映射到实数的函数的集合。我们验证由 ι（x） = x 定义的函数是 操作函数复合的恒等函数。设 f ∈ F。则 f ◦ ι（x） = f（ι（x）） = f（x） 和 ι◦ f（x） = ι（f（x）） = f（x）。 函数 m（x） = 1 是操作函数乘法的恒等函数，a（x） = 0 是函数加法的恒等函数，但函数减法没有恒等元素。

\end{note}

\begin{note}[来源路径]

近世代数/近世代数/无标题/单位元 1aad2efa3f42804aa677fdb25912917d.md

\end{note}

\section*{单位元唯一性证明}

\begin{note}[来源上级主题]

幺半群的定义 定义1.2 定义1.3 (\%E5\%B9\%BA\%E5\%8D\%8A\%E7\%BE\%A4\%E7\%9A\%84\%E5\%AE\%9A\%E4\%B9\%89\%20\%E5\%AE\%9A\%E4\%B9\%891\%202\%20\%E5\%AE\%9A\%E4\%B9\%891\%203\%201cad2efa3f42800ea08fdf383be45a9a.md)

\end{note}

\begin{proof}

数学中唯一性证明通常设存在两个再证明这两个相等

\end{proof}

\begin{note}[来源路径]

近世代数/近世代数/无标题/单位元唯一性证明 1cad2efa3f428021acabe9e17d5b39d5.md

\end{note}

\section*{单射满射}

\begin{note}[来源上级主题]

映射与集合基础知识 (\%E6\%98\%A0\%E5\%B0\%84\%E4\%B8\%8E\%E9\%9B\%86\%E5\%90\%88\%E5\%9F\%BA\%E7\%A1\%80\%E7\%9F\%A5\%E8\%AF\%86\%201aed2efa3f42800c8261da73ac826269.md)

\end{note}

\begin{note}[单射满射]

https://www.notion.so

单射就是原像不同 像不同

映射 原像相等 像相等  像相等原像不一定相等

单射 原像不等 像一定不等 像相等原像一定相等

A到A的映射叫做A的一个变换

\end{note}

\begin{note}[关联图片]
以下图片已纳入本条整理，当前版本优先依据 Markdown 文字整理，图片细节可在审核时对照：

1. image.png\\

\end{note}

\begin{note}[来源路径]

近世代数/近世代数/无标题/单射满射 1a1d2efa3f4280a7ab9bf96427741ce6.md

\end{note}

\section*{变化群的意义}

\begin{note}[来源上级主题]

变换群 (\%E5\%8F\%98\%E6\%8D\%A2\%E7\%BE\%A4\%201c3d2efa3f4280a3a416f8493ad4dae2.md)

\end{note}

\begin{note}[变化群的意义]

可以通过一个简单的例子来理解。设想你有一个正方形，你可以对这个正方形进行各种操作（旋转、翻转等），这些操作我们称为“变换”。现在考虑以下几点：

1.	变换的组合（封闭性）

如果你先对正方形进行一次变换，再进行另一次变换，整体效果依然是一个变换。比如先旋转90°再旋转90°，结果等于旋转180°。这说明当你把两个变换组合在一起时，结果仍然是一个有效的变换。

2.	顺序无关“括号”（结合律）

对于三个变换来说，不论你先组合哪两个，最后的结果是相同的。也就是说，如果你有变换 f、g 和 h，那么 (f \textbackslash{}circ g) \textbackslash{}circ h 与 f \textbackslash{}circ (g \textbackslash{}circ h) 得到的效果相同。这与我们日常操作中“先后顺序”的某些情况相似，只要变换的顺序保持不变，组合的方式（如何分组）不会改变最终结果。

3.	存在一个“无操作”的变换（单位元）

在所有变换中，有一种变换什么也不做（称为单位元），比如旋转0°或者“不翻转”的操作。无论你先做任何变换，再做这个“无操作”的变换，或者反过来，效果都是原来的变换不变。

4.	每个变换都有一个“撤销”动作（逆元）

对于任意一次变换，都可以找到另一个变换来把它“撤销”，使得整体效果恢复到最初状态。比如，如果你旋转正方形90°，那么再旋转270°（或者旋转-90°）就能把正方形转回原来的位置。这个“撤销”的变换就是原变换的逆元。

总的来说，一个群就是满足这四个性质的集合。在这里，每一个对正方形的变换（例如旋转、翻转）都可以看成群的一个元素，它们之间的“运算”就是把两个变换连续执行。因为这些变换组合起来总能得到另一个变换，而且有“无操作”和“撤销操作”，所以它们构成了一个群，通常称为对称群或变换群。

\end{note}

\begin{note}[来源路径]

近世代数/近世代数/无标题/变化群的意义 1b7d2efa3f4280e3baffc1d56b67f68e.md

\end{note}

\section*{变换群}

\begin{note}[来源上级主题]

变换群 (\%E5\%8F\%98\%E6\%8D\%A2\%E7\%BE\%A4\%201c3d2efa3f4280a3a416f8493ad4dae2.md)

\end{note}

\begin{note}[变换群]

子级 项目: 无标题 (\%E6\%97\%A0\%E6\%A0\%87\%E9\%A2\%98\%201bad2efa3f42803389efdacc7eedffb9.md), 无标题 (\%E6\%97\%A0\%E6\%A0\%87\%E9\%A2\%98\%201bad2efa3f42802aa86ec37cd6544448.md)

变换得自己往自己身上映射

变换乘法 恒等变换     变换群里面的元素是变换

变换群里面都是变换

线性变换多项式

\end{note}

\begin{note}[关联图片]
以下图片已纳入本条整理，当前版本优先依据 Markdown 文字整理，图片细节可在审核时对照：

1. image.png\\
2. image 1.png\\
3. image 2.png\\

\end{note}

\begin{note}[来源路径]

近世代数/近世代数/无标题/变换群 1b5d2efa3f4280bc8c6bf24cc87fed3c.md

\end{note}

\section*{变换群}

\begin{note}[来源上级主题]

群论 (\%E7\%BE\%A4\%E8\%AE\%BA\%201c3d2efa3f428064a77dd3ff465d3548.md)

\end{note}

\begin{note}[变换群]

子级 项目: 构造同态同构 (\%E6\%9E\%84\%E9\%80\%A0\%E5\%90\%8C\%E6\%80\%81\%E5\%90\%8C\%E6\%9E\%84\%201b5d2efa3f42802ba396ef3b35139369.md), 变换群 (\%E5\%8F\%98\%E6\%8D\%A2\%E7\%BE\%A4\%201b5d2efa3f4280bc8c6bf24cc87fed3c.md), 变换群证明 (\%E5\%8F\%98\%E6\%8D\%A2\%E7\%BE\%A4\%E8\%AF\%81\%E6\%98\%8E\%201b5d2efa3f42802e858cebf98bd17fd6.md), 凯莱定理任何群都是同一个群的同构 (\%E5\%87\%AF\%E8\%8E\%B1\%E5\%AE\%9A\%E7\%90\%86\%E4\%BB\%BB\%E4\%BD\%95\%E7\%BE\%A4\%E9\%83\%BD\%E6\%98\%AF\%E5\%90\%8C\%E4\%B8\%80\%E4\%B8\%AA\%E7\%BE\%A4\%E7\%9A\%84\%E5\%90\%8C\%E6\%9E\%84\%201b5d2efa3f42805cb494c5d6269c2742.md), 变化群的意义 (\%E5\%8F\%98\%E5\%8C\%96\%E7\%BE\%A4\%E7\%9A\%84\%E6\%84\%8F\%E4\%B9\%89\%201b7d2efa3f4280e3baffc1d56b67f68e.md)

\end{note}

\begin{note}[来源路径]

近世代数/近世代数/无标题/变换群 1c3d2efa3f4280a3a416f8493ad4dae2.md

\end{note}

\section*{变换群证明}

\begin{note}[来源上级主题]

变换群 (\%E5\%8F\%98\%E6\%8D\%A2\%E7\%BE\%A4\%201c3d2efa3f4280a3a416f8493ad4dae2.md)

\end{note}

\begin{proof}

子级 项目: 无标题 (\%E6\%97\%A0\%E6\%A0\%87\%E9\%A2\%98\%201b7d2efa3f42802dbdaecc7643526bbe.md)

逆元

变化乘法就规定了不同变化之间的连接方式

变换群 G是群所以G的每个元素必须有逆元，运算封闭还有单位元

\end{proof}

\begin{note}[关联图片]
以下图片已纳入本条整理，当前版本优先依据 Markdown 文字整理，图片细节可在审核时对照：

1. image.png\\
2. image 1.png\\

\end{note}

\begin{note}[来源路径]

近世代数/近世代数/无标题/变换群证明 1b5d2efa3f42802e858cebf98bd17fd6.md

\end{note}

\section*{只要不相连就可以快速求阶数而且有交换律}

\begin{note}[来源上级主题]

置换群 (\%E7\%BD\%AE\%E6\%8D\%A2\%E7\%BE\%A4\%201bad2efa3f4280228c17c29195a2c525.md)

\end{note}

\begin{note}[只要不相连就可以快速求阶数而且有交换律]

pr bz bon CLL) (¥ )

gk) =

\end{note}

\begin{note}[图片 OCR 摘录，待人工复核]
以下内容来自源图片识别，便于审核时对照：

\textbf{图片 1}（image.png）
\begin{Verbatim}[fontsize=\small]
pr bz bon CLL) (¥ )
gk) =
\end{Verbatim}

\end{note}

\begin{note}[来源路径]

近世代数/近世代数/无标题/只要不相连就可以快速求阶数而且有交换律 1bcd2efa3f4280e9add2d922e72a6f5d.md

\end{note}

\section*{右陪集的定义}

\begin{note}[来源上级主题]

子群的陪集 (\%E5\%AD\%90\%E7\%BE\%A4\%E7\%9A\%84\%E9\%99\%AA\%E9\%9B\%86\%201c3d2efa3f42806a85b0f0113028b207.md)

\end{note}

\begin{definition}[右陪集的定义]

子群 子群类比子集，右陪集是一种类是满足等价关系的元素的集合 等价关系是定义的不是等于

陪集是陪伴着出现的没有对象就不存在陪伴

a和b满足这个关系当且仅当ab逆在H子群中

下面的三条是证明这个关系是等价关系 定义的时候就是H中每一个元素乘a，这里的乘是抽象的

Ha就是H中每一个元素乘a aH一般不是群但是可以与H构成双射

\end{definition}

\begin{note}[关联图片]
以下图片已纳入本条整理，当前版本优先依据 Markdown 文字整理，图片细节可在审核时对照：

1. image.png\\
2. image 1.png\\
3. image 2.png\\

\end{note}

\begin{note}[来源路径]

近世代数/近世代数/无标题/右陪集的定义 1c3d2efa3f428044a7f6c58dda5da09f.md

\end{note}

\section*{同余关系 模n的剩余类}

\begin{note}[来源上级主题]

分类与关系 (\%E5\%88\%86\%E7\%B1\%BB\%E4\%B8\%8E\%E5\%85\%B3\%E7\%B3\%BB\%201aed2efa3f42802491cdeb2b56b5f1f7.md)

\end{note}

\begin{note}[同余关系 模n的剩余类]

a-b能整除n说明a/n和b/n的余数相同

\end{note}

\begin{note}[图片 OCR 摘录，待人工复核]
以下内容来自源图片识别，便于审核时对照：

\textbf{图片 1}（image.png）
\begin{Verbatim}[fontsize=\small]
EM ERM O<neZ, WUE Z
Pie —MSiKA
R:Va,beZ,aRbs nla —b
Ma RAR n WERAR, #48 aRb
itt a=b(n) = (ala RbtRn)
MHAARAARENDRPHWRAR 1
RI.
\end{Verbatim}

\end{note}

\begin{note}[来源路径]

近世代数/近世代数/无标题/同余关系 模n的剩余类 1a8d2efa3f4280ffbd5ffb8a8131fb1b.md

\end{note}

\section*{同态同构的意义}

\begin{note}[来源上级主题]

群的同态同构 (\%E7\%BE\%A4\%E7\%9A\%84\%E5\%90\%8C\%E6\%80\%81\%E5\%90\%8C\%E6\%9E\%84\%201c3d2efa3f428076a462f0397fbfd469.md)

\end{note}

\begin{note}[同态同构的意义]

同态保持运算，同构完全一致

\end{note}

\begin{note}[图片 OCR 摘录，待人工复核]
以下内容来自源图片识别，便于审核时对照：

\textbf{图片 1}（IMG\_2903.jpeg）
\begin{Verbatim}[fontsize=\small]
NikHE, FTAA EH. SE, TEA ge “S|? BT AT AA cS HAAS, SEE, te
HY DAMME A 0 8 | LR A RTE PPE FT TD, FES BES ABT VAS | LHP ARTE.
ANSE “TRAY”. LE, BRERA A DU pea]. WARE ASH MUL
EMAC , MAR MM CesEe— BUS. SEE — PY, RACE BT VA ER [i] FE
TD, BVT DAS ho eS EE es TEU LE, AI ee, PREP, TEAR)
——thtee, ROY Sa (LEM ARAB) POME—DB Ste “PAE AN TA]” .
\end{Verbatim}

\end{note}

\begin{note}[来源路径]

近世代数/近世代数/无标题/同态同构的意义 1cad2efa3f42801b93c4f525f7a94957.md

\end{note}

\section*{同态的定义}

\begin{note}[来源上级主题]

映射与集合基础知识 (\%E6\%98\%A0\%E5\%B0\%84\%E4\%B8\%8E\%E9\%9B\%86\%E5\%90\%88\%E5\%9F\%BA\%E7\%A1\%80\%E7\%9F\%A5\%E8\%AF\%86\%201aed2efa3f42800c8261da73ac826269.md)

\end{note}

\begin{definition}[同态的定义]

同态映射  A这边对两个元素操作A得到的 和A’中华对这两个元素在A‘中的映射操作A 得到的还是映射   鹿丸的影子模仿

同态映射+满射出呢说是同态满射 同态满射才能说A与A‘是同态

同态映射可以非单射，非满射

xy对应的应该是2xy但是A进行操作后是4xy

同态传递结合律分配律交换律

运算符也可以说同态

\end{definition}

\begin{note}[关联图片]
以下图片已纳入本条整理，当前版本优先依据 Markdown 文字整理，图片细节可在审核时对照：

1. IMG\_2919.jpeg\\
2. image.png\\

\end{note}

\begin{note}[来源路径]

近世代数/近世代数/无标题/同态的定义 1a1d2efa3f428051bf6ddd4a0bc015cf.md

\end{note}

\section*{同态的意义}

\begin{note}[来源上级主题]

群的同态同构 (\%E7\%BE\%A4\%E7\%9A\%84\%E5\%90\%8C\%E6\%80\%81\%E5\%90\%8C\%E6\%9E\%84\%201c3d2efa3f428076a462f0397fbfd469.md)

\end{note}

\begin{note}[同态的意义]

群就是要找单位元逆元还有运算封闭

同态已经说明两个群有很多性质是一样的很多东西都能找到对应，同构更是严格表明他们的相似，同构同态就是通过群之间的特殊关系通过拉以知群的线头把未知群需要的东西来出来。

\end{note}

\begin{note}[关联图片]
以下图片已纳入本条整理，当前版本优先依据 Markdown 文字整理，图片细节可在审核时对照：

1. image.png\\
2. image 1.png\\

\end{note}

\begin{note}[来源路径]

近世代数/近世代数/无标题/同态的意义 1b5d2efa3f42801ba4aed9823eb1a0ae.md

\end{note}

\section*{同构}

\begin{note}[来源上级主题]

映射与集合基础知识 (\%E6\%98\%A0\%E5\%B0\%84\%E4\%B8\%8E\%E9\%9B\%86\%E5\%90\%88\%E5\%9F\%BA\%E7\%A1\%80\%E7\%9F\%A5\%E8\%AF\%86\%201aed2efa3f42800c8261da73ac826269.md)

\end{note}

\begin{note}[同构]

自同构

同态的定义

同构=同态+限制必须一一对应既单又满

单和满看能不能一对多，还有逆向也是映射满射

同构只看运算规律，把复杂集合的运算同构简化

自己到自己叫自同构

证明任意ϕ（aob）=ϕ（c）=d

ϕ（a）oϕ（b）=d

假设不是同构不能用同态构造矛盾，特殊点就是切入口

单射反证就是一对多

定义11线性空间的同构

\end{note}

\begin{note}[关联图片]
以下图片已纳入本条整理，当前版本优先依据 Markdown 文字整理，图片细节可在审核时对照：

1. image.png\\
2. image 1.png\\

\end{note}

\begin{note}[来源路径]

近世代数/近世代数/无标题/同构 1a7d2efa3f4280a5b953fd47d4b69430.md

\end{note}

\section*{命题1.11证明}

\begin{note}[来源上级主题]

定义1.13定义子群根本目的是证明它是一个群 (\%E5\%AE\%9A\%E4\%B9\%891\%2013\%E5\%AE\%9A\%E4\%B9\%89\%E5\%AD\%90\%E7\%BE\%A4\%E6\%A0\%B9\%E6\%9C\%AC\%E7\%9B\%AE\%E7\%9A\%84\%E6\%98\%AF\%E8\%AF\%81\%E6\%98\%8E\%E5\%AE\%83\%E6\%98\%AF\%E4\%B8\%80\%E4\%B8\%AA\%E7\%BE\%A4\%201cbd2efa3f4280ee903ac1e652f44dec.md)

\end{note}

\begin{proof}

第二个条件x可以是e，y则是H中任意元素证明了每一个都有逆元再证明x.y属于H即可

\end{proof}

\begin{note}[关联图片]
以下图片已纳入本条整理，当前版本优先依据 Markdown 文字整理，图片细节可在审核时对照：

1. IMG\_2916.jpeg\\

\end{note}

\begin{note}[来源路径]

近世代数/近世代数/无标题/命题1 11证明 1cbd2efa3f4280679acbe2813bc3b4e4.md

\end{note}

\section*{命题1.12特殊线性群是群}

\begin{note}[来源上级主题]

定义1.13定义子群根本目的是证明它是一个群 (\%E5\%AE\%9A\%E4\%B9\%891\%2013\%E5\%AE\%9A\%E4\%B9\%89\%E5\%AD\%90\%E7\%BE\%A4\%E6\%A0\%B9\%E6\%9C\%AC\%E7\%9B\%AE\%E7\%9A\%84\%E6\%98\%AF\%E8\%AF\%81\%E6\%98\%8E\%E5\%AE\%83\%E6\%98\%AF\%E4\%B8\%80\%E4\%B8\%AA\%E7\%BE\%A4\%201cbd2efa3f4280ee903ac1e652f44dec.md)

\end{note}

\begin{proposition}[命题1.12特殊线性群是群]

因为它的定义就是行列式是1，1*1永远是1所以封闭性极好逆元也是自己啥都有了

\end{proposition}

\begin{note}[关联图片]
以下图片已纳入本条整理，当前版本优先依据 Markdown 文字整理，图片细节可在审核时对照：

1. IMG\_2917.jpeg\\

\end{note}

\begin{note}[来源路径]

近世代数/近世代数/无标题/命题1 12特殊线性群是群 1cbd2efa3f4280bc9d1cfc2496d61576.md

\end{note}

\section*{命题1.13群同态单位元对单位元逆元对逆元}

\begin{note}[来源上级主题]

群的同态同构 (\%E7\%BE\%A4\%E7\%9A\%84\%E5\%90\%8C\%E6\%80\%81\%E5\%90\%8C\%E6\%9E\%84\%201b5d2efa3f42803aaaa7f6c2fdcfd0b4.md)

\end{note}

\begin{proposition}[命题1.13群同态单位元对单位元逆元对逆元]

证明思路因为e.e=e所以要添项减少项目的时候e就是很有用e也可以拆成两个互逆的元素相乘把我们想要的东西引入

定义1.5幺半群的同态 同态保持所有运算并且把特殊的元素映射到特殊的元素上面

\end{proposition}

\begin{note}[关联图片]
以下图片已纳入本条整理，当前版本优先依据 Markdown 文字整理，图片细节可在审核时对照：

1. IMG\_2920.jpeg\\

\end{note}

\begin{note}[来源路径]

近世代数/近世代数/无标题/命题1 13群同态单位元对单位元逆元对逆元 1cbd2efa3f4280048341d86cd285853a.md

\end{note}

\section*{命题1.14}

\begin{note}[来源上级主题]

群的同态同构 (\%E7\%BE\%A4\%E7\%9A\%84\%E5\%90\%8C\%E6\%80\%81\%E5\%90\%8C\%E6\%9E\%84\%201c3d2efa3f428076a462f0397fbfd469.md)

\end{note}

\begin{proposition}[命题1.14]

命题1.12特殊线性群是群

\end{proposition}

\begin{note}[图片 OCR 摘录，待人工复核]
以下内容来自源图片识别，便于审核时对照：

\textbf{图片 1}（IMG\_2923.jpeg）
\begin{Verbatim}[fontsize=\small]
PEP R, FAT FEE BURA ELT APIS TERS, FEUEHW Ee PE, FRAT
FTN PRY ASHE. TMT 1 MARCA Lc, PERRI TT SS I PE BEY (FRE )
HEA AS 6

fit jel 1.14

DOA PE CR ESE , FECA, FAT AZS EATEN. AB ARPA VEREE Te? “EP ee ait ae Eb AB 4
DRE LTC Ca. (HERE) PASE (PH). ABA—PROR UL, AEA EAE IALAS RB AA cH MA
eS e? HAE REN. MAARRINAAGE, RAMRAA CE TH, MAIER TH, teat ST
SZ, TRASH, ER HE. TS PE, FATT A AOR is BEY
S#, RSZIBMH AA, PEE SRM AIT HZ. PIRI ZA BE.
\end{Verbatim}

\end{note}

\begin{note}[来源路径]

近世代数/近世代数/无标题/命题1 14 1cbd2efa3f4280569ab4d1fcc291e1bc.md

\end{note}

\section*{命题1.17同态就是拆括号的游戏}

\begin{note}[来源上级主题]

1.2群 (1\%202\%E7\%BE\%A4\%201ecd2efa3f4280cf8df9c881d186f2a4.md)

\end{note}

\begin{proposition}[命题1.17同态就是拆括号的游戏]

同态就是拆括号合括号的游戏

\end{proposition}

\begin{note}[图片 OCR 摘录，待人工复核]
以下内容来自源图片识别，便于审核时对照：

\textbf{图片 1}（IMG\_2939.jpeg）
\begin{Verbatim}[fontsize=\small]
4H 1.17                                            —
EW AA fo EH, PRT RAE fo! ZBAA. Oxy’ € GC’, Mix’ = f(x) = fly).
x'xy=f(x-y), KS TO *y)=x-yaf la): fly). RRART EW.

FOF EL BA, BEB OR”, +), (C",+) RENOIR. FRAT MAT, DOR EAVTREEE SIC FH RNY FB RFA
VERN, POAT. BRA ARF HE, BEATA AT DAE EATS FR RE LR EZ PE? 8
AVES AR, top eT DSS RRL. EIGEN. BAVA MA IPI.
\end{Verbatim}

\end{note}

\begin{note}[来源路径]

近世代数/近世代数/无标题/命题1 17同态就是拆括号的游戏 1d1d2efa3f42800d9bf2efcc79be14a6.md

\end{note}

\section*{命题1.3.3子群陪集要不然相等要不然无交}

\begin{note}[来源上级主题]

1.3有限群 (1\%203\%E6\%9C\%89\%E9\%99\%90\%E7\%BE\%A4\%201ead2efa3f4280c39a49d8b0af9aec03.md)

\end{note}

\begin{proposition}[命题1.3.3子群陪集要不然相等要不然无交]

a1=bh2*h1逆 两边都能加h

比如说剩余类加群\{0,1,2,3,4,5\}子群，有1阶的就是\{0\} ，2阶\{0,3\}生成元是3 a1=2则a1H=\{2,5\}

b1=1 b1H=\{1,4\}   因为这些陪集都互不相交所以他们总能最终填满整个G

\end{proposition}

\begin{note}[关联图片]
以下图片已纳入本条整理，当前版本优先依据 Markdown 文字整理，图片细节可在审核时对照：

1. IMG\_2981.jpeg\\
2. IMG\_2983.jpeg\\

\end{note}

\begin{note}[来源路径]

近世代数/近世代数/无标题/命题1 3 3子群陪集要不然相等要不然无交 1d2d2efa3f42806da9c1eaa41dc2aa66.md

\end{note}

\section*{命题1.34 1.35嵌套子群的配集}

\begin{note}[来源上级主题]

1.3有限群 (1\%203\%E6\%9C\%89\%E9\%99\%90\%E7\%BE\%A4\%201ead2efa3f4280c39a49d8b0af9aec03.md)

\end{note}

\begin{proposition}[命题1.34 1.35嵌套子群的配集]

<是子群的意思

拉格朗日定理

证明1 ｜G｜/｜H｜*｜H｜/｜E｜

证明2 I J表示的就是能让他们形成不相交配集的代表元，需要证明两件事一件事是枚举了所有配集也就能覆盖G，第二件事就是证明集合两两不交

商群性质

不变子群的陪集叫商群

\end{proposition}

\begin{note}[关联图片]
以下图片已纳入本条整理，当前版本优先依据 Markdown 文字整理，图片细节可在审核时对照：

1. IMG\_2988.jpeg\\

\end{note}

\begin{note}[来源路径]

近世代数/近世代数/无标题/命题1 34 1 35嵌套子群的配集 1d2d2efa3f428089be3ecd9b4c4dbf6c.md

\end{note}

\section*{命题1.36良定义}

\begin{note}[来源上级主题]

1.4正规子群 (1\%204\%E6\%AD\%A3\%E8\%A7\%84\%E5\%AD\%90\%E7\%BE\%A4\%201ead2efa3f428095a412e65ae0862a5d.md)

\end{note}

\begin{proposition}[命题1.36良定义]

正规子群就是良定义，良定义在数学里面就是不同表达方式下的取值说一样的，和选取的代表元无关

\end{proposition}

\begin{note}[关联图片]
以下图片已纳入本条整理，当前版本优先依据 Markdown 文字整理，图片细节可在审核时对照：

1. image.png\\

\end{note}

\begin{note}[来源路径]

近世代数/近世代数/无标题/命题1 36良定义 1d8d2efa3f4280f894b7ef35415e3855.md

\end{note}

\section*{命题1.37 商群}

\begin{note}[来源上级主题]

1.4正规子群 (1\%204\%E6\%AD\%A3\%E8\%A7\%84\%E5\%AD\%90\%E7\%BE\%A4\%201ead2efa3f428095a412e65ae0862a5d.md)

\end{note}

\begin{proposition}[命题1.37 商群]

命题1.36良定义 N是正规子群，商群是N的陪集组成的群，是aN，bN组成的群，aN和bN都是不相交的

命题1.3.3子群陪集要不然相等要不然无交

定义1.28商集与指数 G/N就是N所有的陪集，指数就是陪集的数量

\end{proposition}

\begin{note}[关联图片]
以下图片已纳入本条整理，当前版本优先依据 Markdown 文字整理，图片细节可在审核时对照：

1. image.png\\

\end{note}

\begin{note}[来源路径]

近世代数/近世代数/无标题/命题1 37 商群 1d8d2efa3f4280a2ab8bc71b5decdd72.md

\end{note}

\section*{命题1.38正规子群的交集还是正规子群}

\begin{note}[来源上级主题]

1.4正规子群 (1\%204\%E6\%AD\%A3\%E8\%A7\%84\%E5\%AD\%90\%E7\%BE\%A4\%201ead2efa3f428095a412e65ae0862a5d.md)

\end{note}

\begin{proposition}[命题1.38正规子群的交集还是正规子群]

Aria 1.38

\& (Nidier HH G MH IEM FFF, MEAN REG RE G AEA a, BP

(Ni dG                                             (1.75)

tel                                                          6

WE 4, RN TETRA ERRATA. RPA STRAW ARKRTRHEWAZILERA-RN, HK

BMRA, PHR=EAR HH WAR T.

ALARM REE EAE. RFA BHR. Fae Gne Nice Ni, RN RAMEE ana! € je Nie

fERiel, WineN;. HEN AG, RMA ana'eN;. Ash, ana e€\textbackslash{}jie Nico KRIEHT Nie Ni AG.

WBA ARH, BEAL AT DAUEBA, 3-7 SE ABET DAZE MIELE, RE SCA TA EL FEY TE

PULP RENI ACER , TMG SACRA EP IERL SE. FEL, FR Barat, RRR eae AT DAE SS TTB

FR TTD 4 AS — SE TELE BI. TSE, SF EERE (G,-), EPL SHE \{e\} AUR -SHE G ARIES HE, BN

\end{proposition}

\begin{note}[图片 OCR 摘录，待人工复核]
以下内容来自源图片识别，便于审核时对照：

\textbf{图片 1}（image.png）
\begin{Verbatim}[fontsize=\small]
Aria 1.38
& (Nidier HH G MH IEM FFF, MEAN REG RE G AEA a, BP
(Ni dG                                             (1.75)
tel                                                          6
WE 4, RN TETRA ERRATA. RPA STRAW ARKRTRHEWAZILERA-RN, HK
BMRA, PHR=EAR HH WAR T.
ALARM REE EAE. RFA BHR. Fae Gne Nice Ni, RN RAMEE ana! € je Nie
fERiel, WineN;. HEN AG, RMA ana'eN;. Ash, ana e€\jie Nico KRIEHT Nie Ni AG.
WBA ARH, BEAL AT DAUEBA, 3-7 SE ABET DAZE MIELE, RE SCA TA EL FEY TE
PULP RENI ACER , TMG SACRA EP IERL SE. FEL, FR Barat, RRR eae AT DAE SS TTB
FR TTD 4 AS — SE TELE BI. TSE, SF EERE (G,-), EPL SHE {e} AUR -SHE G ARIES HE, BN
\end{Verbatim}

\end{note}

\begin{note}[来源路径]

近世代数/近世代数/无标题/命题1 38正规子群的交集还是正规子群 1d8d2efa3f428001ba66de6c2e2fde67.md

\end{note}

\section*{命题1.39 1.40e和自己都是自己的正规子群交换群的子群都是正规子群}

\begin{note}[来源上级主题]

1.4正规子群 (1\%204\%E6\%AD\%A3\%E8\%A7\%84\%E5\%AD\%90\%E7\%BE\%A4\%201ead2efa3f428095a412e65ae0862a5d.md)

\end{note}

\begin{proposition}[命题1.39 1.40e和自己都是自己的正规子群交换群的子群都是正规子群]

待根据图片进一步人工校对。

\end{proposition}

\begin{note}[来源路径]

近世代数/近世代数/无标题/命题1 39 1 40e和自己都是自己的正规子群交换群的子群都是正规子群 1d8d2efa3f4280beb4f3c1126c211d2b.md

\end{note}

\section*{命题1.5证明A生成的子幺半群是包含A的最小子幺半群}

\begin{note}[来源上级主题]

1.1幺半群 (1\%201\%E5\%B9\%BA\%E5\%8D\%8A\%E7\%BE\%A4\%201cad2efa3f4280628789c382df95fdc5.md)

\end{note}

\begin{proof}

由一个子集A生成的子结构就是包含了A的所有子结构的交集

所有幺半群都有e自然在交集了吗，证明x，y有x*y就有因为每一个群x，y在交集里面所以每个半幺群里面都有，所以每个半幺群里面也有x*y所以x*y也在交集里面

\end{proof}

\begin{note}[关联图片]
以下图片已纳入本条整理，当前版本优先依据 Markdown 文字整理，图片细节可在审核时对照：

1. IMG\_2906.jpeg\\

\end{note}

\begin{note}[来源路径]

近世代数/近世代数/无标题/命题1 5证明A生成的子幺半群是包含A的最小子幺半群 1cad2efa3f4280d7aeccdd9429149cb4.md

\end{note}

\section*{命题1.7证明方法把单位元拆解然后用结合律引出新的东西唯一性证明}

\begin{note}[来源上级主题]

1.1幺半群 (1\%201\%E5\%B9\%BA\%E5\%8D\%8A\%E7\%BE\%A4\%201cad2efa3f4280628789c382df95fdc5.md)

\end{note}

\begin{proof}

fiz jel: 1.7

A (S,:) EP AER. PGR x eS HATA, WHE. LHL, Ry, y eSMREMUA, 2

y=y's                                                                                                        -

UW Rik yw Mex He. My-x=e, x-y =e, RNERRAA-POWHER ‘RREW” RIE

y=y'o

y=y-e=y-x-y =e-y'=y’                     (1.21)

HPRNAAY ce z¥tirt, WR (OX) SOE (AARNBRA MEST).

FRAT KI ZA APES ES TA CRAB CN) (ADRAC) , SPOPENS RRS, ATONE

EK RY. TMI AAS, RHE Re, GEAR ME (AN eee, Ee a CPE

JERE HO ZRIA ANS ad). SSE, PEPER Ae, FRE A ABE De, TA RE OE

FEA ES PRA, ASAT SEAR AY. RAE

\end{proof}

\begin{note}[图片 OCR 摘录，待人工复核]
以下内容来自源图片识别，便于审核时对照：

\textbf{图片 1}（IMG\_2908.jpeg）
\begin{Verbatim}[fontsize=\small]
fiz jel: 1.7

A (S,:) EP AER. PGR x eS HATA, WHE. LHL, Ry, y eSMREMUA, 2

y=y's                                                                                                        -
UW Rik yw Mex He. My-x=e, x-y =e, RNERRAA-POWHER ‘RREW” RIE
y=y'o

y=y-e=y-x-y =e-y'=y’                     (1.21)
HPRNAAY ce z¥tirt, WR (OX) SOE (AARNBRA MEST).

FRAT KI ZA APES ES TA CRAB CN) (ADRAC) , SPOPENS RRS, ATONE
EK RY. TMI AAS, RHE Re, GEAR ME (AN eee, Ee a CPE
JERE HO ZRIA ANS ad). SSE, PEPER Ae, FRE A ABE De, TA RE OE
FEA ES PRA, ASAT SEAR AY. RAE
\end{Verbatim}

\end{note}

\begin{note}[来源路径]

近世代数/近世代数/无标题/命题1 7证明方法把单位元拆解然后用结合律引出新的东西唯一性证明 1cbd2efa3f428061b991eda66a78e31d.md

\end{note}

\section*{商群性质}

\begin{note}[商群性质]

Rata ROI ARTE:

DH G/N hh= [G: N].

2) |G RARE, URE

lel

G/N ‘t= [G: N] = IN]

3) ARR AUR, ESC ROBY OT

)NQG.H e=eN=N 8B GIN

wees, a Now aN wrx.

\end{note}

\begin{note}[图片 OCR 摘录，待人工复核]
以下内容来自源图片识别，便于审核时对照：

\textbf{图片 1}（image.png）
\begin{Verbatim}[fontsize=\small]
Rata ROI ARTE:
DH G/N hh= [G: N].
2) |G RARE, URE
lel
G/N ‘t= [G: N] = IN]
3) ARR AUR, ESC ROBY OT
)NQG.H e=eN=N 8B GIN
wees, a Now aN wrx.
\end{Verbatim}

\end{note}

\begin{note}[来源路径]

近世代数/近世代数/无标题/商群性质 1cad2efa3f42800d9fd2d13889aa2844.md

\end{note}

\section*{子群}

\begin{note}[子群]

子级 项目: 定理1子群充要条件 (\%E5\%AE\%9A\%E7\%90\%861\%E5\%AD\%90\%E7\%BE\%A4\%E5\%85\%85\%E8\%A6\%81\%E6\%9D\%A1\%E4\%BB\%B6\%201bdd2efa3f428096a849dfb1aa61075c.md), 定理2非空子集充要条件 (\%E5\%AE\%9A\%E7\%90\%862\%E9\%9D\%9E\%E7\%A9\%BA\%E5\%AD\%90\%E9\%9B\%86\%E5\%85\%85\%E8\%A6\%81\%E6\%9D\%A1\%E4\%BB\%B6\%201bdd2efa3f4280eca5fdd0e948316100.md), 定理3子群充要条件 非空有限子群 (\%E5\%AE\%9A\%E7\%90\%863\%E5\%AD\%90\%E7\%BE\%A4\%E5\%85\%85\%E8\%A6\%81\%E6\%9D\%A1\%E4\%BB\%B6\%20\%E9\%9D\%9E\%E7\%A9\%BA\%E6\%9C\%89\%E9\%99\%90\%E5\%AD\%90\%E7\%BE\%A4\%201bdd2efa3f4280bfbf36f4060b3612ad.md), 循环子群 (\%E5\%BE\%AA\%E7\%8E\%AF\%E5\%AD\%90\%E7\%BE\%A4\%201bdd2efa3f428088afeae785ed4258af.md), 学子群的意义 (\%E5\%AD\%A6\%E5\%AD\%90\%E7\%BE\%A4\%E7\%9A\%84\%E6\%84\%8F\%E4\%B9\%89\%201c6d2efa3f428082a2a0c319d06b2a4e.md)

继承运算

子群也得是群而且运算和父空间一样

\end{note}

\begin{note}[关联图片]
以下图片已纳入本条整理，当前版本优先依据 Markdown 文字整理，图片细节可在审核时对照：

1. image.png\\
2. image 1.png\\
3. image 2.png\\

\end{note}

\begin{note}[来源路径]

近世代数/近世代数/无标题/子群 1bdd2efa3f42801b985de1cf59e5ed5f.md

\end{note}

\section*{子群的阶定义拉格朗日定理}

\begin{note}[来源上级主题]

1.3有限群 (1\%203\%E6\%9C\%89\%E9\%99\%90\%E7\%BE\%A4\%201ead2efa3f4280c39a49d8b0af9aec03.md)

\end{note}

\begin{theorem}[子群的阶定义拉格朗日定理]

EM 1.26                 =

SHRG ATH, MAA, ie |Al, CLAHMHREK). GSH AMRF IZ |A| = 0,    .

ASH BER SE TY IPE «

fir WH 1.32

GHAG HTH, RHYME GC AU, BP

wilicl                  (1.54)

a

ROC AA REY RE ABS ERY), EL EY BO A EY A. CS LT DAR Fi

REP A EFL (HT DA ee BN A EPEC). BEEN Sie, Be ee |

\end{theorem}

\begin{note}[图片 OCR 摘录，待人工复核]
以下内容来自源图片识别，便于审核时对照：

\textbf{图片 1}（image.png）
\begin{Verbatim}[fontsize=\small]
EM 1.26                 =
SHRG ATH, MAA, ie |Al, CLAHMHREK). GSH AMRF IZ |A| = 0,    .
ASH BER SE TY IPE «
fir WH 1.32
GHAG HTH, RHYME GC AU, BP
wilicl                  (1.54)
a
ROC AA REY RE ABS ERY), EL EY BO A EY A. CS LT DAR Fi
REP A EFL (HT DA ee BN A EPEC). BEEN Sie, Be ee |
\end{Verbatim}

\end{note}

\begin{note}[来源路径]

近世代数/近世代数/无标题/子群的阶定义拉格朗日定理 1d2d2efa3f4280abb951d3f9b92ef54b.md

\end{note}

\section*{子群的陪集}

\begin{note}[子群的陪集]

子级 项目: 陪集的意义 (\%E9\%99\%AA\%E9\%9B\%86\%E7\%9A\%84\%E6\%84\%8F\%E4\%B9\%89\%201c6d2efa3f42801ba4c4d1f60871c4bc.md), 右陪集的定义 (\%E5\%8F\%B3\%E9\%99\%AA\%E9\%9B\%86\%E7\%9A\%84\%E5\%AE\%9A\%E4\%B9\%89\%201c3d2efa3f428044a7f6c58dda5da09f.md), 例子 (\%E4\%BE\%8B\%E5\%AD\%90\%201c4d2efa3f428049872dd61c51dda43e.md), 陪集的定义 (\%E9\%99\%AA\%E9\%9B\%86\%E7\%9A\%84\%E5\%AE\%9A\%E4\%B9\%89\%201c4d2efa3f42804e9574d5a966680083.md), 左陪集 (\%E5\%B7\%A6\%E9\%99\%AA\%E9\%9B\%86\%201c3d2efa3f4280bebcadc5fc4efa6e1e.md), 定理1 左右陪集个数相等，有限还相等或都是无限 (\%E5\%AE\%9A\%E7\%90\%861\%20\%E5\%B7\%A6\%E5\%8F\%B3\%E9\%99\%AA\%E9\%9B\%86\%E4\%B8\%AA\%E6\%95\%B0\%E7\%9B\%B8\%E7\%AD\%89\%EF\%BC\%8C\%E6\%9C\%89\%E9\%99\%90\%E8\%BF\%98\%E7\%9B\%B8\%E7\%AD\%89\%E6\%88\%96\%E9\%83\%BD\%E6\%98\%AF\%E6\%97\%A0\%E9\%99\%90\%201c3d2efa3f4280489fb7c3b3a78e65b1.md), 定义． 陪集个数为指数 (\%E5\%AE\%9A\%E4\%B9\%89\%EF\%BC\%8E\%20\%E9\%99\%AA\%E9\%9B\%86\%E4\%B8\%AA\%E6\%95\%B0\%E4\%B8\%BA\%E6\%8C\%87\%E6\%95\%B0\%201c3d2efa3f4280d4ae10d01394813e30.md), 定理2 (\%E5\%AE\%9A\%E7\%90\%862\%201c3d2efa3f42808d86d9d910407f5566.md), 定理3有限群的阶都是父群因子的阶 (\%E5\%AE\%9A\%E7\%90\%863\%E6\%9C\%89\%E9\%99\%90\%E7\%BE\%A4\%E7\%9A\%84\%E9\%98\%B6\%E9\%83\%BD\%E6\%98\%AF\%E7\%88\%B6\%E7\%BE\%A4\%E5\%9B\%A0\%E5\%AD\%90\%E7\%9A\%84\%E9\%98\%B6\%201c3d2efa3f4280c2992cc6ef4f5024c8.md), 奇偶排列素数阶的群一定是循环群 (\%E5\%A5\%87\%E5\%81\%B6\%E6\%8E\%92\%E5\%88\%97\%E7\%B4\%A0\%E6\%95\%B0\%E9\%98\%B6\%E7\%9A\%84\%E7\%BE\%A4\%E4\%B8\%80\%E5\%AE\%9A\%E6\%98\%AF\%E5\%BE\%AA\%E7\%8E\%AF\%E7\%BE\%A4\%201bdd2efa3f428081b6cdc32b80a5f208.md)

\end{note}

\begin{note}[来源路径]

近世代数/近世代数/无标题/子群的陪集 1c3d2efa3f42806a85b0f0113028b207.md

\end{note}

\section*{子群里面的元素当陪集的x，则这个配集还是原来子群xH=H}

\begin{note}[来源上级主题]

1.3有限群 (1\%203\%E6\%9C\%89\%E9\%99\%90\%E7\%BE\%A4\%201ead2efa3f4280c39a49d8b0af9aec03.md)

\end{note}

\begin{note}[子群里面的元素当陪集的x，则这个配集还是原来子群xH=H]

这里就把前面的g可以分为H里面的和G-H里面的，核心原因是子群对乘法封闭

\end{note}

\begin{note}[关联图片]
以下图片已纳入本条整理，当前版本优先依据 Markdown 文字整理，图片细节可在审核时对照：

1. IMG\_2987.jpeg\\

\end{note}

\begin{note}[来源路径]

近世代数/近世代数/无标题/子群里面的元素当陪集的x，则这个配集还是原来子群xH=H 1d2d2efa3f4280d29dd3d8346705ba89.md

\end{note}

\section*{学子群的意义}

\begin{note}[来源上级主题]

子群 (\%E5\%AD\%90\%E7\%BE\%A4\%201bdd2efa3f42801b985de1cf59e5ed5f.md)

\end{note}

\begin{note}[学子群的意义]

1, SEAREE ROHAN

SREB “POE”. BATRA ATE, BATEIAT RPA eA

Qa, ANDRO AE, ON, AF RT MERE OR, B

AU HRA, AMUBTRMAEN MERRIER.

2. fa HC EBRAD AE

ARRRLOGM, BTIUAE-TFRARASIC, REBORN. Bl

i, TPR OREAN, ASSAM MAF MLE MP REE EEN

\#215 SORT

3. BAO ASS BE

FH, CHALMTH, QOBRASMSAHEM, BUMATH, BATTLE LE

ROAR, TARAS RHA, IATARERTAUR (HR), BMITEMIIA, BCH

RDPB SLR CARES ABE EE

4. RES BLA

EXIM, FRA CRA BEB” LBD” WOMEN. HIM, ZEA. ME

SAUCER, —MANEAMMEEE TORRE TF MME, BT

BY, BATRA BigP PRS A AME RN

5. MI OIF ARO

FRARUEAT WESRBAMNLA, REN, BUI TRORMETFER

CAMA, BATTEN FR RG, MRNBCS MASTER.

82, PREBTLEERRESANWK, LRBUBES, SECM SHMA

(OUEST, DRELS) ML. BRAS IFH, RRNA

BRA, SPANAIR.

\end{note}

\begin{note}[图片 OCR 摘录，待人工复核]
以下内容来自源图片识别，便于审核时对照：

\textbf{图片 1}（IMG\_2826.jpeg）
\begin{Verbatim}[fontsize=\small]
1, SEAREE ROHAN

SREB “POE”. BATRA ATE, BATEIAT RPA eA
Qa, ANDRO AE, ON, AF RT MERE OR, B
AU HRA, AMUBTRMAEN MERRIER.

2. fa HC EBRAD AE

ARRRLOGM, BTIUAE-TFRARASIC, REBORN. Bl
i, TPR OREAN, ASSAM MAF MLE MP REE EEN
#215 SORT

3. BAO ASS BE

FH, CHALMTH, QOBRASMSAHEM, BUMATH, BATTLE LE
ROAR, TARAS RHA, IATARERTAUR (HR), BMITEMIIA, BCH
RDPB SLR CARES ABE EE

4. RES BLA

EXIM, FRA CRA BEB” LBD” WOMEN. HIM, ZEA. ME
SAUCER, —MANEAMMEEE TORRE TF MME, BT
BY, BATRA BigP PRS A AME RN

5. MI OIF ARO

FRARUEAT WESRBAMNLA, REN, BUI TRORMETFER
CAMA, BATTEN FR RG, MRNBCS MASTER.
82, PREBTLEERRESANWK, LRBUBES, SECM SHMA
(OUEST, DRELS) ML. BRAS IFH, RRNA
BRA, SPANAIR.
\end{Verbatim}

\end{note}

\begin{note}[来源路径]

近世代数/近世代数/无标题/学子群的意义 1c6d2efa3f428082a2a0c319d06b2a4e.md

\end{note}

\section*{定义1.11n阶一般线性群}

\begin{note}[来源上级主题]

1.2群 (1\%202\%E7\%BE\%A4\%201ecd2efa3f4280cf8df9c881d186f2a4.md)

\end{note}

\begin{definition}[定义1.11n阶一般线性群]

BAT F ARLE nen PUK LER RAY RAAF, RA (KH LhAy) n SAAB, IVE (GL(n,R),-). W

FPR TES AGH TARAR, Ast

GL(n,R) = \{A € M(n,R) : det(A) \# 0\}                               (1.26)  .

FER, DOWN WAZE- MEF, AAAS 2h n « n SEE PAEXTHETE PM AR RE (PEL A AE) , BE POR FETT FY

Wc, WEE THRACE. 49K, FTA RBCEN, ABPACEY , SHE AR PERE

fe Bl) FAA PETES EAS BGA IRA VERE

\end{definition}

\begin{note}[图片 OCR 摘录，待人工复核]
以下内容来自源图片识别，便于审核时对照：

\textbf{图片 1}（IMG\_2913.jpeg）
\begin{Verbatim}[fontsize=\small]
BAT F ARLE nen PUK LER RAY RAAF, RA (KH LhAy) n SAAB, IVE (GL(n,R),-). W
FPR TES AGH TARAR, Ast
GL(n,R) = {A € M(n,R) : det(A) # 0}                               (1.26)  .
FER, DOWN WAZE- MEF, AAAS 2h n « n SEE PAEXTHETE PM AR RE (PEL A AE) , BE POR FETT FY
Wc, WEE THRACE. 49K, FTA RBCEN, ABPACEY , SHE AR PERE
fe Bl) FAA PETES EAS BGA IRA VERE
\end{Verbatim}

\textbf{图片 2}（image.png）
\begin{Verbatim}[fontsize=\small]
FATE — PY, PRUE AACE GL (n, R) LEAP IRA HERE SL(n, R) PRET A Za PN? DLE AG
yy, AA
det : GL(n,R) > R*                                             (1.89)
Fe MRIS, HL ker(det) = SL(n,R), SWE — ERE, BTA
GL(n,R)/SL(n, R) ~ RX                                                  (1.90)
\end{Verbatim}

\end{note}

\begin{note}[来源路径]

近世代数/近世代数/无标题/定义1 11n阶一般线性群 1cbd2efa3f4280dda90ce3802a0d1200.md

\end{note}

\section*{定义1.12n阶特殊线性群行列式是1}

\begin{note}[来源上级主题]

1.2群 (1\%202\%E7\%BE\%A4\%201ecd2efa3f4280cf8df9c881d186f2a4.md)

\end{note}

\begin{definition}[定义1.12n阶特殊线性群行列式是1]

we SL 1.12

(KH LAG) WAFER ERE Wy ARLE ATF ASAE 1g nen RHEE MAY RABE, IDE (SL(n,R),:), PP

SL(n,R) = \{A € M(n,R) : det(A) = 1\}             (127) |

TERE, KMESGEM AAR REND, AEA ILHAM ee TE. SEP, Ee SP. BEATE

AAS FREI ARE? TERN FRAT 2A FZ PRE AR EAB EE

\end{definition}

\begin{note}[图片 OCR 摘录，待人工复核]
以下内容来自源图片识别，便于审核时对照：

\textbf{图片 1}（IMG\_2914.jpeg）
\begin{Verbatim}[fontsize=\small]
we SL 1.12
(KH LAG) WAFER ERE Wy ARLE ATF ASAE 1g nen RHEE MAY RABE, IDE (SL(n,R),:), PP
SL(n,R) = {A € M(n,R) : det(A) = 1}             (127) |
TERE, KMESGEM AAR REND, AEA ILHAM ee TE. SEP, Ee SP. BEATE
AAS FREI ARE? TERN FRAT 2A FZ PRE AR EAB EE
\end{Verbatim}

\end{note}

\begin{note}[来源路径]

近世代数/近世代数/无标题/定义1 12n阶特殊线性群行列式是1 1cbd2efa3f42805e8f6be357d2dce70f.md

\end{note}

\section*{定义1.13定义子群根本目的是证明它是一个群}

\begin{note}[来源上级主题]

1.2群 (1\%202\%E7\%BE\%A4\%201ecd2efa3f4280cf8df9c881d186f2a4.md)

\end{note}

\begin{proof}

子级 项目: 命题1.11证明 (\%E5\%91\%BD\%E9\%A2\%981\%2011\%E8\%AF\%81\%E6\%98\%8E\%201cbd2efa3f4280679acbe2813bc3b4e4.md), 命题1.12特殊线性群是群 (\%E5\%91\%BD\%E9\%A2\%981\%2012\%E7\%89\%B9\%E6\%AE\%8A\%E7\%BA\%BF\%E6\%80\%A7\%E7\%BE\%A4\%E6\%98\%AF\%E7\%BE\%A4\%201cbd2efa3f4280bc9d1cfc2496d61576.md)

\end{proof}

\begin{note}[关联图片]
以下图片已纳入本条整理，当前版本优先依据 Markdown 文字整理，图片细节可在审核时对照：

1. IMG\_2915.jpeg\\

\end{note}

\begin{note}[来源路径]

近世代数/近世代数/无标题/定义1 13定义子群根本目的是证明它是一个群 1cbd2efa3f4280ee903ac1e652f44dec.md

\end{note}

\section*{定义1.15群同态下定义 核是映射到e‘的原像集 imf是像集}

\begin{note}[来源上级主题]

1.2群 (1\%202\%E7\%BE\%A4\%201ecd2efa3f4280cf8df9c881d186f2a4.md)

\end{note}

\begin{definition}[定义1.15群同态下定义 核是映射到e‘的原像集 imf是像集]

核是映射到对方单位元对应的原像集，像是G能映射到G‘的所有G’的元素。但是因为同态不要求单满所以imf（f）至少G的子集

要证两件事1有单位元2.x，y在这里面那么xy-1次方也在，用kerf的定义也就是验证这个东西能不能映射到e‘上。e’本身就是一个群，反正要把运算通过同态的性质转化到已知的群上。

\end{definition}

\begin{note}[关联图片]
以下图片已纳入本条整理，当前版本优先依据 Markdown 文字整理，图片细节可在审核时对照：

1. IMG\_2933.jpeg\\
2. IMG\_2934.jpeg\\
3. IMG\_2935.jpeg\\

\end{note}

\begin{note}[来源路径]

近世代数/近世代数/无标题/定义1 15群同态下定义 核是映射到e‘的原像集 imf是像集 1d1d2efa3f42808daa56cdc3a30925f6.md

\end{note}

\section*{定义1.16群同态}

\begin{note}[来源上级主题]

1.2群 (1\%202\%E7\%BE\%A4\%201ecd2efa3f4280cf8df9c881d186f2a4.md)

\end{note}

\begin{definition}[定义1.16群同态]

左乘，同态最好的就是可以拆括号让x和x逆相遇

\end{definition}

\begin{note}[图片 OCR 摘录，待人工复核]
以下内容来自源图片识别，便于审核时对照：

\textbf{图片 1}（IMG\_2938.jpeg）
\begin{Verbatim}[fontsize=\small]
we SL 1.16              —

Sf (G) 4 (Gs) RAMA, BUNA SRA

.             A, Milt f R-DMAAY f Rho, thf R-PHMAS f RH
a
\end{Verbatim}

\textbf{图片 2}（IMG\_2936.jpeg）
\begin{Verbatim}[fontsize=\small]
& fi: (G, ‘) > (G’, *) 7a —/ FFA AS,  ny) f VBA SAWS ker(f) = {e}. ALU,  —/ # a A
Bay HAA HK RF LAY.                                 .
VED) BR f 2h, MARA fle)=e'’, Alte f(x) =e’, WHABHNERRIN-CAx=e, RREWT
\end{Verbatim}

\textbf{图片 3}（IMG\_2937.jpeg）
\begin{Verbatim}[fontsize=\small]
KEP. (RAFAZERM)

BDFARMA BR. RAVE ker(f) = {e'}. ik xx’ eG, HA f(x) = f(x’), RA RAGE x =x".
ERE, ANAM ER SO’), BS COLO) = fax) =e. WAAKE FILM, PED QAA xx! =e.
EPRAHGR x’, RMR xox’. RRL RAM.

f                             ;            f
-                    G’                                        G’               ~                         G’
Pe] 1.2: ESLRRIIAS, WRAL AS.
VERE, ZE3C/MIEBAIUPE BE, FRITH RT SAMIR CE, HEBIREPER (ETE) FECL, BE
TABULA (AGE, REPGRAMARIN). AK, RNASE UINR MAC, HARA WR cde
HL. URE, BCS BLL ROR TRA HE (UDETE). ARRAS, JCM UTE RE Ha eh nit
HET REIS, BRAY SRS. BONE RMI ZOE SD
\end{Verbatim}

\end{note}

\begin{note}[来源路径]

近世代数/近世代数/无标题/定义1 16群同态 1d1d2efa3f4280dfbb71fb6ae82fca39.md

\end{note}

\section*{定义1.18两个群直积还是群}

\begin{note}[来源上级主题]

1.2群 (1\%202\%E7\%BE\%A4\%201ecd2efa3f4280cf8df9c881d186f2a4.md)

\end{note}

\begin{definition}[定义1.18两个群直积还是群]

EL 1.18                                —

4 (G1), (G’,-2) RANA, AME C11 BR, Kit (GxG’,*). WF wy) A yWIEEXKG’, K

NV LRA ADR AR, A

(xy) *@,y) =x, y 2’)                      (1.39) .

fir el 1.18

iW) HAMM: BAG \#1 FHA, CO £2. THA, WExG HRFRPZBLREXM, WExe st

* = (41,-2) Rta 3t Hl BY.

BLT: RRMA, (c,e’) EHRW BET. HER (x,y) € GxG’, BINA (wy) \# (0, e’) = (x-1e, 9-2e’) =

(x,y), A-Wthe AH, RHE T (ee) LHRH BLT.

\end{definition}

\begin{note}[图片 OCR 摘录，待人工复核]
以下内容来自源图片识别，便于审核时对照：

\textbf{图片 1}（IMG\_2940.jpeg）
\begin{Verbatim}[fontsize=\small]
EL 1.18                                —

4 (G1), (G’,-2) RANA, AME C11 BR, Kit (GxG’,*). WF wy) A yWIEEXKG’, K

NV LRA ADR AR, A

(xy) *@,y) =x, y 2’)                      (1.39) .

fir el 1.18
iW) HAMM: BAG #1 FHA, CO £2. THA, WExG HRFRPZBLREXM, WExe st
* = (41,-2) Rta 3t Hl BY.

BLT: RRMA, (c,e’) EHRW BET. HER (x,y) € GxG’, BINA (wy) # (0, e’) = (x-1e, 9-2e’) =
(x,y), A-Wthe AH, RHE T (ee) LHRH BLT.
\end{Verbatim}

\textbf{图片 2}（IMG\_2941.jpeg）
\begin{Verbatim}[fontsize=\small]
7G: MAA ERA, UTHER ye GxG’, wlLy ley) He. EHRAUM, RESM.
BEA SIME ELAR, Bet AT VASSAL ae SC R-II. ERE a,
SOTA VAT ERE, ECE. ARITA SI ty)
\end{Verbatim}

\end{note}

\begin{note}[来源路径]

近世代数/近世代数/无标题/定义1 18两个群直积还是群 1d1d2efa3f4280189fb9c52e96214d0c.md

\end{note}

\section*{定义1.19一族群的值积 定义1.20还需再看}

\begin{note}[来源上级主题]

1.3有限群 (1\%203\%E6\%9C\%89\%E9\%99\%90\%E7\%BE\%A4\%201ead2efa3f4280c39a49d8b0af9aec03.md)

\end{note}

\begin{definition}[定义1.19一族群的值积 定义1.20还需再看]

xe SL 1.19

(Gi, -iicr HARA, HP 1 eK. MTA LEN MARA (Te Gi *), APBTRAMRER.

MF (xdier, Wdier € je Gi, RAVEN

(ai ier * Wdier = OG i Yidier                                     (1.40)  *

JEWS, FAN

fir jel 1.19

\& (Gi, ier HARA, MCA AAR ([je, Gi, *) LEV FF.                                              “

UE HEA SRSA, KAM ROB -KER. HMESHSOEEEAWN. BE (edie, M W)ie H\#

Toe (x; "ier

IZ LY Bil-F A a nn SSO (R, +) FEBIAY CR", +). FETE A EY AS EM EY IY

TORE IE EIA. AK, TERA TEU TOL BF, FTA ERIE CHINE) WRG,

FRAT (BME) «

FRIAR STR RAR FRA, EAR EI PA. ERMAN EE, SE PEASE TAB FRAP

PFE

TEX IE, (Gi, tier HE RAE, fe 1 eR, FETE BUN hs 7 RESON

pj: | [Gi > G;                           (1.41)

tel

WH Giier, FETE CAA HY

pj(Xi)ier) =X;                                                (1.42)

AY cK, FAUNA BAER RR Amie. 49K, SETA RRR CAR) 28

WOR Aie AFAR ROR FRR (AR) EL. AB, Fe Sei RA SSS: MCRAE EER A

fib EAS. AK, FUN A WMA.

\end{definition}

\begin{note}[图片 OCR 摘录，待人工复核]
以下内容来自源图片识别，便于审核时对照：

\textbf{图片 1}（IMG\_2942.jpeg）
\begin{Verbatim}[fontsize=\small]
xe SL 1.19
(Gi, -iicr HARA, HP 1 eK. MTA LEN MARA (Te Gi *), APBTRAMRER.
MF (xdier, Wdier € je Gi, RAVEN
(ai ier * Wdier = OG i Yidier                                     (1.40)  *
JEWS, FAN
fir jel 1.19
& (Gi, ier HARA, MCA AAR ([je, Gi, *) LEV FF.                                              “
UE HEA SRSA, KAM ROB -KER. HMESHSOEEEAWN. BE (edie, M W)ie H#
Toe (x; "ier
IZ LY Bil-F A a nn SSO (R, +) FEBIAY CR", +). FETE A EY AS EM EY IY
TORE IE EIA. AK, TERA TEU TOL BF, FTA ERIE CHINE) WRG,
FRAT (BME) «
FRIAR STR RAR FRA, EAR EI PA. ERMAN EE, SE PEASE TAB FRAP
PFE
TEX IE, (Gi, tier HE RAE, fe 1 eR, FETE BUN hs 7 RESON
pj: | [Gi > G;                           (1.41)
tel
WH Giier, FETE CAA HY
pj(Xi)ier) =X;                                                (1.42)
AY cK, FAUNA BAER RR Amie. 49K, SETA RRR CAR) 28
WOR Aie AFAR ROR FRR (AR) EL. AB, Fe Sei RA SSS: MCRAE EER A
fib EAS. AK, FUN A WMA.
\end{Verbatim}

\textbf{图片 2}（IMG\_2943.jpeg）
\begin{Verbatim}[fontsize=\small]
fri 1.20                                                  “
% (Gi, tier RHA, GET RIERA, Wa pj: [],<,Gi > Gj RP AAA.                   ‘
WEA eA ay AY, (AKA AER TSC EE. NGL APRRRE Lvl, CHEMIE TA BEA Ei] —- Mi
ae “ELBA, GEMMA HA, BES ASN PA eT TF .” (APLAR RI AE GE , EAN AS a AEB A RE I AS
WE, APBD A BA EE, ABLE BR TT A Se oy
UE) + (ier, Wier € Tier Gi, Ml
Pi((XiNier) =Xj,  Pi(Nidier) =Y;                (1.43)
Pj(Xaier * (Vidier) = Pj( ii Ydier) =X; +7 V7 = PJ (Xier) +7 Pj (Ndier)                (1.44)
\end{Verbatim}

\end{note}

\begin{note}[来源路径]

近世代数/近世代数/无标题/定义1 19一族群的值积 定义1 20还需再看 1d1d2efa3f4280bbb5e6c6260092ce70.md

\end{note}

\section*{定义1.2.4子集生成群证明}

\begin{note}[来源上级主题]

1.3有限群 (1\%203\%E6\%9C\%89\%E9\%99\%90\%E7\%BE\%A4\%201ead2efa3f4280c39a49d8b0af9aec03.md)

\end{note}

\begin{proof}

定义1.6由子集生成的自幺半群

\end{proof}

\begin{note}[图片 OCR 摘录，待人工复核]
以下内容来自源图片识别，便于审核时对照：

\textbf{图片 1}（IMG\_2948.jpeg）
\begin{Verbatim}[fontsize=\small]
ESL 1.24
4 (G,) RP #H, SCG. Nh SARA, IDE (S), CLA
(S)=( {Hc G:H>S,H<G}             (1.47)
&
MOOR, Fe Ser EAE) (S) WASE ETM.
fir 1.24
& (G,+) e—-/FF, ALSCG, WM (S) <G.                                             “
ULV) AE, RN RETR KASCH HL, ERAPHCPHA.
Ri EX, (SS) LHMALST SHG HTHRSRMR AREA.
Bi: ESAS HARSH, KEMNNWKREMES BT.
FIKH APE: Bik x,y e(S), ER-+PESTY SHFRA, WxyeH. HAH RTH, KryedH, Pry
xy €(S).
HAMM: 5 EMR -AAHAW, BIER (S) PRAM TRES—PREW A, EIR
(S)
AIK, GHP S AE HE, SEA SS Waa) PH
AAA Ee Hc 8 AE EY, DORE A PEEL BCAA
\end{Verbatim}

\end{note}

\begin{note}[来源路径]

近世代数/近世代数/无标题/定义1 2 4子集生成群证明 1d1d2efa3f42805fa62fd5a0548e8041.md

\end{note}

\section*{定义1.20元素的阶}

\begin{note}[来源上级主题]

群的基础知识 (\%E7\%BE\%A4\%E7\%9A\%84\%E5\%9F\%BA\%E7\%A1\%80\%E7\%9F\%A5\%E8\%AF\%86\%201aed2efa3f428009ac4eca74062880ff.md)

\end{note}

\begin{definition}[定义1.20元素的阶]

子级 项目: a和a的逆阶数一样 (a\%E5\%92\%8Ca\%E7\%9A\%84\%E9\%80\%86\%E9\%98\%B6\%E6\%95\%B0\%E4\%B8\%80\%E6\%A0\%B7\%201b1d2efa3f4280c9848ef6cf1a5a43d6.md), 有限群每一个元素都有元 (\%E6\%9C\%89\%E9\%99\%90\%E7\%BE\%A4\%E6\%AF\%8F\%E4\%B8\%80\%E4\%B8\%AA\%E5\%85\%83\%E7\%B4\%A0\%E9\%83\%BD\%E6\%9C\%89\%E5\%85\%83\%201b1d2efa3f428077b99af9bef5e16a15.md), 阶数与奇偶 (\%E9\%98\%B6\%E6\%95\%B0\%E4\%B8\%8E\%E5\%A5\%87\%E5\%81\%B6\%201b1d2efa3f4280fe9903f391009c0bee.md), 阶的例题 (\%E9\%98\%B6\%E7\%9A\%84\%E4\%BE\%8B\%E9\%A2\%98\%201b1d2efa3f42801ca927ea10853179f8.md)

乘法不是真正的乘法

a的n次方=aoaoa 就是n次运算，运算可以定义

阶是最小正整数 阶的符号和绝对值一样

有限群每个元素的阶均有限 a和阶和a逆的阶一样

定义1.22 x的幂

交换群可以拆括号

交换群 阿贝尔群

\end{definition}

\begin{note}[关联图片]
以下图片已纳入本条整理，当前版本优先依据 Markdown 文字整理，图片细节可在审核时对照：

1. IMG\_2944.jpeg\\

\end{note}

\begin{note}[来源路径]

近世代数/近世代数/无标题/定义1 20元素的阶 1afd2efa3f42804abe31e7c2129a941d.md

\end{note}

\section*{定义1.22 x的幂}

\begin{note}[来源上级主题]

1.3有限群 (1\%203\%E6\%9C\%89\%E9\%99\%90\%E7\%BE\%A4\%201ead2efa3f4280c39a49d8b0af9aec03.md)

\end{note}

\begin{definition}[定义1.22 x的幂]

x的幂是群同态f得是同态映射

\end{definition}

\begin{note}[图片 OCR 摘录，待人工复核]
以下内容来自源图片识别，便于审核时对照：

\textbf{图片 1}（IMG\_2945.jpeg）
\begin{Verbatim}[fontsize=\small]
fiz je: 1.22
4S (G,) RH, HBX EG. NS: (Z4) 9G), ZLA f(n) =x", R-DHAA.      .
MSE b, BAAR RE SG POIRIER, PAGS eM EJS Be I. PECL, FRAT lth LEX.
cE SL 1.22
&(G,) RPAH, HxeeG. $neN,, MAIL x= (x!)", FeV x =e.             .
SUE, BRAVA Hy Lik a ETE HA
UY) Re xeG. &mneZ, RA KRAGEW f(m+n) = f(m): f(n), WE x" ax” - x",
BEERS, MR mneEN,, UK LEASE RUA ER. % m Kn * 0, All FA AL 7c BY Ei
HABA. HRM TUARAAE, Bm <0, Wm!’ =—m, Mama” = (rel).
#n<0, in’ =—-n, WRB, x= (e)", Bama (1) RE mn EM, FERAIET.
HO<n<m, Mam sx = (NN mata)" Ol)", RE-E-h, AMER
i, tHE.
An>m', Ax FAR ('), AAWEARES HO AERAW, RERKIEM, CARS
HERE, DOHLREWISORANATR, AE, TSE LITA EE AT AEN. ae
wee [FH REAR TE BEET PRT IS. FR BIE, BORE “SEZ” AERO A IN LI, TE
\end{Verbatim}

\textbf{图片 2}（IMG\_2946.jpeg）
\begin{Verbatim}[fontsize=\small]
fi 1.23
4 (G,:) =F, xe G. &mneZ, Wx = (x")",                                      ‘
VED) TEA EAA. HERE.
RPK, AVSIM HER, FTA Ree HE, MAPLES, that
{x” sn € Z}                                          (1.45)
fe OS PRE, BRA HH x AE CEH.
\end{Verbatim}

\end{note}

\begin{note}[来源路径]

近世代数/近世代数/无标题/定义1 22 x的幂 1d1d2efa3f42801fbf26e402d99d9fbb.md

\end{note}

\section*{定义1.23由x生成子群}

\begin{note}[来源上级主题]

1.3有限群 (1\%203\%E6\%9C\%89\%E9\%99\%90\%E7\%BE\%A4\%201ead2efa3f4280c39a49d8b0af9aec03.md)

\end{note}

\begin{definition}[定义1.23由x生成子群]

定义1.6由子集生成的自幺半群

定义1.20元素的阶 阶只是n次方可以等于e没有说n次方的过程，这n个元素就是群中所有的元素

\end{definition}

\begin{note}[关联图片]
以下图片已纳入本条整理，当前版本优先依据 Markdown 文字整理，图片细节可在审核时对照：

1. IMG\_2947.jpeg\\
2. IMG\_2949.jpeg\\

\end{note}

\begin{note}[来源路径]

近世代数/近世代数/无标题/定义1 23由x生成子群 1d1d2efa3f428063a543c4b613e77eac.md

\end{note}

\section*{定义1.28商集与指数}

\begin{note}[来源上级主题]

1.3有限群 (1\%203\%E6\%9C\%89\%E9\%99\%90\%E7\%BE\%A4\%201ead2efa3f4280c39a49d8b0af9aec03.md)

\end{note}

\begin{definition}[定义1.28商集与指数]

命题1.3.3子群陪集要不然相等要不然无交

指数就是左配集个数

\end{definition}

\begin{note}[关联图片]
以下图片已纳入本条整理，当前版本优先依据 Markdown 文字整理，图片细节可在审核时对照：

1. IMG\_2984.jpeg\\

\end{note}

\begin{note}[来源路径]

近世代数/近世代数/无标题/定义1 28商集与指数 1d2d2efa3f4280779550ef45dc732f24.md

\end{note}

\section*{定义1.29正规子群}

\begin{note}[来源上级主题]

1.4正规子群 (1\%204\%E6\%AD\%A3\%E8\%A7\%84\%E5\%AD\%90\%E7\%BE\%A4\%201ead2efa3f428095a412e65ae0862a5d.md)

\end{note}

\begin{definition}[定义1.29正规子群]

得给商集 群结构，a和a‘是不同代表元，正规子群才是三角，普通子群只是小于等于

证明

\end{definition}

\begin{note}[关联图片]
以下图片已纳入本条整理，当前版本优先依据 Markdown 文字整理，图片细节可在审核时对照：

1. IMG\_2989.jpeg\\
2. IMG\_2990.jpeg\\

\end{note}

\begin{note}[来源路径]

近世代数/近世代数/无标题/定义1 29正规子群 1d2d2efa3f428009a478d7345a44f7a0.md

\end{note}

\section*{定义1.3群中的幂}

\begin{note}[来源上级主题]

1.1幺半群 (1\%201\%E5\%B9\%BA\%E5\%8D\%8A\%E7\%BE\%A4\%201cad2efa3f4280628789c382df95fdc5.md)

\end{note}

\begin{definition}[定义1.3群中的幂]

牢记幂只是把多个乘法换了一种写法表示了，乘法只是一种运算

\end{definition}

\begin{note}[关联图片]
以下图片已纳入本条整理，当前版本优先依据 Markdown 文字整理，图片细节可在审核时对照：

1. IMG\_2898.jpeg\\

\end{note}

\begin{note}[来源路径]

近世代数/近世代数/无标题/定义1 3群中的幂 1cad2efa3f4280fa8cbac5abca50f603.md

\end{note}

\section*{定义1.4幺半群的子群还是幺半群}

\begin{note}[来源上级主题]

1.1幺半群 (1\%201\%E5\%B9\%BA\%E5\%8D\%8A\%E7\%BE\%A4\%201cad2efa3f4280628789c382df95fdc5.md)

\end{note}

\begin{definition}[定义1.4幺半群的子群还是幺半群]

eC ML 1.4

\&(S,-) RP APR, ST CS, KAI (T, +) \%\&(S,-) HDF ASH, BeeT, AT ERAFHA, BP

ecT                                        (1.12)

Vx,yeET,xxyeT                                    (1.13)

\&

ee CARA MEI, SOUMEET PA ABE ERE PHF

fir jel 1.4

VEY) Ro TBH HRN, FRSA (AME) RHR, RRARNSHHMRREARLH. BH,

BeEtT SPCRR HR, SAMT RRB (TTR). BER, RW, CUTMAS PRARE

Bi, AMT T PAP eB H.

HAA, ZS EBERT MY RC “TRIAS”. TRASH E “PRS” . TA HEA

RW “REIRIAS” , “SRTAAS”, SSR. OW TILE pe fh Bk — EN (ATR), FTF AA a,

ALEC LA DER. GH, FRAT APES.

\end{definition}

\begin{note}[图片 OCR 摘录，待人工复核]
以下内容来自源图片识别，便于审核时对照：

\textbf{图片 1}（IMG\_2900.jpeg）
\begin{Verbatim}[fontsize=\small]
eC ML 1.4
&(S,-) RP APR, ST CS, KAI (T, +) %&(S,-) HDF ASH, BeeT, AT ERAFHA, BP
ecT                                        (1.12)
Vx,yeET,xxyeT                                    (1.13)
&
ee CARA MEI, SOUMEET PA ABE ERE PHF
fir jel 1.4
VEY) Ro TBH HRN, FRSA (AME) RHR, RRARNSHHMRREARLH. BH,
BeEtT SPCRR HR, SAMT RRB (TTR). BER, RW, CUTMAS PRARE
Bi, AMT T PAP eB H.
HAA, ZS EBERT MY RC “TRIAS”. TRASH E “PRS” . TA HEA
RW “REIRIAS” , “SRTAAS”, SSR. OW TILE pe fh Bk — EN (ATR), FTF AA a,
ALEC LA DER. GH, FRAT APES.
\end{Verbatim}

\end{note}

\begin{note}[来源路径]

近世代数/近世代数/无标题/定义1 4幺半群的子群还是幺半群 1cad2efa3f4280a186e5ff06d952436b.md

\end{note}

\section*{定义1.5幺半群的同态}

\begin{note}[来源上级主题]

1.1幺半群 (1\%201\%E5\%B9\%BA\%E5\%8D\%8A\%E7\%BE\%A4\%201cad2efa3f4280628789c382df95fdc5.md)

\end{note}

\begin{definition}[定义1.5幺半群的同态]

这里是把自然数中的x映射到了x的n次方这个元素是在另一个幺半群里面的

这两个要素证明的就是同态的条件

\end{definition}

\begin{note}[关联图片]
以下图片已纳入本条整理，当前版本优先依据 Markdown 文字整理，图片细节可在审核时对照：

1. IMG\_2902.jpeg\\

\end{note}

\begin{note}[来源路径]

近世代数/近世代数/无标题/定义1 5幺半群的同态 1cad2efa3f4280878f7ace8395a0a565.md

\end{note}

\section*{定义1.6由子集生成的自幺半群}

\begin{note}[来源上级主题]

1.1幺半群 (1\%201\%E5\%B9\%BA\%E5\%8D\%8A\%E7\%BE\%A4\%201cad2efa3f4280628789c382df95fdc5.md)

\end{note}

\begin{definition}[定义1.6由子集生成的自幺半群]

是包含A的所有子幺半群，如果A本身就是幺半群它不需要添加项就可以生成子幺半群。但是如果A本身元素构成不了幺半群你们就要引入新的元素显然新的元素+A是包含A的，所有A生成的子幺半群里面都有全部的A。所有的交集也就A+最少构成子幺半群需要添加的元素，这个最少代价形成的幺半群就是A生成的子幺半群

\end{definition}

\begin{note}[关联图片]
以下图片已纳入本条整理，当前版本优先依据 Markdown 文字整理，图片细节可在审核时对照：

1. IMG\_2904.jpeg\\

\end{note}

\begin{note}[来源路径]

近世代数/近世代数/无标题/定义1 6由子集生成的自幺半群 1cad2efa3f4280c8b8e4ca2c6bd26b3d.md

\end{note}

\section*{定义1.7幺半群的同构}

\begin{note}[来源上级主题]

1.1幺半群 (1\%201\%E5\%B9\%BA\%E5\%8D\%8A\%E7\%BE\%A4\%201cad2efa3f4280628789c382df95fdc5.md)

\end{note}

\begin{definition}[定义1.7幺半群的同构]

同构概念可以类比全等

双射，原函数同态

\end{definition}

\begin{note}[图片 OCR 摘录，待人工复核]
以下内容来自源图片识别，便于审核时对照：

\textbf{图片 1}（IMG\_2907.jpeg）
\begin{Verbatim}[fontsize=\small]
cE SL 1.7                                                 —
(Rik (S,-), (T,*) EAR AHH, BP SOT SPOR, BNI Sf R-PAERAM, Hf e—-PRM
Ht, Het AA.
f RH                              (1.17)
Vx,y € S, f(x-y) = f(x) * f(y)                      (1.18)
fle) =e’                    (1.19)
HY, e foe’ ZH (S,-) Fo (T, *) B12.                                                                   .
KR ie ES, TRIS AE. EE ERB, BE OSE RAR, AU HE. TT
FR ABORT f AERA, ARBOR fo! AERIS (A f ROU, RDA fo! ae FETE). RK
x ERM AN SP 3) tise 4 & EE I) PA ee et A EEA AS. AIT, BES AE RET AY Ae at H
PRET. AKI AE, REM LT A, PS, SEELEY. TEAM A SEE BOE FY
BEE HY ee FSR RE, FBR SEL, Te BETH PSE RA ETA AS.
fii jel 1.6
Gf:(S) 9 (1,9) RPA ERA, MPT OSR-PEZERAA. Ast, fp uap4tew
Hy
a
Sear bk, e-SUEW AE AEN. I EE EF EI Be PEI BCE A. 4A “TR” ei],
it “le” WAN SEG AKA, MS KA Te ARE. AE AA, KE ADM? ae aE
HAT RRE 5
\end{Verbatim}

\end{note}

\begin{note}[来源路径]

近世代数/近世代数/无标题/定义1 7幺半群的同构 1cbd2efa3f4280e49eafe83547a86380.md

\end{note}

\section*{定义1.8幺半群的逆元}

\begin{note}[来源上级主题]

1.2群 (1\%202\%E7\%BE\%A4\%201ecd2efa3f4280cf8df9c881d186f2a4.md)

\end{note}

\begin{definition}[定义1.8幺半群的逆元]

ET                            eee

A (S,-) EP AFH, ES. MANTRA VTA, HAMS

Aye S,x-y=y-x=e                        (1.20)

Sip y RARA x agtk, ite x!

\&

\end{definition}

\begin{note}[图片 OCR 摘录，待人工复核]
以下内容来自源图片识别，便于审核时对照：

\textbf{图片 1}（IMG\_2911.jpeg）
\begin{Verbatim}[fontsize=\small]
ET                            eee
A (S,-) EP AFH, ES. MANTRA VTA, HAMS
Aye S,x-y=y-x=e                        (1.20)
Sip y RARA x agtk, ite x!
&
\end{Verbatim}

\end{note}

\begin{note}[来源路径]

近世代数/近世代数/无标题/定义1 8幺半群的逆元 1cbd2efa3f4280ab8ab0c76329b68410.md

\end{note}

\section*{定义1.9幺半群加逆元就是群}

\begin{note}[来源上级主题]

1.2群 (1\%202\%E7\%BE\%A4\%201ecd2efa3f4280cf8df9c881d186f2a4.md)

\end{note}

\begin{definition}[定义1.9幺半群加逆元就是群]

这里有着为什么加逆符号这个需要改变x，y的位置

\end{definition}

\begin{note}[图片 OCR 摘录，待人工复核]
以下内容来自源图片识别，便于审核时对照：

\textbf{图片 1}（IMG\_2909.png）
\begin{Verbatim}[fontsize=\small]
=a                                      eee
(Gy) RRA FH, PN C R—-DF, BG PHA CHP TUN. REZ, SRG La—*
=A HH, WAI (G,-) EAA, RGR HMA, HRPABLHRASH, SEP, HADATE
FLAT. Fit—-bRA RIL, MASI, Fe G LAN—-PAAAH, MAMIE (G,-) EP, 4
Vx,y,2€G,x-(y-z)=(x-y)-z                           (1.22)
deeG,Vx €G,x-e=e-x=x                       (1.23)
Vx €G,dayeG,x-y=y-x=e                    (1.24)
&
YER, OPH, FOSSA G, DARBY ALK, BATE ScAE group. WHb, FRATGE x Wy wiocid
Vex! —“S ASRWOHEIE EE, BED TC RI BETC AME I (MXP BEAR BI PAHIC )
ABA Wi CHY Wie Ee — FE ce A AWE? RE EN
fir jel 1.8
&(G,) RH, Saxe G, Wel) =x.
ry
VED) WHR, RNAyax!, FHAx- yay-x=e. RNB y lax, WkHey- xax-y=e, LAR
Me RRIEAT RON REA, FOTANA HH.
fit jel: 1.9
&(G,) RMA, SxyeG, M(eyytaytat,                         .
UE RATA ERIE. —HH, AIA MAO, (wy) Cl x !)=e; F-FH, ABABA A-W
WEA, RE RR we yytsyt ext.
FONT RE ACHE, FRM AEs A CHEE. PRL SEERA DUAR OE ACHE OT HR, AC EEA
AA HS BN Beg ee bn] DL ANB
6
\end{Verbatim}

\end{note}

\begin{note}[来源路径]

近世代数/近世代数/无标题/定义1 9幺半群加逆元就是群 1cbd2efa3f4280789737e0626b159cd6.md

\end{note}

\section*{定义． 陪集个数为指数}

\begin{note}[来源上级主题]

子群的陪集 (\%E5\%AD\%90\%E7\%BE\%A4\%E7\%9A\%84\%E9\%99\%AA\%E9\%9B\%86\%201c3d2efa3f42806a85b0f0113028b207.md)

\end{note}

\begin{definition}[定义． 陪集个数为指数]

EM — PEG WPS FEA OR CBR ZEB SE) BY PhS HY A EG

ENE £8

P Wa Be ATT EF a Br Se RE BA LS BS eB. A eB SA eB SEY A

HE, Suze BOTT AP FA a Be Se 9 He Ze BY DA cB SR

\end{definition}

\begin{note}[图片 OCR 摘录，待人工复核]
以下内容来自源图片识别，便于审核时对照：

\textbf{图片 1}（image.png）
\begin{Verbatim}[fontsize=\small]
EM — PEG WPS FEA OR CBR ZEB SE) BY PhS HY A EG
ENE £8

P Wa Be ATT EF a Br Se RE BA LS BS eB. A eB SA eB SEY A
HE, Suze BOTT AP FA a Be Se 9 He Ze BY DA cB SR
\end{Verbatim}

\textbf{图片 2}（image 1.png）
\begin{Verbatim}[fontsize=\small]
S12 —} FH HSH DB—SAGR Ho Zi RES —— BRN.
WE AA                    g: h—>ha
Zz H 5 Ha fA —— PRY. A
G) A Kj —P 70h ASE WK ha
Gi) Ha K#—P otha 2H W7ch WR;
Git) RW hia=h,a, PHAh,=h2. ws.
AIX 5, RTWB RES HEH 2 AE 3.
\end{Verbatim}

\end{note}

\begin{note}[来源路径]

近世代数/近世代数/无标题/定义． 陪集个数为指数 1c3d2efa3f4280d4ae10d01394813e30.md

\end{note}

\section*{定理1 左右陪集个数相等，有限还相等或都是无限}

\begin{note}[来源上级主题]

子群的陪集 (\%E5\%AD\%90\%E7\%BE\%A4\%E7\%9A\%84\%E9\%99\%AA\%E9\%9B\%86\%201c3d2efa3f42806a85b0f0113028b207.md)

\end{note}

\begin{theorem}[定理1 左右陪集个数相等，有限还相等或都是无限]

先构造左陪集和右陪集的一个映射

i是证明a，b的选取不影响

\end{theorem}

\begin{note}[关联图片]
以下图片已纳入本条整理，当前版本优先依据 Markdown 文字整理，图片细节可在审核时对照：

1. image.png\\
2. image 1.png\\

\end{note}

\begin{note}[来源路径]

近世代数/近世代数/无标题/定理1 左右陪集个数相等，有限还相等或都是无限 1c3d2efa3f4280489fb7c3b3a78e65b1.md

\end{note}

\section*{定理1不变子群的判定}

\begin{theorem}[定理1不变子群的判定]

正规子群的意义就是为了让看看陪集是不是群

\end{theorem}

\begin{note}[图片 OCR 摘录，待人工复核]
以下内容来自源图片识别，便于审核时对照：

\textbf{图片 1}（image.png）
\begin{Verbatim}[fontsize=\small]
EH 1 —-THGW—-ThTHN B2—TAET HY SO Mh
Fh BE RE
aNa '=N
FGM TER— Phot a AX                                  |
\end{Verbatim}

\end{note}

\begin{note}[来源路径]

近世代数/近世代数/无标题/定理1不变子群的判定 1cad2efa3f42806f994edf8fc55ae209.md

\end{note}

\section*{定理1子群充要条件}

\begin{note}[来源上级主题]

子群 (\%E5\%AD\%90\%E7\%BE\%A4\%201bdd2efa3f42801b985de1cf59e5ed5f.md)

\end{note}

\begin{theorem}[定理1子群充要条件]

子群单位元是原群的单位元，子群的逆元也是原群的逆元

继承最重要的

\end{theorem}

\begin{note}[关联图片]
以下图片已纳入本条整理，当前版本优先依据 Markdown 文字整理，图片细节可在审核时对照：

1. image.png\\

\end{note}

\begin{note}[来源路径]

近世代数/近世代数/无标题/定理1子群充要条件 1bdd2efa3f428096a849dfb1aa61075c.md

\end{note}

\section*{定理2}

\begin{note}[来源上级主题]

子群的陪集 (\%E5\%AD\%90\%E7\%BE\%A4\%E7\%9A\%84\%E9\%99\%AA\%E9\%9B\%86\%201c3d2efa3f42806a85b0f0113028b207.md)

\end{note}

\begin{theorem}[定理2]

N=nj

\end{theorem}

\begin{note}[图片 OCR 摘录，待人工复核]
以下内容来自源图片识别，便于审核时对照：

\textbf{图片 1}（image.png）
\begin{Verbatim}[fontsize=\small]
EH 2 (he H E—-TARBRG M—TSFR. BAH WB n MEEG B
Wie 7 AREER C MPT N, FFA
N=n}.
\end{Verbatim}

\textbf{图片 2}（image 1.png）
\begin{Verbatim}[fontsize=\small]
1A GHUBIN MAAR, A Mbt n Mi Abe A RIE BA. GC hh
N *7cRA MARR, MAS. B—-TERBA n Soc, PU
e 54 «@
\end{Verbatim}

\textbf{图片 3}（image 2.png）
\begin{Verbatim}[fontsize=\small]
tH AlLagrange xe
cM RRG HRA HARA)
BEAN CA BR ERTCI) A HW FE G PAFR SK.
WHE (G: H),
fll G=S, H ={(1),(12)}
H(I)={(1),(02)        H(12)= (12) ()}.
(13) = {(13),(123)}                 (23) = {(23), (132)}
H(123)= {(123),(13)},                  H (132) = {(132),(23)}.
(G:H)=3
\end{Verbatim}

\end{note}

\begin{note}[来源路径]

近世代数/近世代数/无标题/定理2 1c3d2efa3f42808d86d9d910407f5566.md

\end{note}

\section*{定理2不变子群判定方式2}

\begin{theorem}[定理2不变子群判定方式2]

定理1不变子群的判定

\end{theorem}

\begin{note}[图片 OCR 摘录，待人工复核]
以下内容来自源图片识别，便于审核时对照：

\textbf{图片 1}（image.png）
\begin{Verbatim}[fontsize=\small]
EEL —MEGH-TTEN B-TAR FRR
ADEE:
a€G, nEN>ana'€N
EM ARIPO, JEL ALR. RITE
EWERAN. BEXTRA MU. BAM GC HE —*F 70 a
Kit
(1)               aNa"'CN
RH, Ana 'the GW, RNA
a 'NaCN, a (a™'!Na) a@'!CaNa™!
(2)               NCaNa™!
FR O(1) fl (2)       aNa7'=N
Am aeHE 1, NEROTH. IES
\end{Verbatim}

\end{note}

\begin{note}[来源路径]

近世代数/近世代数/无标题/定理2不变子群判定方式2 1cad2efa3f42808ca4e0c22a6b95ca07.md

\end{note}

\section*{定理2非空子集充要条件}

\begin{note}[来源上级主题]

子群 (\%E5\%AD\%90\%E7\%BE\%A4\%201bdd2efa3f42801b985de1cf59e5ed5f.md)

\end{note}

\begin{theorem}[定理2非空子集充要条件]

子群

定理3子群充要条件 非空有限子群

定理1子群充要条件

定理2就是简化的定理1

\end{theorem}

\begin{note}[关联图片]
以下图片已纳入本条整理，当前版本优先依据 Markdown 文字整理，图片细节可在审核时对照：

1. image.png\\

\end{note}

\begin{note}[来源路径]

近世代数/近世代数/无标题/定理2非空子集充要条件 1bdd2efa3f4280eca5fdd0e948316100.md

\end{note}

\section*{定理3子群充要条件 非空有限子群}

\begin{note}[来源上级主题]

子群 (\%E5\%AD\%90\%E7\%BE\%A4\%201bdd2efa3f42801b985de1cf59e5ed5f.md)

\end{note}

\begin{theorem}[定理3子群充要条件 非空有限子群]

EHS —tT\#HGCH-—TIES ARTE H ERG MNH—-hTBRW REA

4 Fe

a, be H>abc€ H.

\end{theorem}

\begin{note}[图片 OCR 摘录，待人工复核]
以下内容来自源图片识别，便于审核时对照：

\textbf{图片 1}（image.png）
\begin{Verbatim}[fontsize=\small]
EHS —tT#HGCH-—TIES ARTE H ERG MNH—-hTBRW REA
4 Fe
a, be H>abc€ H.
\end{Verbatim}

\end{note}

\begin{note}[来源路径]

近世代数/近世代数/无标题/定理3子群充要条件 非空有限子群 1bdd2efa3f4280bfbf36f4060b3612ad.md

\end{note}

\section*{对称群}

\begin{note}[来源上级主题]

置换群 (\%E7\%BD\%AE\%E6\%8D\%A2\%E7\%BE\%A4\%201bad2efa3f4280228c17c29195a2c525.md)

\end{note}

\begin{note}[对称群]

EM — PAA n SIMRAN EAB PRE BY HY i n WOT RR. XA

» AD «

\end{note}

\begin{note}[图片 OCR 摘录，待人工复核]
以下内容来自源图片识别，便于审核时对照：

\textbf{图片 1}（image.png）
\begin{Verbatim}[fontsize=\small]
EM — PAA n SIMRAN EAB PRE BY HY i n WOT RR. XA
» AD «
\end{Verbatim}

\textbf{图片 2}（image 1.png）
\begin{Verbatim}[fontsize=\small]
HIRATA S, RRA.

40 Se RR AG» on SUN RFA nh. BRA DIEW, RA
BYfEn SINTER, MER SRP we PMR, RATS a1 EH
AM, AnH RE, BET WI, Ba. HN, PRAn—-1 PARE, HF
A, FEA Bl

n(n—1)2*1l=n!
AAR TAA ER. DORE, BTA
\end{Verbatim}

\end{note}

\begin{note}[来源路径]

近世代数/近世代数/无标题/对称群 1bad2efa3f42809caaf1ff2a6c16b526.md

\end{note}

\section*{封闭 二元运算}

\begin{note}[来源上级主题]

代数运算 (\%E4\%BB\%A3\%E6\%95\%B0\%E8\%BF\%90\%E7\%AE\%97\%201aed2efa3f4280d597defe353a625299.md)

\end{note}

\begin{note}[封闭 二元运算]

乘法与乘积 运算的所有元素都是A里面的元素

实数集+-*是二元运算

二元运算也称为A上的代数运算 如在有理域上的代数运算

定义 对集合 S 进行二元运算 ∗ 是一个将 S×S 映射到 S 的函数。对于集合 S×S 中每对 （a,b），我们将用 a∗b 表示 S 中元素 ∗（（a,b））。 ■ 直观上，我们可以将 S 上二元运算 ∗ 视为将 S 的每个有序对 （a,b） 映射到 S 的元素 a∗b。 二元运算是指我们将 S 的元素对映射到 S。我们也可以将三元运算定义为将 S 的三元组映射到 S 的函数，但我们不需要这种类型的运算。在本书中，我们通常会省略二元运算这个词，使用运算这个词来表示二元运算。

实数域上+-✖️复合都是封闭————对于函数

\end{note}

\begin{note}[关联图片]
以下图片已纳入本条整理，当前版本优先依据 Markdown 文字整理，图片细节可在审核时对照：

1. image.png\\

\end{note}

\begin{note}[来源路径]

近世代数/近世代数/无标题/封闭 二元运算 1a1d2efa3f4280c1abc9d4dc388cc80c.md

\end{note}

\section*{左陪集}

\begin{note}[来源上级主题]

子群的陪集 (\%E5\%AD\%90\%E7\%BE\%A4\%E7\%9A\%84\%E9\%99\%AA\%E9\%9B\%86\%201c3d2efa3f42806a85b0f0113028b207.md)

\end{note}

\begin{note}[左陪集]

A ERE ER RAR \textasciitilde{} :

a\textasciitilde{}b 4AM ab" CH it

HAMAS. EMR MMe—-SRA\textasciitilde{}’:

a\textasciitilde{}'b SAM4b acH i,

PRARMWE—-RAWBH, \textasciitilde{}hE—-TSER KA. MAKTSHRKA, RMN

AAS G HA-*FAR.

EM HSRKA\textasciitilde{} MRENKUYMTH HA MAB. BAe HA

AAS aH RRM.

FL, E— FER WE: aH WIA MANUS

ah (hEH)

ESA G ATC.

\end{note}

\begin{note}[图片 OCR 摘录，待人工复核]
以下内容来自源图片识别，便于审核时对照：

\textbf{图片 1}（image.png）
\begin{Verbatim}[fontsize=\small]
A ERE ER RAR ~ :
a~b 4AM ab" CH it
HAMAS. EMR MMe—-SRA~’:
a~'b SAM4b acH i,
PRARMWE—-RAWBH, ~hE—-TSER KA. MAKTSHRKA, RMN
AAS G HA-*FAR.
EM HSRKA~ MRENKUYMTH HA MAB. BAe HA
AAS aH RRM.
FL, E— FER WE: aH WIA MANUS
ah (hEH)
ESA G ATC.
\end{Verbatim}

\end{note}

\begin{note}[来源路径]

近世代数/近世代数/无标题/左陪集 1c3d2efa3f4280bebcadc5fc4efa6e1e.md

\end{note}

\section*{左陪集定义}

\begin{note}[来源上级主题]

1.3有限群 (1\%203\%E6\%9C\%89\%E9\%99\%90\%E7\%BE\%A4\%201ead2efa3f4280c39a49d8b0af9aec03.md)

\end{note}

\begin{definition}[左陪集定义]

\& GRA, H<G RY TH, AEG. NaH RH MH AK (Ha sl), CLA

aH = \{ax : x € H\}                                    (1.55) *

thteb, Ha STH A AYZzCBASR, Wi a Ace A PCR OTE. ps, eH = A th

Fe PACS, AREA AE, TE A ETE RLS REI AE, FATT VA rE SOB SEY ARIA , TT CENT IC MERAY BRE

eH = A yee BI PACED

FRNTARAL LAT BREE) Ze FE

\end{definition}

\begin{note}[图片 OCR 摘录，待人工复核]
以下内容来自源图片识别，便于审核时对照：

\textbf{图片 1}（image.png）
\begin{Verbatim}[fontsize=\small]
& GRA, H<G RY TH, AEG. NaH RH MH AK (Ha sl), CLA
aH = {ax : x € H}                                    (1.55) *
thteb, Ha STH A AYZzCBASR, Wi a Ace A PCR OTE. ps, eH = A th
Fe PACS, AREA AE, TE A ETE RLS REI AE, FATT VA rE SOB SEY ARIA , TT CENT IC MERAY BRE
eH = A yee BI PACED
FRNTARAL LAT BREE) Ze FE
\end{Verbatim}

\end{note}

\begin{note}[来源路径]

近世代数/近世代数/无标题/左陪集定义 1d2d2efa3f4280fdbd2fd7e28903354a.md

\end{note}

\section*{幺半群的定义 定义1.2 定义1.3}

\begin{note}[来源上级主题]

1.1幺半群 (1\%201\%E5\%B9\%BA\%E5\%8D\%8A\%E7\%BE\%A4\%201cad2efa3f4280628789c382df95fdc5.md)

\end{note}

\begin{definition}[幺半群的定义 定义1.2 定义1.3]

子级 项目: 单位元唯一性证明 (\%E5\%8D\%95\%E4\%BD\%8D\%E5\%85\%83\%E5\%94\%AF\%E4\%B8\%80\%E6\%80\%A7\%E8\%AF\%81\%E6\%98\%8E\%201cad2efa3f428021acabe9e17d5b39d5.md)

幺半群只要求有单位元，还有满足结合律

幺半群的单位元是唯一的

\end{definition}

\begin{note}[关联图片]
以下图片已纳入本条整理，当前版本优先依据 Markdown 文字整理，图片细节可在审核时对照：

1. IMG\_2894.jpeg\\

\end{note}

\begin{note}[来源路径]

近世代数/近世代数/无标题/幺半群的定义 定义1 2 定义1 3 1cad2efa3f42800ea08fdf383be45a9a.md

\end{note}

\section*{广义结合律与数学归纳法}

\begin{note}[来源上级主题]

1.1幺半群 (1\%201\%E5\%B9\%BA\%E5\%8D\%8A\%E7\%BE\%A4\%201cad2efa3f4280628789c382df95fdc5.md)

\end{note}

\begin{note}[广义结合律与数学归纳法]

数学归纳法先对1一个证明成立然后假设k成立证明k成立的时候可以借鉴证明特殊的做法再证明k+1成立，因为k就是1为了推广写成k

这里需要证明的原因就是为了1.3命题为什么这里需要证明因为针对不是加法或者乘法而是所有的二元运算包括自定义的

\end{note}

\begin{note}[关联图片]
以下图片已纳入本条整理，当前版本优先依据 Markdown 文字整理，图片细节可在审核时对照：

1. IMG\_2899.jpeg\\

\end{note}

\begin{note}[来源路径]

近世代数/近世代数/无标题/广义结合律与数学归纳法 1cad2efa3f42803e8230cae3ea9f46af.md

\end{note}

\section*{引理1.1幺半群的逆元构成群}

\begin{note}[来源上级主题]

1.10阿贝尔群 (1\%2010\%E9\%98\%BF\%E8\%B4\%9D\%E5\%B0\%94\%E7\%BE\%A4\%201cbd2efa3f4280bfa73dc117dd0d632c.md)

\end{note}

\begin{lemma}[引理1.1幺半群的逆元构成群]

5 BH 1.1

FRA TERE A ABR PY AT BCR “PL”, AE G ATA Ge PRA (tei EE APE).

WEY) PREOERSRAK ES, REE. WEST, AlkhecG. MP ZEWG PETRA

(Gay) Ht, MILF ARAN. Mix eG, Wx 2TH CA, KRiRyeS, Hx -yay-x=eRLE

HERRIRRAAREy ELSES +). RPRANEEH yeG, Wy TH, WMRELAW, AAr ESCH

wo Phy eG. 3x, MILT (G,-) eH.

Bilin, FERMI, (Z,-) SrA Awaz NEAL (+1,-), SEAS, (A Bee.

\end{lemma}

\begin{note}[图片 OCR 摘录，待人工复核]
以下内容来自源图片识别，便于审核时对照：

\textbf{图片 1}（image.png）
\begin{Verbatim}[fontsize=\small]
5 BH 1.1

FRA TERE A ABR PY AT BCR “PL”, AE G ATA Ge PRA (tei EE APE).
WEY) PREOERSRAK ES, REE. WEST, AlkhecG. MP ZEWG PETRA
(Gay) Ht, MILF ARAN. Mix eG, Wx 2TH CA, KRiRyeS, Hx -yay-x=eRLE
HERRIRRAAREy ELSES +). RPRANEEH yeG, Wy TH, WMRELAW, AAr ESCH
wo Phy eG. 3x, MILT (G,-) eH.

Bilin, FERMI, (Z,-) SrA Awaz NEAL (+1,-), SEAS, (A Bee.
\end{Verbatim}

\end{note}

\begin{note}[来源路径]

近世代数/近世代数/无标题/引理1 1幺半群的逆元构成群 1edd2efa3f4280479e2be2b835b83ca6.md

\end{note}

\section*{引理1.4正规子群}

\begin{note}[来源上级主题]

1.4正规子群 (1\%204\%E6\%AD\%A3\%E8\%A7\%84\%E5\%AD\%90\%E7\%BE\%A4\%201ead2efa3f428095a412e65ae0862a5d.md)

\end{note}

\begin{lemma}[引理1.4正规子群]

3) 5H 1.4                                         —

4 (G,:) RAF, ALN <G, MFA PRS

1 NGA EMR, BP

Vaé€G,aN = Na                                       (1.72)

2.

Va € G,aNa"! CN                                    (1.73)

3.

Va € G,W¥n € N,ana"!€N                                   (1.74)  5

UEW) BARBOSA EMREALESH. RN REEAB—-DAEK RIKI APEI

—FH, RIVRRNEGHEMTH. SacG, MlaN=Na. ARARa!|HH-+HEASCKA, RMN

4245] Y aNa! CN.

B-HB, RNMARBXIDAH. Sa G, I aNa! CN BB) aN CNa, aN (a) CN Al

NacaN,. Aik, aN=Na.

SESE, GAO APPA BERTRAM, ESHER — SAR EAR

— aL ae, WR TE LPP FE RIK EPL

\end{lemma}

\begin{note}[图片 OCR 摘录，待人工复核]
以下内容来自源图片识别，便于审核时对照：

\textbf{图片 1}（image.png）
\begin{Verbatim}[fontsize=\small]
3) 5H 1.4                                         —
4 (G,:) RAF, ALN <G, MFA PRS
1 NGA EMR, BP
Vaé€G,aN = Na                                       (1.72)
2.
Va € G,aNa"! CN                                    (1.73)
3.
Va € G,W¥n € N,ana"!€N                                   (1.74)  5
UEW) BARBOSA EMREALESH. RN REEAB—-DAEK RIKI APEI
—FH, RIVRRNEGHEMTH. SacG, MlaN=Na. ARARa!|HH-+HEASCKA, RMN
4245] Y aNa! CN.
B-HB, RNMARBXIDAH. Sa G, I aNa! CN BB) aN CNa, aN (a) CN Al
NacaN,. Aik, aN=Na.
SESE, GAO APPA BERTRAM, ESHER — SAR EAR
— aL ae, WR TE LPP FE RIK EPL
\end{Verbatim}

\end{note}

\begin{note}[来源路径]

近世代数/近世代数/无标题/引理1 4正规子群 1d8d2efa3f4280599a3fd672105fee0c.md

\end{note}

\section*{拉格朗日定理}

\begin{theorem}[拉格朗日定理]

jESE 2 (Lagrangese HE) ARH CG ,

H<G> W |el=|a\textbackslash{}(G:A) .

WER AN a <Gc, PO BRA,

Ail W fEG PARR ANISAHBARE. i

(G:H)=rhG=a,H Va,HU---UaH

HEBEL, | a,H |=|4,H |= ||

Ary,

|G |=|4,H |+|a,H |+--+ 4,0 |=r|H |= |A\textbackslash{}(G: H)

\end{theorem}

\begin{note}[图片 OCR 摘录，待人工复核]
以下内容来自源图片识别，便于审核时对照：

\textbf{图片 1}（image.png）
\begin{Verbatim}[fontsize=\small]
jESE 2 (Lagrangese HE) ARH CG ,
H<G> W |el=|a\(G:A) .
WER AN a <Gc, PO BRA,
Ail W fEG PARR ANISAHBARE. i
(G:H)=rhG=a,H Va,HU---UaH
HEBEL, | a,H |=|4,H |= ||
Ary,
|G |=|4,H |+|a,H |+--+ 4,0 |=r|H |= |A\(G: H)
\end{Verbatim}

\end{note}

\begin{note}[来源路径]

近世代数/近世代数/无标题/拉格朗日定理 1cad2efa3f4280aab6f1ee857a182752.md

\end{note}

\section*{拉格朗日定理}

\begin{note}[来源上级主题]

1.3有限群 (1\%203\%E6\%9C\%89\%E9\%99\%90\%E7\%BE\%A4\%201ead2efa3f4280c39a49d8b0af9aec03.md)

\end{note}

\begin{theorem}[拉格朗日定理]

G的阶等于商群的数量成子群的阶比如

循环群\{0，1，2，3，4，5\} 子群\{0,3\} 他生成的商群\{1,4\} \{2,5\} \{0,3\} 6=3*2

子群的阶一定是父群阶的因数

命题1.30循环群的生成元的幂指数和循环群阶数互素 生成元的幂指数和子群的阶完全不是一个东西

\end{theorem}

\begin{note}[关联图片]
以下图片已纳入本条整理，当前版本优先依据 Markdown 文字整理，图片细节可在审核时对照：

1. IMG\_2985.jpeg\\
2. IMG\_2986.jpeg\\

\end{note}

\begin{note}[来源路径]

近世代数/近世代数/无标题/拉格朗日定理 1d2d2efa3f42801fa5fdf15cafbdcb4e.md

\end{note}

\section*{指数是2的子群一定是不变的}

\begin{note}[指数是2的子群一定是不变的]

也就是要不然g就在eN里面左乘右乘没区别，要不然因为已经有一个eN了右陪集，左陪集只有一个一个座位所以他们只能是一个人

\end{note}

\begin{note}[关联图片]
以下图片已纳入本条整理，当前版本优先依据 Markdown 文字整理，图片细节可在审核时对照：

1. IMG\_2888.jpeg\\

\end{note}

\begin{note}[来源路径]

近世代数/近世代数/无标题/指数是2的子群一定是不变的 1cad2efa3f4280a49122f16b85f73100.md

\end{note}

\section*{无标题}

\begin{note}[来源上级主题]

变换群证明 (\%E5\%8F\%98\%E6\%8D\%A2\%E7\%BE\%A4\%E8\%AF\%81\%E6\%98\%8E\%201b5d2efa3f42802e858cebf98bd17fd6.md)

\end{note}

\begin{note}[无标题]

待根据图片进一步人工校对。

\end{note}

\begin{note}[来源路径]

近世代数/近世代数/无标题/无标题 1b7d2efa3f42802dbdaecc7643526bbe.md

\end{note}

\section*{无标题}

\begin{note}[来源上级主题]

变换群 (\%E5\%8F\%98\%E6\%8D\%A2\%E7\%BE\%A4\%201b5d2efa3f4280bc8c6bf24cc87fed3c.md)

\end{note}

\begin{note}[无标题]

待根据图片进一步人工校对。

\end{note}

\begin{note}[来源路径]

近世代数/近世代数/无标题/无标题 1bad2efa3f42802aa86ec37cd6544448.md

\end{note}

\section*{无标题}

\begin{note}[来源上级主题]

变换群 (\%E5\%8F\%98\%E6\%8D\%A2\%E7\%BE\%A4\%201b5d2efa3f4280bc8c6bf24cc87fed3c.md)

\end{note}

\begin{note}[无标题]

待根据图片进一步人工校对。

\end{note}

\begin{note}[来源路径]

近世代数/近世代数/无标题/无标题 1bad2efa3f42803389efdacc7eedffb9.md

\end{note}

\section*{无标题}

\begin{note}[来源上级主题]

群同构第三定理 (\%E7\%BE\%A4\%E5\%90\%8C\%E6\%9E\%84\%E7\%AC\%AC\%E4\%B8\%89\%E5\%AE\%9A\%E7\%90\%86\%201ead2efa3f428024abb0fc36b0b0317a.md)

\end{note}

\begin{note}[无标题]

待根据图片进一步人工校对。

\end{note}

\begin{note}[来源路径]

近世代数/近世代数/无标题/无标题 1ead2efa3f4280d8976efa8f012a8275.md

\end{note}

\section*{映射与集合基础知识}

\begin{note}[来源上级主题]

基本概念 (\%E5\%9F\%BA\%E6\%9C\%AC\%E6\%A6\%82\%E5\%BF\%B5\%201c3d2efa3f4280648bdcd081c8508c78.md)

\end{note}

\begin{note}[映射与集合基础知识]

子级 项目: 群 (\%E7\%BE\%A4\%201a0d2efa3f4280e5ab23eef494b59a4c.md), 集合的积 (\%E9\%9B\%86\%E5\%90\%88\%E7\%9A\%84\%E7\%A7\%AF\%201a1d2efa3f4280f19e43c59633af2d85.md), 集合的映射 (\%E9\%9B\%86\%E5\%90\%88\%E7\%9A\%84\%E6\%98\%A0\%E5\%B0\%84\%201a1d2efa3f428022b697d8ac2040bd12.md), 单射满射 (\%E5\%8D\%95\%E5\%B0\%84\%E6\%BB\%A1\%E5\%B0\%84\%201a1d2efa3f4280a7ab9bf96427741ce6.md), 同态的定义 (\%E5\%90\%8C\%E6\%80\%81\%E7\%9A\%84\%E5\%AE\%9A\%E4\%B9\%89\%201a1d2efa3f428051bf6ddd4a0bc015cf.md), 同构 (\%E5\%90\%8C\%E6\%9E\%84\%201a7d2efa3f4280a5b953fd47d4b69430.md)

\end{note}

\begin{note}[来源路径]

近世代数/近世代数/无标题/映射与集合基础知识 1aed2efa3f42800c8261da73ac826269.md

\end{note}

\section*{映射的核与像}

\begin{note}[来源上级主题]

群的同态同构 (\%E7\%BE\%A4\%E7\%9A\%84\%E5\%90\%8C\%E6\%80\%81\%E5\%90\%8C\%E6\%9E\%84\%201c3d2efa3f428076a462f0397fbfd469.md)

\end{note}

\begin{note}[映射的核与像]

2M 1.15                                                    =“

\& f i (G,-) > (G’,*) Z-PARAA, MAMIE f KER, ite ker(f) 5 im(f), HA

ker(f) = \{x €G: f(x) =e\} CG                      (1.35)

im(f) = \{y €G’: ae EG, y= f(x\} =\{f(@®) :x Ee G\}] CG’                    (1.36)

\&

FIX, FBS Bere CEA, MEERA. Kink, Boe (Hh, PUTRI TASH.

TEMUWIA DENY, Bea Be BEEN.

9

1.2 \#\#

,

G                                                G’

Al 1.1: BSRWKAREEA.

\end{note}

\begin{note}[图片 OCR 摘录，待人工复核]
以下内容来自源图片识别，便于审核时对照：

\textbf{图片 1}（IMG\_2924.png）
\begin{Verbatim}[fontsize=\small]
2M 1.15                                                    =“
& f i (G,-) > (G’,*) Z-PARAA, MAMIE f KER, ite ker(f) 5 im(f), HA

ker(f) = {x €G: f(x) =e} CG                      (1.35)

im(f) = {y €G’: ae EG, y= f(x} ={f(@®) :x Ee G}] CG’                    (1.36)
&
FIX, FBS Bere CEA, MEERA. Kink, Boe (Hh, PUTRI TASH.

TEMUWIA DENY, Bea Be BEEN.
9
1.2 ##
,
G                                                G’
Al 1.1: BSRWKAREEA.
\end{Verbatim}

\end{note}

\begin{note}[来源路径]

近世代数/近世代数/无标题/映射的核与像 1cbd2efa3f428060a500d68691e7bbdf.md

\end{note}

\section*{有限群每一个元素都有元}

\begin{note}[来源上级主题]

定义1.20元素的阶 (\%E5\%AE\%9A\%E4\%B9\%891\%2020\%E5\%85\%83\%E7\%B4\%A0\%E7\%9A\%84\%E9\%98\%B6\%201afd2efa3f42804abe31e7c2129a941d.md)

\end{note}

\begin{note}[有限群每一个元素都有元]

不能不相等是因为如果都不相等就会出现无限了

\end{note}

\begin{note}[图片 OCR 摘录，待人工复核]
以下内容来自源图片识别，便于审核时对照：

\textbf{图片 1}（image.png）
\begin{Verbatim}[fontsize=\small]
4. 74 PY — PCY BT ob PR.
fe SGHR-TARRM CORGHE—-UR, BA
FRMAA. RUGELILME i fDi, Bat=ol. Fy
oi Rem,         -
(1)         ai ize
poy, FPTEIEM A I-7, fh (1) RO. Bee
ERA, fF a"=e, XP, TOMBE.
\end{Verbatim}

\end{note}

\begin{note}[来源路径]

近世代数/近世代数/无标题/有限群每一个元素都有元 1b1d2efa3f428077b99af9bef5e16a15.md

\end{note}

\section*{构造同态同构}

\begin{note}[来源上级主题]

变换群 (\%E5\%8F\%98\%E6\%8D\%A2\%E7\%BE\%A4\%201c3d2efa3f4280a3a416f8493ad4dae2.md)

\end{note}

\begin{note}[构造同态同构]

G@ 2H, AMRARHKRERTUR—H, HRBEBMBHR AME, EURFET

AIC B25 ON TREK.

RR

1. fel\# (Homomorphism)

» ABBR RRA:

o(a-b) = \$(a) * o(b)

- wig, (R. +) A(R.) BBA), (ABATED SAAR SY O(w) = CER EA,

(FHA.

2. fal\#3 (Isomorphism)

> RI BARRA, PRADA.

HF

1. ZRF REFERS

- 18 (Z, +) BR MIDAR, (Z/rZ, +) BH n MKB.

» MO: Z92/nZ 414 2modn,

- REMEABAIDA, (EZ BIDAR, Mite 2/2 BIRKS MN.

- HAMM, PENERY (ANEREHH).

2. BETRERERM

» (R, +) 5 (R", x) BAAR, (BERR O(@) =e BT:

(a +b) =e =e" eb = (a) - 6b)

3. GRAFRAAES

-@G=(Z,+) Mm H=(Z,\textasciitilde{}x),

-@GH1+1=2, Ge AH, 1) O01) DASF 2), BBM IIARIE (IRIE

FES TRIR ST, OES AREY).

+ BTCA Bl iE ARE ASIEF MERRY

BS

-HYATRANRGR, HRWEHR SREY, YAIR BAKA.

-ASTERMH, (WM BAKA.

-AMERMS, BME BRAARH RM, MPSS Tel, tA

jZILEKA.

UMRMABRAHRS MBS, TH -SRTS REVERSE.

\end{note}

\begin{note}[图片 OCR 摘录，待人工复核]
以下内容来自源图片识别，便于审核时对照：

\textbf{图片 1}（IMG\_2708.jpeg）
\begin{Verbatim}[fontsize=\small]
G@ 2H, AMRARHKRERTUR—H, HRBEBMBHR AME, EURFET
AIC B25 ON TREK.
RR
1. fel# (Homomorphism)
» ABBR RRA:
o(a-b) = $(a) * o(b)
- wig, (R. +) A(R.) BBA), (ABATED SAAR SY O(w) = CER EA,
(FHA.
2. fal#3 (Isomorphism)
> RI BARRA, PRADA.
HF
1. ZRF REFERS
- 18 (Z, +) BR MIDAR, (Z/rZ, +) BH n MKB.
» MO: Z92/nZ 414 2modn,
- REMEABAIDA, (EZ BIDAR, Mite 2/2 BIRKS MN.
- HAMM, PENERY (ANEREHH).
2. BETRERERM
» (R, +) 5 (R", x) BAAR, (BERR O(@) =e BT:
(a +b) =e =e" eb = (a) - 6b)
3. GRAFRAAES
-@G=(Z,+) Mm H=(Z,~x),
-@GH1+1=2, Ge AH, 1) O01) DASF 2), BBM IIARIE (IRIE
FES TRIR ST, OES AREY).
+ BTCA Bl iE ARE ASIEF MERRY
BS
-HYATRANRGR, HRWEHR SREY, YAIR BAKA.
-ASTERMH, (WM BAKA.
-AMERMS, BME BRAARH RM, MPSS Tel, tA
jZILEKA.
UMRMABRAHRS MBS, TH -SRTS REVERSE.
\end{Verbatim}

\end{note}

\begin{note}[来源路径]

近世代数/近世代数/无标题/构造同态同构 1b5d2efa3f42802ba396ef3b35139369.md

\end{note}

\section*{正规子群定义}

\begin{definition}[正规子群定义]

循环群 Na=aN是用于验证的

任何两个群的两个平凡子群都是不变子群

交换群的子群都是不变子群

素数阶的群都是循环群

代表选取不影响是重点

\end{definition}

\begin{note}[关联图片]
以下图片已纳入本条整理，当前版本优先依据 Markdown 文字整理，图片细节可在审核时对照：

1. image.png\\
2. IMG\_2882.jpeg\\

\end{note}

\begin{note}[来源路径]

近世代数/近世代数/无标题/正规子群定义 1cad2efa3f4280be908adb4819bd31ae.md

\end{note}

\section*{正规子群的意义就是为了让看看陪集是不是群}

\begin{note}[正规子群的意义就是为了让看看陪集是不是群]

加群逆元就是相反数，乘群逆元就是倒数

\end{note}

\begin{note}[图片 OCR 摘录，待人工复核]
以下内容来自源图片识别，便于审核时对照：

\textbf{图片 1}（IMG\_2883.jpeg）
\begin{Verbatim}[fontsize=\small]
28H Z7(BRE) RAFH 3Z = {..., 6, -3,0, 3,6, ...}

-WE-BM o, ERS 0 +3Z RTMBS o ARF a (3) HBR, GI, 4
a=1ff,

1432 = {...,—-5,—2,1,4,7,...}.

- MED A—TEM a =4, $B4518R83, AN41-1=-3 BF 32, As

44+ 3Z=1432Z.

SUHEA 1 Al 4 ATEN BE 1432 ORT. HRS, YARNS ot =a

BY (XB SE 1+3Z2=4+3Z), RNRANE 154 RF Re, SB

HiRBEMKRI.

-EZH, SH=3Z,

.1+H= {1,4,7,...,-2,—5,...};

44+ H= {4,7,10,...,1,—-2,...};

~>AF 1+ H=44+4, P14 PEMA.

LML—TAFRERAKRTHA AGF.
\end{Verbatim}

\textbf{图片 2}（IMG\_1075.jpeg）
\begin{Verbatim}[fontsize=\small]
|               -           s         a |  . Ce  C                                                      sa                                                           mie
|                aa                                               -                       ciSthS-n-
has,                                                           -             =                                        ——
2
-                                           =<                                                                                       |
:           mee   @ B1Bn  eH     @ 3H B       |
-                          ai     ol                  ii
EEE                                                   <
a 0  na)  =¢   H  a ei  =>      d  |      i= oO | h   {                                   ———.                 ;
& late eo  ee Jw eet.   iL '     ) <—_ ae KE : LLS Phin jy a    >)  am   “    ——                        XY  4         +   ay
eee re     Be                                                                                                                                ty
——      is    -   oe     cee tieenengtilindie Sameeenetilines taal  eet ee     tt tw.    mt                                                               :   ,                        “7                                                                                                     j u
V-                             WVih               j             j                                                                                                                    i
oo                        iM. « a4    ae       pee os   a      aay   ‘ee ee Bek                                                                                                           oor                        i    ‘      ~*~,
Baar >d            4   Fe Sh    pea Sa ee Nee! ae    A     _  ee,     ——      see ne   oy pa                       aie               i      eto ai                       f                    |
[= Se. Goa [te | Ey,               ame        |
.          TD pret ee      Fic         4                                                                                          come
J                                                             iss        3  See | ls, j PD oe   hue bi  ve J]   P  ian “4                                                                                                        Zo                              wt
—    ——                              }      Re            L                Poe                 ni re 5                                  ‘       \      ~
\end{Verbatim}

\textbf{图片 3}（IMG\_2884.jpeg）
\begin{Verbatim}[fontsize=\small]
RGAH, HAGWT#H, AaH=d'H, RRRBa5¢0 BFR-TABR, hit
Bik, CNS HEX FES.

BRA, aH = d HNS RAE

aH=adH = a'a eH.

RA a TUS Ka S HR SBA, Bla’ =ahtERTR he A,
RF, RADAR o A’ SARE CH NNR. Hein, FR
BRANT ARAL EAIZBEN RI, ENZEAARAKAZD.

Alt,cH = HWSatte2 oS ATe—BR, HMA ENZBRWCRT.
\end{Verbatim}

\end{note}

\begin{note}[来源路径]

近世代数/近世代数/无标题/正规子群的意义就是为了让看看陪集是不是群 1cad2efa3f42800b9c25e9e1bc63b83c.md

\end{note}

\section*{每一个分类都决定了A的一个等价关系}

\begin{note}[来源上级主题]

分类与关系 (\%E5\%88\%86\%E7\%B1\%BB\%E4\%B8\%8E\%E5\%85\%B3\%E7\%B3\%BB\%201aed2efa3f42802491cdeb2b56b5f1f7.md)

\end{note}

\begin{note}[每一个分类都决定了A的一个等价关系]

等价关系

产生一个分类就能产生一个等价关系

因为一个分类里面都是有同一种性质的元素

把这种性质作为等价的判定自然能够构造出一种等价关系

\end{note}

\begin{note}[关联图片]
以下图片已纳入本条整理，当前版本优先依据 Markdown 文字整理，图片细节可在审核时对照：

1. image.png\\

\end{note}

\begin{note}[来源路径]

近世代数/近世代数/无标题/每一个分类都决定了A的一个等价关系 1a7d2efa3f4280eeabf8d5f851e06cc2.md

\end{note}

\section*{消去律成立与证明G是一个群 重点}

\begin{note}[来源上级主题]

群的基础知识 (\%E7\%BE\%A4\%E7\%9A\%84\%E5\%9F\%BA\%E7\%A1\%80\%E7\%9F\%A5\%E8\%AF\%86\%201aed2efa3f428009ac4eca74062880ff.md)

\end{note}

\begin{proof}

封闭有限集 有限有限不需要找单位元和逆元

消去律是未来解方程

群可以用消去律

\end{proof}

\begin{note}[关联图片]
以下图片已纳入本条整理，当前版本优先依据 Markdown 文字整理，图片细节可在审核时对照：

1. image.png\\
2. IMG\_0586.jpeg\\

\end{note}

\begin{note}[来源路径]

近世代数/近世代数/无标题/消去律成立与证明G是一个群 重点 1aed2efa3f42802bbde2ff09c5d7753f.md

\end{note}

\section*{看maki的习题课是最好的，构建关系}

\begin{note}[看maki的习题课是最好的，构建关系]

待根据图片进一步人工校对。

\end{note}

\begin{note}[来源路径]

近世代数/近世代数/无标题/看maki的习题课是最好的，构建关系 1c3d2efa3f42800db761fc4973d5b01e.md

\end{note}

\section*{等价关系}

\begin{note}[来源上级主题]

分类与关系 (\%E5\%88\%86\%E7\%B1\%BB\%E4\%B8\%8E\%E5\%85\%B3\%E7\%B3\%BB\%201aed2efa3f42802491cdeb2b56b5f1f7.md)

\end{note}

\begin{note}[等价关系]

关系的定义 二元关系

aRa

aRb 则bRa

aRc bRa 则bRc  显然这里是R代表更广义的=除了等于的意思还表示（a，c）属于集合R

比如大于号就不是等价关系

\end{note}

\begin{note}[关联图片]
以下图片已纳入本条整理，当前版本优先依据 Markdown 文字整理，图片细节可在审核时对照：

1. image.png\\

\end{note}

\begin{note}[来源路径]

近世代数/近世代数/无标题/等价关系 1a7d2efa3f4280a0ae90db267cbf3328.md

\end{note}

\section*{结合律 交换律 分配律}

\begin{note}[来源上级主题]

代数运算 (\%E4\%BB\%A3\%E6\%95\%B0\%E8\%BF\%90\%E7\%AE\%97\%201aed2efa3f4280d597defe353a625299.md)

\end{note}

\begin{note}[结合律 交换律 分配律]

https://www.notion.so

保证映射就行 第一分配律把z分进去左侧 第二右侧

\end{note}

\begin{note}[关联图片]
以下图片已纳入本条整理，当前版本优先依据 Markdown 文字整理，图片细节可在审核时对照：

1. image.png\\
2. image 1.png\\

\end{note}

\begin{note}[来源路径]

近世代数/近世代数/无标题/结合律 交换律 分配律 1a1d2efa3f428005bcfac88aa2af2b6e.md

\end{note}

\section*{置换与置换群}

\begin{note}[来源上级主题]

置换群 (\%E7\%BD\%AE\%E6\%8D\%A2\%E7\%BE\%A4\%201bad2efa3f4280228c17c29195a2c525.md)

\end{note}

\begin{note}[置换与置换群]

有限集+变换=置换

一一变化为置换

\end{note}

\begin{note}[图片 OCR 摘录，待人工复核]
以下内容来自源图片识别，便于审核时对照：

\textbf{图片 1}（image.png）
\begin{Verbatim}[fontsize=\small]
ex — PARR —S—— BRM K—-T ER.

— 45 BRR HY FS ERE A — EY i — PS BRR.

RINB—TA RSG A, A An WSF 1, Ary ‘ts An. A 85 #2, A
EY) 4> A FS RE MB G.
\end{Verbatim}

\end{note}

\begin{note}[来源路径]

近世代数/近世代数/无标题/置换与置换群 1bad2efa3f4280ed8022e05463b6b6f8.md

\end{note}

\section*{置换书写方式}

\begin{note}[来源上级主题]

置换群 (\%E7\%BD\%AE\%E6\%8D\%A2\%E7\%BE\%A4\%201bad2efa3f4280228c17c29195a2c525.md)

\end{note}

\begin{note}[置换书写方式]

例子

1分别对应1，2，3    2分别对应1，2，3

单位元就是123映射到123

\end{note}

\begin{note}[关联图片]
以下图片已纳入本条整理，当前版本优先依据 Markdown 文字整理，图片细节可在审核时对照：

1. image.png\\
2. image 1.png\\

\end{note}

\begin{note}[来源路径]

近世代数/近世代数/无标题/置换书写方式 1bcd2efa3f4280e08df6cc729ec19b71.md

\end{note}

\section*{置换群}

\begin{note}[来源上级主题]

群论 (\%E7\%BE\%A4\%E8\%AE\%BA\%201c3d2efa3f428064a77dd3ff465d3548.md)

\end{note}

\begin{note}[置换群]

子级 项目: 置换与置换群 (\%E7\%BD\%AE\%E6\%8D\%A2\%E4\%B8\%8E\%E7\%BD\%AE\%E6\%8D\%A2\%E7\%BE\%A4\%201bad2efa3f4280ed8022e05463b6b6f8.md), 对称群 (\%E5\%AF\%B9\%E7\%A7\%B0\%E7\%BE\%A4\%201bad2efa3f42809caaf1ff2a6c16b526.md), n次对称群S\_n的阶是n！ (n\%E6\%AC\%A1\%E5\%AF\%B9\%E7\%A7\%B0\%E7\%BE\%A4S\_n\%E7\%9A\%84\%E9\%98\%B6\%E6\%98\%AFn\%EF\%BC\%81\%201bad2efa3f4280749db7c5ede391c5b3.md), 循环置换  (\%E5\%BE\%AA\%E7\%8E\%AF\%E7\%BD\%AE\%E6\%8D\%A2\%201bad2efa3f4280109703cec9c3fb3605.md), 只要不相连就可以快速求阶数而且有交换律 (\%E5\%8F\%AA\%E8\%A6\%81\%E4\%B8\%8D\%E7\%9B\%B8\%E8\%BF\%9E\%E5\%B0\%B1\%E5\%8F\%AF\%E4\%BB\%A5\%E5\%BF\%AB\%E9\%80\%9F\%E6\%B1\%82\%E9\%98\%B6\%E6\%95\%B0\%E8\%80\%8C\%E4\%B8\%94\%E6\%9C\%89\%E4\%BA\%A4\%E6\%8D\%A2\%E5\%BE\%8B\%201bcd2efa3f4280e9add2d922e72a6f5d.md), 置换书写方式 (\%E7\%BD\%AE\%E6\%8D\%A2\%E4\%B9\%A6\%E5\%86\%99\%E6\%96\%B9\%E5\%BC\%8F\%201bcd2efa3f4280e08df6cc729ec19b71.md), 证明 可以写成积 (\%E8\%AF\%81\%E6\%98\%8E\%20\%E5\%8F\%AF\%E4\%BB\%A5\%E5\%86\%99\%E6\%88\%90\%E7\%A7\%AF\%201bcd2efa3f428091b11de9806d6b6c01.md)

\end{note}

\begin{note}[来源路径]

近世代数/近世代数/无标题/置换群 1bad2efa3f4280228c17c29195a2c525.md

\end{note}

\section*{群}

\begin{note}[来源上级主题]

映射与集合基础知识 (\%E6\%98\%A0\%E5\%B0\%84\%E4\%B8\%8E\%E9\%9B\%86\%E5\%90\%88\%E5\%9F\%BA\%E7\%A1\%80\%E7\%9F\%A5\%E8\%AF\%86\%201aed2efa3f42800c8261da73ac826269.md)

\end{note}

\begin{note}[群]

无标题

\end{note}

\begin{note}[来源路径]

近世代数/近世代数/无标题/群 1a0d2efa3f4280e5ab23eef494b59a4c.md

\end{note}

\section*{群同构第一定理1.2}

\begin{theorem}[群同构第一定理1.2]

子级 项目: 证明 (\%E8\%AF\%81\%E6\%98\%8E\%201ead2efa3f4280df93bace7a315e8dc8.md)

文本: 22

定义1.15群同态下定义 核是映射到e‘的原像集 imf是像集 如果G和G‘是群同态，ker（f）是G的一个子群

ker（f）是映射到e的原像集，G上由ker（f）的陪集组成群和im(f)也就是值域群上同构

如果f是满射，im（f）就是G‘，如果f是单同态ker（f）=\{e\}，

如果G有限则他的阶等于ker（f）的阶乘im（f）的乘积

命题1.37 商群

\end{theorem}

\begin{note}[关联图片]
以下图片已纳入本条整理，当前版本优先依据 Markdown 文字整理，图片细节可在审核时对照：

1. ef3c3956-eb85-4347-b28d-6e6a16bee2e0.png\\
2. image.png\\

\end{note}

\begin{note}[来源路径]

近世代数/近世代数/无标题/群同构第一定理1 2 1d8d2efa3f4280faaa25d06a12ca3ddc.md

\end{note}

\section*{群同构第三定理}

\begin{theorem}[群同构第三定理]

子级 项目: 证明 (\%E8\%AF\%81\%E6\%98\%8E\%201ead2efa3f42803d9277fa296434fc29.md), 无标题 (\%E6\%97\%A0\%E6\%A0\%87\%E9\%A2\%98\%201ead2efa3f4280d8976efa8f012a8275.md)

\end{theorem}

\begin{note}[关联图片]
以下图片已纳入本条整理，当前版本优先依据 Markdown 文字整理，图片细节可在审核时对照：

1. image.png\\

\end{note}

\begin{note}[来源路径]

近世代数/近世代数/无标题/群同构第三定理 1ead2efa3f428024abb0fc36b0b0317a.md

\end{note}

\section*{群同构第二定理}

\begin{theorem}[群同构第二定理]

子级 项目: 证明 (\%E8\%AF\%81\%E6\%98\%8E\%201ead2efa3f428093999af29d0c6c20e3.md)

N是正规子群，H是真子群 H和N的交集是H的正规子群，N是H中元素乘N构造陪集的正规子群

H 交N构成的子群的陪集组成的商群和H中元素与N组合构成的陪集的商群是同构的

\end{theorem}

\begin{note}[关联图片]
以下图片已纳入本条整理，当前版本优先依据 Markdown 文字整理，图片细节可在审核时对照：

1. image.png\\

\end{note}

\begin{note}[来源路径]

近世代数/近世代数/无标题/群同构第二定理 1ead2efa3f42806da76fe0cd08cb72be.md

\end{note}

\section*{群的同态同构}

\begin{note}[来源上级主题]

群的同态同构 (\%E7\%BE\%A4\%E7\%9A\%84\%E5\%90\%8C\%E6\%80\%81\%E5\%90\%8C\%E6\%9E\%84\%201c3d2efa3f428076a462f0397fbfd469.md)

\end{note}

\begin{note}[群的同态同构]

子级 项目: 命题1.13群同态单位元对单位元逆元对逆元 (\%E5\%91\%BD\%E9\%A2\%981\%2013\%E7\%BE\%A4\%E5\%90\%8C\%E6\%80\%81\%E5\%8D\%95\%E4\%BD\%8D\%E5\%85\%83\%E5\%AF\%B9\%E5\%8D\%95\%E4\%BD\%8D\%E5\%85\%83\%E9\%80\%86\%E5\%85\%83\%E5\%AF\%B9\%E9\%80\%86\%E5\%85\%83\%201cbd2efa3f4280048341d86cd285853a.md)

同态的定义

同构

满射可以多对一 同构必须一一对应

\end{note}

\begin{note}[关联图片]
以下图片已纳入本条整理，当前版本优先依据 Markdown 文字整理，图片细节可在审核时对照：

1. image.png\\

\end{note}

\begin{note}[来源路径]

近世代数/近世代数/无标题/群的同态同构 1b5d2efa3f42803aaaa7f6c2fdcfd0b4.md

\end{note}

\section*{群的同态同构}

\begin{note}[来源上级主题]

群论 (\%E7\%BE\%A4\%E8\%AE\%BA\%201c3d2efa3f428064a77dd3ff465d3548.md)

\end{note}

\begin{note}[群的同态同构]

子级 项目: 群的同态同构 (\%E7\%BE\%A4\%E7\%9A\%84\%E5\%90\%8C\%E6\%80\%81\%E5\%90\%8C\%E6\%9E\%84\%201b5d2efa3f42803aaaa7f6c2fdcfd0b4.md), 同态的意义 (\%E5\%90\%8C\%E6\%80\%81\%E7\%9A\%84\%E6\%84\%8F\%E4\%B9\%89\%201b5d2efa3f42801ba4aed9823eb1a0ae.md), 同态同构的意义 (\%E5\%90\%8C\%E6\%80\%81\%E5\%90\%8C\%E6\%9E\%84\%E7\%9A\%84\%E6\%84\%8F\%E4\%B9\%89\%201cad2efa3f42801b93c4f525f7a94957.md), 命题1.14 (\%E5\%91\%BD\%E9\%A2\%981\%2014\%201cbd2efa3f4280569ab4d1fcc291e1bc.md), 映射的核与像 (\%E6\%98\%A0\%E5\%B0\%84\%E7\%9A\%84\%E6\%A0\%B8\%E4\%B8\%8E\%E5\%83\%8F\%201cbd2efa3f428060a500d68691e7bbdf.md)

\end{note}

\begin{note}[来源路径]

近世代数/近世代数/无标题/群的同态同构 1c3d2efa3f428076a462f0397fbfd469.md

\end{note}

\section*{群的基础知识}

\begin{note}[来源上级主题]

群论 (\%E7\%BE\%A4\%E8\%AE\%BA\%201c3d2efa3f428064a77dd3ff465d3548.md)

\end{note}

\begin{note}[群的基础知识]

子级 项目: 群的定义 (\%E7\%BE\%A4\%E7\%9A\%84\%E5\%AE\%9A\%E4\%B9\%89\%201a8d2efa3f428012bcbcf151531edd91.md), 逆映射 (\%E9\%80\%86\%E6\%98\%A0\%E5\%B0\%84\%201a9d2efa3f42802b8a64f4efb72349db.md), 单位元 (\%E5\%8D\%95\%E4\%BD\%8D\%E5\%85\%83\%201aad2efa3f42804aa677fdb25912917d.md), 逆元 (\%E9\%80\%86\%E5\%85\%83\%201aad2efa3f4280ffa71ef933ebbf0d7f.md), 交换群 阿贝尔群 (\%E4\%BA\%A4\%E6\%8D\%A2\%E7\%BE\%A4\%20\%E9\%98\%BF\%E8\%B4\%9D\%E5\%B0\%94\%E7\%BE\%A4\%201aed2efa3f4280cdbde1e3cd83d58bc8.md), 群的等价判别 (\%E7\%BE\%A4\%E7\%9A\%84\%E7\%AD\%89\%E4\%BB\%B7\%E5\%88\%A4\%E5\%88\%AB\%201aed2efa3f4280fa9785ff77d205d695.md), 剩余类加群 (\%E5\%89\%A9\%E4\%BD\%99\%E7\%B1\%BB\%E5\%8A\%A0\%E7\%BE\%A4\%201aed2efa3f4280428135f1eab5e513d0.md), 群的意义 (\%E7\%BE\%A4\%E7\%9A\%84\%E6\%84\%8F\%E4\%B9\%89\%201aed2efa3f4280cf8e95f1b5d51c950a.md), 消去律成立与证明G是一个群 重点 (\%E6\%B6\%88\%E5\%8E\%BB\%E5\%BE\%8B\%E6\%88\%90\%E7\%AB\%8B\%E4\%B8\%8E\%E8\%AF\%81\%E6\%98\%8EG\%E6\%98\%AF\%E4\%B8\%80\%E4\%B8\%AA\%E7\%BE\%A4\%20\%E9\%87\%8D\%E7\%82\%B9\%201aed2efa3f42802bbde2ff09c5d7753f.md), 群的性质 (\%E7\%BE\%A4\%E7\%9A\%84\%E6\%80\%A7\%E8\%B4\%A8\%201aed2efa3f4280818433c4856f272a94.md), 定义1.20元素的阶 (\%E5\%AE\%9A\%E4\%B9\%891\%2020\%E5\%85\%83\%E7\%B4\%A0\%E7\%9A\%84\%E9\%98\%B6\%201afd2efa3f42804abe31e7c2129a941d.md), 群的阶与元素的阶 (\%E7\%BE\%A4\%E7\%9A\%84\%E9\%98\%B6\%E4\%B8\%8E\%E5\%85\%83\%E7\%B4\%A0\%E7\%9A\%84\%E9\%98\%B6\%201d3d2efa3f4280c89aacce78ed1a9c9f.md)

\end{note}

\begin{note}[来源路径]

近世代数/近世代数/无标题/群的基础知识 1aed2efa3f428009ac4eca74062880ff.md

\end{note}

\section*{群的定义}

\begin{note}[来源上级主题]

群的基础知识 (\%E7\%BE\%A4\%E7\%9A\%84\%E5\%9F\%BA\%E7\%A1\%80\%E7\%9F\%A5\%E8\%AF\%86\%201aed2efa3f428009ac4eca74062880ff.md)

\end{note}

\begin{definition}[群的定义]

乘法封闭 满足结合律  叫做乘法不一定是乘法是可以自定义的

定义Agroup⟨G,∗⟩是一个集合G，在二元运算∗下封闭，满足以下公理：

比如aob=a+b

有左单位元有左逆元

先证明是代数运算

代数运算定义

封闭 二元运算

线性空间定义 比线性空间少了数乘的约束更加一般

\end{definition}

\begin{note}[关联图片]
以下图片已纳入本条整理，当前版本优先依据 Markdown 文字整理，图片细节可在审核时对照：

1. image.png\\

\end{note}

\begin{note}[来源路径]

近世代数/近世代数/无标题/群的定义 1a8d2efa3f428012bcbcf151531edd91.md

\end{note}

\section*{群的性质}

\begin{note}[来源上级主题]

群的基础知识 (\%E7\%BE\%A4\%E7\%9A\%84\%E5\%9F\%BA\%E7\%A1\%80\%E7\%9F\%A5\%E8\%AF\%86\%201aed2efa3f428009ac4eca74062880ff.md)

\end{note}

\begin{note}[群的性质]

定理1左逆元也是右逆元

a逆是a的左逆元 a‘是a逆的左逆元再证明a’和a一样

群中元素满足结合律

左单位元也是右单位元

群的乘法适合消去律

\end{note}

\begin{note}[关联图片]
以下图片已纳入本条整理，当前版本优先依据 Markdown 文字整理，图片细节可在审核时对照：

1. image.png\\

\end{note}

\begin{note}[来源路径]

近世代数/近世代数/无标题/群的性质 1aed2efa3f4280818433c4856f272a94.md

\end{note}

\section*{群的意义}

\begin{note}[来源上级主题]

群的基础知识 (\%E7\%BE\%A4\%E7\%9A\%84\%E5\%9F\%BA\%E7\%A1\%80\%E7\%9F\%A5\%E8\%AF\%86\%201aed2efa3f428009ac4eca74062880ff.md)

\end{note}

\begin{note}[群的意义]

群论的诞生不仅源于对具体数学问题的解决需求，也体现了数学家对整体结构和对称性认识的不断深化。下面是更详细的背景和动因：

1. 多项式方程与伽罗瓦理论

- 多项式求解问题：

早期数学家在研究三次、四次乃至更高次多项式的求解方法时，逐渐发现这些方程的根之间存在着某种排列和对称关系。特别是伽罗瓦在19世纪初提出的伽罗瓦理论，通过研究根的置换群，判断一个多项式是否能用有限次根式求解。这一理论促使数学家将这些置换抽象成群，进而奠定了群论的基础。

2. 几何中的对称性

- 几何变换：

从古希腊时期起，人们就注意到各种几何图形（如正多边形、正多面体等）的对称性。旋转、反射、平移等几何变换具有很强的规律性，这种规律性可以用群论来描述。通过抽象化和统一讨论各种对称操作，群论为几何学提供了一个强大的工具，也促进了对称性在数学中的普遍应用。

3. 数学结构的统一与抽象化

- 抽象代数学的发展：

数学家逐步认识到，不同领域中许多问题都可以归纳为研究某种代数运算的性质。群只是众多代数结构中的一种，但它的公理体系（封闭性、结合律、单位元、逆元）简单而又富有力量，能够揭示对象间的内在联系。群论的抽象化推动了环、域、向量空间等更复杂结构的建立，使数学理论趋向统一和抽象。

4. 物理学与其他自然科学的应用

- 对称性与守恒定律：

在物理学中，对称性与守恒定律密切相关。比如，物理定律在某些变换下保持不变，这正是群论描述对称性的核心思想。现代物理学（如量子力学、粒子物理）中，群论被用来构建理论模型，解释粒子之间的相互作用，甚至预测新的物理现象。这一跨学科的应用进一步彰显了群论的普适性和重要性。

5. 数学美感与逻辑严谨

- 简洁而深远的公理体系：

群论仅用少量公理便能推导出丰富的结论，这种逻辑上的精炼和美感吸引了许多数学家深入研究和推广这一理论。正是这种抽象与统一的精神，让群论不仅成为解决实际问题的工具，也成为研究数学本质的重要领域。

总结

群论的发明可以看作数学发展过程中对“对称性”与“结构”认识不断深入的结果。它起初是为了解决具体的数学问题（如多项式方程求解和几何对称性）而提出，但其后迅速扩展成为一种统一描述和研究各种抽象结构的强大工具。无论是在数学内部构建严密理论，还是在物理和其他科学领域中解释自然现象，群论都发挥了不可替代的重要作用。

\end{note}

\begin{note}[来源路径]

近世代数/近世代数/无标题/群的意义 1aed2efa3f4280cf8e95f1b5d51c950a.md

\end{note}

\section*{群的等价判别}

\begin{note}[来源上级主题]

群的基础知识 (\%E7\%BE\%A4\%E7\%9A\%84\%E5\%9F\%BA\%E7\%A1\%80\%E7\%9F\%A5\%E8\%AF\%86\%201aed2efa3f428009ac4eca74062880ff.md)

\end{note}

\begin{note}[群的等价判别]

ER KGET AA MAGE AEE ©

INSEa RA, PATE:

ll. ABE:     Va,b,ceG, \#8

(ach)oc=aoc(boc)

Ill. G HARM AT a.b

Ti\#e ax =b, ya =b HECH AR

WK CG RAFRBGEE oo TERME.

\end{note}

\begin{note}[图片 OCR 摘录，待人工复核]
以下内容来自源图片识别，便于审核时对照：

\textbf{图片 1}（image.png）
\begin{Verbatim}[fontsize=\small]
ER KGET AA MAGE AEE ©
INSEa RA, PATE:
ll. ABE:     Va,b,ceG, #8
(ach)oc=aoc(boc)
Ill. G HARM AT a.b
Ti#e ax =b, ya =b HECH AR
WK CG RAFRBGEE oo TERME.
\end{Verbatim}

\end{note}

\begin{note}[来源路径]

近世代数/近世代数/无标题/群的等价判别 1aed2efa3f4280fa9785ff77d205d695.md

\end{note}

\section*{群的阶与元素的阶}

\begin{note}[来源上级主题]

群的基础知识 (\%E7\%BE\%A4\%E7\%9A\%84\%E5\%9F\%BA\%E7\%A1\%80\%E7\%9F\%A5\%E8\%AF\%86\%201aed2efa3f428009ac4eca74062880ff.md)

\end{note}

\begin{note}[群的阶与元素的阶]

群的阶是群中元素的个数

定义1.20元素的阶

群中每一个元素的阶可以是不一样的，

非循环群Klein 比如一个生于类加群\{0,1,2,3,4,5\} 2的阶就是3因为加法2+2+2=6也就是0类也就是这里面的e，但是2不是生成元，它没法生成所有的元素。但是由2生成的子群一定是循环群，当然子群还能由别打选取方式产生。就像不是所有群都是循环群但是所有子群的阶数都是父群阶的因子

拉格朗日定理

命题1.30循环群的生成元的幂指数和循环群阶数互素

\end{note}

\begin{note}[来源路径]

近世代数/近世代数/无标题/群的阶与元素的阶 1d3d2efa3f4280c89aacce78ed1a9c9f.md

\end{note}

\section*{群论}

\begin{note}[群论]

子级 项目: 群的基础知识 (\%E7\%BE\%A4\%E7\%9A\%84\%E5\%9F\%BA\%E7\%A1\%80\%E7\%9F\%A5\%E8\%AF\%86\%201aed2efa3f428009ac4eca74062880ff.md), 群的同态同构 (\%E7\%BE\%A4\%E7\%9A\%84\%E5\%90\%8C\%E6\%80\%81\%E5\%90\%8C\%E6\%9E\%84\%201c3d2efa3f428076a462f0397fbfd469.md), 变换群 (\%E5\%8F\%98\%E6\%8D\%A2\%E7\%BE\%A4\%201c3d2efa3f4280a3a416f8493ad4dae2.md), 置换群 (\%E7\%BD\%AE\%E6\%8D\%A2\%E7\%BE\%A4\%201bad2efa3f4280228c17c29195a2c525.md), 循环群 (\%E5\%BE\%AA\%E7\%8E\%AF\%E7\%BE\%A4\%201c3d2efa3f42802a8213f9e0014506cb.md)

\end{note}

\begin{note}[来源路径]

近世代数/近世代数/无标题/群论 1c3d2efa3f428064a77dd3ff465d3548.md

\end{note}

\section*{证明}

\begin{note}[来源上级主题]

群同构第三定理 (\%E7\%BE\%A4\%E5\%90\%8C\%E6\%9E\%84\%E7\%AC\%AC\%E4\%B8\%89\%E5\%AE\%9A\%E7\%90\%86\%201ead2efa3f428024abb0fc36b0b0317a.md)

\end{note}

\begin{proof}

UEY) BRE eH, N/M CG/M, RNEHERN/M ESB, ZEA. RN RAE MAN, KIL E

BAW. One€N, meM, Iilnmn'eM. Al N/M 24\#.

AARWADBZA, PUSRA N/M <G/M, EPFRANAUAIEA EME (RAR AWE

W, AEP OHA MEET RARER EMS RHR, RRA), RLS ERAN. SnMEN/M(nEN),

gM €G/M(geG), Il

(gM)(nM)(gM)\textasciitilde{}! = (gng7!)M € \{nM :n€N\}=N/M                            (1.95)

Auk N/M «G/M,

AA, RNBUX Sf: G/M G/N, XA

ff (gM) = gN                                                              (1.96)

BIA REX. RR eM=g'M, Wg 'e’ eM, Ke!’ EN, My gN=e'N.

(gMg’M) = f(gg’M) = gg'N = gNg’N = f(gM)g(g'M)                               (1.97)

WARILFHELAW. EM gNe€G/N(geG), Ml f(gM) =gN.

Ke, BRA?

ker(f) = \{gM : f(gM) = gN=eN\}=\{gM:geN\}=N/M                    (1.98)

ieee eH, ORS

(G/M)/(N/M) =\textasciitilde{} G/N                                (1.99)

\end{proof}

\begin{note}[图片 OCR 摘录，待人工复核]
以下内容来自源图片识别，便于审核时对照：

\textbf{图片 1}（image.png）
\begin{Verbatim}[fontsize=\small]
UEY) BRE eH, N/M CG/M, RNEHERN/M ESB, ZEA. RN RAE MAN, KIL E
BAW. One€N, meM, Iilnmn'eM. Al N/M 24#.
AARWADBZA, PUSRA N/M <G/M, EPFRANAUAIEA EME (RAR AWE
W, AEP OHA MEET RARER EMS RHR, RRA), RLS ERAN. SnMEN/M(nEN),
gM €G/M(geG), Il
(gM)(nM)(gM)~! = (gng7!)M € {nM :n€N}=N/M                            (1.95)
Auk N/M «G/M,
AA, RNBUX Sf: G/M G/N, XA
ff (gM) = gN                                                              (1.96)
BIA REX. RR eM=g'M, Wg 'e’ eM, Ke!’ EN, My gN=e'N.
(gMg’M) = f(gg’M) = gg'N = gNg’N = f(gM)g(g'M)                               (1.97)
WARILFHELAW. EM gNe€G/N(geG), Ml f(gM) =gN.
Ke, BRA?
ker(f) = {gM : f(gM) = gN=eN}={gM:geN}=N/M                    (1.98)
ieee eH, ORS
(G/M)/(N/M) =~ G/N                                (1.99)
\end{Verbatim}

\end{note}

\begin{note}[来源路径]

近世代数/近世代数/无标题/证明 1ead2efa3f42803d9277fa296434fc29.md

\end{note}

\section*{证明}

\begin{note}[来源上级主题]

群同构第二定理 (\%E7\%BE\%A4\%E5\%90\%8C\%E6\%9E\%84\%E7\%AC\%AC\%E4\%BA\%8C\%E5\%AE\%9A\%E7\%90\%86\%201ead2efa3f42806da76fe0cd08cb72be.md)

\end{note}

\begin{proof}

正规子群一共干的就是一件事从父群里面任意挑出a和对应的a的逆，再证明正规子群里面的元素可以达成aNa的逆=N，而证明的时候就要用增加a和a的逆把想要的形式凑出来在用拆括号把他们分开把需要的保留下来其他的去除掉，先把父群构成对应的映射，然后证明核就是H交N因为需要hN=eN所以h属于N

群同构第一定理1.2

\end{proof}

\begin{note}[关联图片]
以下图片已纳入本条整理，当前版本优先依据 Markdown 文字整理，图片细节可在审核时对照：

1. image.png\\

\end{note}

\begin{note}[来源路径]

近世代数/近世代数/无标题/证明 1ead2efa3f428093999af29d0c6c20e3.md

\end{note}

\section*{证明}

\begin{note}[来源上级主题]

群同构第一定理1.2 (\%E7\%BE\%A4\%E5\%90\%8C\%E6\%9E\%84\%E7\%AC\%AC\%E4\%B8\%80\%E5\%AE\%9A\%E7\%90\%861\%202\%201d8d2efa3f4280faaa25d06a12ca3ddc.md)

\end{note}

\begin{proof}

ker（f）被拿出来就是因为它能映射到e‘只要映射到G’这一项就会变成e’然后消失掉

要证明aN=Na就是证明aNa−1*a*−1=N然后因为同态就能把括号拆开，然后因为N有特殊的性质两边都能变成e‘然后就证明完毕了

证明同构映射，值域有多少个值ker（f）的陪集就有几个，所以一一对应就能构造同构

aN=a‘N代表他们只是不同的代表元，所以N=\$a\textasciicircum{}\{-1\}\$a‘N这也是为什么直接变成了e’

定义1.11n阶一般线性群

\end{proof}

\begin{note}[关联图片]
以下图片已纳入本条整理，当前版本优先依据 Markdown 文字整理，图片细节可在审核时对照：

1. image.png\\
2. image 1.png\\
3. image 2.png\\

\end{note}

\begin{note}[来源路径]

近世代数/近世代数/无标题/证明 1ead2efa3f4280df93bace7a315e8dc8.md

\end{note}

\section*{证明 可以写成积}

\begin{note}[来源上级主题]

置换群 (\%E7\%BD\%AE\%E6\%8D\%A2\%E7\%BE\%A4\%201bad2efa3f4280228c17c29195a2c525.md)

\end{note}

\begin{proof}

先证明集合中有重复的因为不是无限的；再让第一次出现重复的时候截断，证明k是不空的

\end{proof}

\begin{note}[关联图片]
以下图片已纳入本条整理，当前版本优先依据 Markdown 文字整理，图片细节可在审核时对照：

1. image.png\\
2. image 1.png\\

\end{note}

\begin{note}[来源路径]

近世代数/近世代数/无标题/证明 可以写成积 1bcd2efa3f428091b11de9806d6b6c01.md

\end{note}

\section*{运算表}

\begin{note}[来源上级主题]

代数运算 (\%E4\%BB\%A3\%E6\%95\%B0\%E8\%BF\%90\%E7\%AE\%97\%201aed2efa3f4280d597defe353a625299.md)

\end{note}

\begin{note}[运算表]

有限集的时候用运算表

https://www.notion.so

\end{note}

\begin{note}[关联图片]
以下图片已纳入本条整理，当前版本优先依据 Markdown 文字整理，图片细节可在审核时对照：

1. image.png\\

\end{note}

\begin{note}[来源路径]

近世代数/近世代数/无标题/运算表 1a1d2efa3f4280498713e856fc016324.md

\end{note}

\section*{逆元}

\begin{note}[来源上级主题]

群的基础知识 (\%E7\%BE\%A4\%E7\%9A\%84\%E5\%9F\%BA\%E7\%A1\%80\%E7\%9F\%A5\%E8\%AF\%86\%201aed2efa3f428009ac4eca74062880ff.md)

\end{note}

\begin{note}[逆元]

定义如果∗是S上的一个运算，且有恒等元e，那么对任意x∈S，x的逆元是x′，使得x∗x′=x′∗x=e。

\end{note}

\begin{note}[来源路径]

近世代数/近世代数/无标题/逆元 1aad2efa3f4280ffa71ef933ebbf0d7f.md

\end{note}

\section*{逆映射}

\begin{note}[来源上级主题]

群的基础知识 (\%E7\%BE\%A4\%E7\%9A\%84\%E5\%9F\%BA\%E7\%A1\%80\%E7\%9F\%A5\%E8\%AF\%86\%201aed2efa3f428009ac4eca74062880ff.md)

\end{note}

\begin{note}[逆映射]

如果X×Y的一个子集是X到Y的一一映射φ，那么φ中的每个x∈X都作为第一个成员出现在一个有序对中，每个y∈Y都作为第二个成员出现在一个有序对中。因此，如果我们交换φ中所有有序对（x，y）的第一个成员和第二个成员，得到一个有序对（y，x）的集合，我们得到一个Y×X的子集，它给出一个Y到X的一一映射。这个函数被称为φ的逆函数，用φ-1表示。总结一下，如果φ将X一对一地映射到Y，且φ（x）=y，那么φ-1将Y一对一地映射到X，且φ-1（y）=x。

\end{note}

\begin{note}[来源路径]

近世代数/近世代数/无标题/逆映射 1a9d2efa3f42802b8a64f4efb72349db.md

\end{note}

\section*{阶数与奇偶}

\begin{note}[来源上级主题]

定义1.20元素的阶 (\%E5\%AE\%9A\%E4\%B9\%891\%2020\%E5\%85\%83\%E7\%B4\%A0\%E7\%9A\%84\%E9\%98\%B6\%201afd2efa3f42804abe31e7c2129a941d.md)

\end{note}

\begin{note}[阶数与奇偶]

由a和a的逆出现偶数，有a的逆是群的定义规定的

\end{note}

\begin{note}[图片 OCR 摘录，待人工复核]
以下内容来自源图片识别，便于审核时对照：

\textbf{图片 1}（image.png）
\begin{Verbatim}[fontsize=\small]
2. fE—S ARE GT AF 2 och SRE EL.
RGAG IM CMB nihBka he, Bat 4a.
Wy
a® =aa"'=e

5n>2hRkFE. 2K, RNRA-MTANHAF 2
Hy 7c ¢ faq’,

HP GEAR, MGH BAF 2 Mace moth Hw,
PUG BAH t+ hee eR.
\end{Verbatim}

\end{note}

\begin{note}[来源路径]

近世代数/近世代数/无标题/阶数与奇偶 1b1d2efa3f4280fe9903f391009c0bee.md

\end{note}

\section*{阶的例题}

\begin{note}[来源上级主题]

定义1.20元素的阶 (\%E5\%AE\%9A\%E4\%B9\%891\%2020\%E5\%85\%83\%E7\%B4\%A0\%E7\%9A\%84\%E9\%98\%B6\%201afd2efa3f42804abe31e7c2129a941d.md)

\end{note}

\begin{example}[阶的例题]

1. BGHGHR—TPrcmB AWE x’ =e, BRAG ft

weit.

B@ ScoMbEGHERATtI. HB

(ab) (ab) = (ab)? =e

Fa Fi Tl

-              (ab) (ba) = ab? a=aea=a? =e

THA (ab) (ab) = (ab) (6a), FIBRE, B

ab= ba

PTELG ERR.                         -

\end{example}

\begin{note}[图片 OCR 摘录，待人工复核]
以下内容来自源图片识别，便于审核时对照：

\textbf{图片 1}（image.png）
\begin{Verbatim}[fontsize=\small]
1. BGHGHR—TPrcmB AWE x’ =e, BRAG ft
weit.
B@ ScoMbEGHERATtI. HB
(ab) (ab) = (ab)? =e
Fa Fi Tl
-              (ab) (ba) = ab? a=aea=a? =e
THA (ab) (ab) = (ab) (6a), FIBRE, B
ab= ba
PTELG ERR.                         -
\end{Verbatim}

\textbf{图片 2}（image 1.png）
\begin{Verbatim}[fontsize=\small]
filal=m, MSAa" =e, 9b(at)” =(a")y* =e* =e, Alja|<m,
tt la =n=>(a")" =e>(a")' =e>a" =e" =e.. | <n, |xEtin<m,
Hm<n=>me=nn,  MH jal =|a ‘|.
\end{Verbatim}

\textbf{图片 3}（image 2.png）
\begin{Verbatim}[fontsize=\small]
4. —A RY Poche PR.
fe SGHA-T#ERRMGEGHWE—TR, BA
ARAB ARARAIE. FETE IEE iy fii, (iai=al,
oi eB,         ”
(1)         ai-i=e
DOM, APTEIERORI-7, fi (1) Re. Ah e eR
ERA, fa"=e, XP, MB E™.
\end{Verbatim}

\textbf{图片 4}（image 3.png）
\begin{Verbatim}[fontsize=\small]
2. fE—-PARBBOAKF 2 Hoch SMe eB.
EGG OM OMB a) niko HB, Bia '#a,
FN
a? =aa"'=e

5n>2HRaFe. DH, RNRA-MTAMBKF 2
Hy 7c @ #ila~*,

HPG EAR, MGH BAT 2 Mace RH,
PLU G Hic Hoch + ke — EER.
\end{Verbatim}

\textbf{图片 5}（image 4.png）
\begin{Verbatim}[fontsize=\small]
3. BECR-TRERRN ARH. EG BGS F 2
HG BA See — fe a

2A, GRAAF 2 WEAR. (8
GRATE LNT BE RMGe. THAT ABE
a, ACERT 2 Nac PEE AH,
\end{Verbatim}

\end{note}

\begin{note}[来源路径]

近世代数/近世代数/无标题/阶的例题 1b1d2efa3f42801ca927ea10853179f8.md

\end{note}

\section*{陪集的定义}

\begin{note}[来源上级主题]

子群的陪集 (\%E5\%AD\%90\%E7\%BE\%A4\%E7\%9A\%84\%E9\%99\%AA\%E9\%9B\%86\%201c3d2efa3f42806a85b0f0113028b207.md)

\end{note}

\begin{definition}[陪集的定义]

@) Cit, Ae v

PUTA LURT BERR “OBR” (ER RR”), BBRILMBADEe.

BARS

Rik —TH G MEH FH H. PIB “RES”, HERG HHNRTBERE 9 KF

B BITE HY, 3l— MRA.

+ AR: AIA HW PASH, Fl

gH = \{g-h: he H\}.

+ BR: AI AR H PNB, Fl

Hg =\{h-g:he€ H\}.

ANAiRAREN?

BEMA—kA, BLARS aA, REACRHGCPHTR. FH ARKAP—-T

“SUA” FR. ME, MRMA— he (A RRE HB), Amie HLA

BFS TRAE (HMB 9), MENHMNMRAME—T “BR”. AEB”

PRED LARA AIC H BK “FS0)” BPR.

— Ma PHF

IF: BZ FH 2Z

+ HZ MBAR: (4, —3,—2,-1,0,1,2,3,4, \},

+ FH 27 BAA BM: (4, —2,0.2.4.\},

WEBER «, EMAMRA

a+2Z=\{a+2k:keZ\}.

RD:

“4a=0mt, pee O+ 2Z = \{,—4,—2,0,2,4,\} (RB 27 AE).

-4a=1h, be 1+ 2Z = \{,—3,-1,1,3,5,\},

dE: MABRRRSERH—MERH, HETSNIRETAR. iL, eS

SIF, PABMEA SEM (BF 0+2Z), BALE (RF 1+ 2Z),

RHE

1. DEH

BEAT LIL GC ORKAFTRESHDD. hiteit, STTRIMRTRIBEN

BR, MAREREZEAREARMTA.

2. RA—R

SMARNTAT (Bl “ANY ) SFR HAMA) (QR CG SABRE). RAH

AAA ENTERIC: FH HAG (Ady) BREE GC BND.

3. BMA

Base A BD FR A] 7 BE CG PF RH WA “(B” A “oH”, TEPIIERB (Quotient Group)

AY HEE SC HELE FA.

BS

+ CER MEMABRPN TTR I, OSH HW BURA BSAMNMES.

+ ARAN ABRNEMAW, RADREHAATE.

+ BSE ME GC DERAFMERWRD, SPRONA)MBSFTTEH HAWK.

\#5 2X MAGI F BER BT BaF IAS “PSR” XS. MRMATA RIA,

TRIAS HE [A] !

\end{definition}

\begin{note}[图片 OCR 摘录，待人工复核]
以下内容来自源图片识别，便于审核时对照：

\textbf{图片 1}（IMG\_2818.jpeg）
\begin{Verbatim}[fontsize=\small]
@) Cit, Ae v
PUTA LURT BERR “OBR” (ER RR”), BBRILMBADEe.
BARS
Rik —TH G MEH FH H. PIB “RES”, HERG HHNRTBERE 9 KF
B BITE HY, 3l— MRA.
+ AR: AIA HW PASH, Fl
gH = {g-h: he H}.
+ BR: AI AR H PNB, Fl
Hg ={h-g:he€ H}.
ANAiRAREN?
BEMA—kA, BLARS aA, REACRHGCPHTR. FH ARKAP—-T
“SUA” FR. ME, MRMA— he (A RRE HB), Amie HLA
BFS TRAE (HMB 9), MENHMNMRAME—T “BR”. AEB”
PRED LARA AIC H BK “FS0)” BPR.
— Ma PHF
IF: BZ FH 2Z
+ HZ MBAR: (4, —3,—2,-1,0,1,2,3,4, },
+ FH 27 BAA BM: (4, —2,0.2.4.},
WEBER «, EMAMRA
a+2Z={a+2k:keZ}.
RD:
“4a=0mt, pee O+ 2Z = {,—4,—2,0,2,4,} (RB 27 AE).
-4a=1h, be 1+ 2Z = {,—3,-1,1,3,5,},
dE: MABRRRSERH—MERH, HETSNIRETAR. iL, eS
SIF, PABMEA SEM (BF 0+2Z), BALE (RF 1+ 2Z),
RHE
1. DEH
BEAT LIL GC ORKAFTRESHDD. hiteit, STTRIMRTRIBEN
BR, MAREREZEAREARMTA.
2. RA—R
SMARNTAT (Bl “ANY ) SFR HAMA) (QR CG SABRE). RAH
AAA ENTERIC: FH HAG (Ady) BREE GC BND.
3. BMA
Base A BD FR A] 7 BE CG PF RH WA “(B” A “oH”, TEPIIERB (Quotient Group)
AY HEE SC HELE FA.
BS
+ CER MEMABRPN TTR I, OSH HW BURA BSAMNMES.
+ ARAN ABRNEMAW, RADREHAATE.
+ BSE ME GC DERAFMERWRD, SPRONA)MBSFTTEH HAWK.
#5 2X MAGI F BER BT BaF IAS “PSR” XS. MRMATA RIA,
TRIAS HE [A] !
\end{Verbatim}

\end{note}

\begin{note}[来源路径]

近世代数/近世代数/无标题/陪集的定义 1c4d2efa3f42804e9574d5a966680083.md

\end{note}

\section*{陪集的意义}

\begin{note}[来源上级主题]

子群的陪集 (\%E5\%AD\%90\%E7\%BE\%A4\%E7\%9A\%84\%E9\%99\%AA\%E9\%9B\%86\%201c3d2efa3f42806a85b0f0113028b207.md)

\end{note}

\begin{note}[陪集的意义]

1. SER (BERD):

Sa eE— Sa G URAF HH HI, G IRIS ABF TAMAS “BER”. PRD,

ABR Meu

gH=\{gh:he H\}, geG.

IEE BE RIE 3S BER GRC RIES FG. XAPRID AS Ee ERR A BAS t9, FS

eS MARIA) BST FH AA, MMA ERBHIR BRA Ete Btn.

2. WEAAHItRBA A ERE:

AUFABE RAID, BTLAUEAABE G AUB) (TTA) STH AWMS EG AHIR

SANFRAR. eh:

IG| =(G: H]- |HI,

Bh |G: A] 2 WARK. RATT RMOMNDEIER, MAWAAT FROIN

ERRETHHM.

3. PIERRE TEAL BF:

4—S 3 H REMFRN (ERE SARS), TER C/H, BBR

RARE, PRETRRHAWEBNORRAMA BRA ARIER.

4. FEAR EEF RA ARTE :

BRERA Phi HeSE AE. tow, BURARERTRALNIER, Bilal

BARES, MXZNVESMRVNABARUND SIR. HM, RPDHNARTD

ATT AREAL it Blo) eA,

5. RR SCETG] A:

E-KAR MH, URBRAR, DTRREHWAWS, RNR Shek TAM

WLR. BUR, we RESERN A, ECT ATLA ARO RRABF “SA” eS, M

Ti iai 10 (ol AY SR AIS FE

AZ, BRNAD PRIM ee (FE) HRIUIRE THA, SHICH—TMER

RBOMBANLA, CMREBICUEAP AA BER (MOAR BAA EEE. RMI), ih

BETS SMA ARG PR RES SSE.

\end{note}

\begin{note}[图片 OCR 摘录，待人工复核]
以下内容来自源图片识别，便于审核时对照：

\textbf{图片 1}（IMG\_2827.jpeg）
\begin{Verbatim}[fontsize=\small]
1. SER (BERD):

Sa eE— Sa G URAF HH HI, G IRIS ABF TAMAS “BER”. PRD,
ABR Meu

gH={gh:he H}, geG.

IEE BE RIE 3S BER GRC RIES FG. XAPRID AS Ee ERR A BAS t9, FS
eS MARIA) BST FH AA, MMA ERBHIR BRA Ete Btn.

2. WEAAHItRBA A ERE:

AUFABE RAID, BTLAUEAABE G AUB) (TTA) STH AWMS EG AHIR
SANFRAR. eh:

IG| =(G: H]- |HI,

Bh |G: A] 2 WARK. RATT RMOMNDEIER, MAWAAT FROIN
ERRETHHM.

3. PIERRE TEAL BF:

4—S 3 H REMFRN (ERE SARS), TER C/H, BBR
RARE, PRETRRHAWEBNORRAMA BRA ARIER.

4. FEAR EEF RA ARTE :

BRERA Phi HeSE AE. tow, BURARERTRALNIER, Bilal
BARES, MXZNVESMRVNABARUND SIR. HM, RPDHNARTD
ATT AREAL it Blo) eA,

5. RR SCETG] A:

E-KAR MH, URBRAR, DTRREHWAWS, RNR Shek TAM
WLR. BUR, we RESERN A, ECT ATLA ARO RRABF “SA” eS, M
Ti iai 10 (ol AY SR AIS FE

AZ, BRNAD PRIM ee (FE) HRIUIRE THA, SHICH—TMER
RBOMBANLA, CMREBICUEAP AA BER (MOAR BAA EEE. RMI), ih
BETS SMA ARG PR RES SSE.
\end{Verbatim}

\end{note}

\begin{note}[来源路径]

近世代数/近世代数/无标题/陪集的意义 1c6d2efa3f42801ba4c4d1f60871c4bc.md

\end{note}

\section*{集合A的任一个等价关系都可确定A的一个分类}

\begin{note}[来源上级主题]

分类与关系 (\%E5\%88\%86\%E7\%B1\%BB\%E4\%B8\%8E\%E5\%85\%B3\%E7\%B3\%BB\%201aed2efa3f42802491cdeb2b56b5f1f7.md)

\end{note}

\begin{note}[集合A的任一个等价关系都可确定A的一个分类]

等价关系

条件给的是一个等价关系 要证明的是确定了一个分类

分类

[a]是所有和a等价的元素组成的集合 等价有

2不同类的子集相交时 如果a集合和b集合不等价那么a集合和b集合的交集是空集，因为如果a集合和b集合不等价还有交集x a的集和中都是与x等价的元素  b的集合中也是都是和x等价的元素所以a集合和b集合等价   🟰的威力

\end{note}

\begin{note}[关联图片]
以下图片已纳入本条整理，当前版本优先依据 Markdown 文字整理，图片细节可在审核时对照：

1. image.png\\

\end{note}

\begin{note}[来源路径]

近世代数/近世代数/无标题/集合A的任一个等价关系都可确定A的一个分类 1a8d2efa3f4280d68789ddcfda4d73ac.md

\end{note}

\section*{集合的映射}

\begin{note}[来源上级主题]

映射与集合基础知识 (\%E6\%98\%A0\%E5\%B0\%84\%E4\%B8\%8E\%E9\%9B\%86\%E5\%90\%88\%E5\%9F\%BA\%E7\%A1\%80\%E7\%9F\%A5\%E8\%AF\%86\%201aed2efa3f42800c8261da73ac826269.md)

\end{note}

\begin{note}[集合的映射]

1.48 4,4,,°:,4,,D WOR;

2. 4,,4,,°°°, A, ARPA BERR

3. BRA — iE SE PIC E PR;

4. —*70)R eA ME WIR;

5. PAN RABY AED WY 7c.

\end{note}

\begin{note}[图片 OCR 摘录，待人工复核]
以下内容来自源图片识别，便于审核时对照：

\textbf{图片 1}（IMG\_2628.jpeg）
\begin{Verbatim}[fontsize=\small]
1.48 4,4,,°:,4,,D WOR;
2. 4,,4,,°°°, A, ARPA BERR

3. BRA — iE SE PIC E PR;
4. —*70)R eA ME WIR;

5. PAN RABY AED WY 7c.
\end{Verbatim}

\end{note}

\begin{note}[来源路径]

近世代数/近世代数/无标题/集合的映射 1a1d2efa3f428022b697d8ac2040bd12.md

\end{note}

\section*{集合的积}

\begin{note}[来源上级主题]

映射与集合基础知识 (\%E6\%98\%A0\%E5\%B0\%84\%E4\%B8\%8E\%E9\%9B\%86\%E5\%90\%88\%E5\%9F\%BA\%E7\%A1\%80\%E7\%9F\%A5\%E8\%AF\%86\%201aed2efa3f42800c8261da73ac826269.md)

\end{note}

\begin{note}[集合的积]

WA, 4,,-, 4, (n AEBS) BRERA. H-W

MAA, 4, BU RW wR A

(a,,4,,°°°,4,),a,€ A,, PNA

\{(@,,4,,°:°,4,)|a, €A,,i=1, 2,:--n\}

BRA A,, A,,--, 4, BR, WE

A x A, xx A.

\end{note}

\begin{note}[图片 OCR 摘录，待人工复核]
以下内容来自源图片识别，便于审核时对照：

\textbf{图片 1}（IMG\_2627.jpeg）
\begin{Verbatim}[fontsize=\small]
WA, 4,,-, 4, (n AEBS) BRERA. H-W
MAA, 4, BU RW wR A
(a,,4,,°°°,4,),a,€ A,, PNA

{(@,,4,,°:°,4,)|a, €A,,i=1, 2,:--n}
BRA A,, A,,--, 4, BR, WE
A x A, xx A.
\end{Verbatim}

\end{note}

\begin{note}[来源路径]

近世代数/近世代数/无标题/集合的积 1a1d2efa3f4280f19e43c59633af2d85.md

\end{note}
